% !TeX root = ../Book2.tex
\BookPart{变分法、弱解与椭圆方程}{b2:part:03}
% 本编导读（章级标题，不计入编号）
\BookPartGuide

上一编的显式核和函数展开在规则区域上十分有效；当区域不规则、系数随位置变化，或源项只具有有限可积性时，直接公式往往无法取得。本编改从函数空间中的能量和测试等式出发，依次回答三个问题：解是否存在，是否唯一并稳定依赖数据，存在的弱解又有多光滑。

\ProseParagraphBreak

\chapref{b2:ch:14}从允许扰动推导 Euler--Lagrange 方程，再用直接方法构造凸能量的极小元。\chapref{b2:ch:15}把对称能量推广为一般双线性型，由 Lax--Milgram 定理建立弱解，并处理有界可测系数。\chapref{b2:ch:16}以差商和截断估计研究正则性，分别说明系数、边界、维数和源项空间的作用。下图的箭头表示本编实际使用的前置关系。
\[
\begin{tikzcd}[column sep=2.7em,row sep=1.5em]
\text{第一卷弱收敛与 Sobolev 理论}\arrow[d]\arrow[dr]\arrow[drr]&&\\
\text{第\ref{b2:ch:14}章}\arrow[r]&\text{第\ref{b2:ch:15}章}\arrow[r]&\text{第\ref{b2:ch:16}章}
\end{tikzcd}
\]
本编反复使用第一卷第八至九章的 Hilbert 与弱拓扑工具、第二十二章的零边界 Sobolev 空间，以及第二十四至二十五章的嵌入与紧性。\chapref{b2:ch:15}的 Fredholm 选读小节还将结合第一卷附录 C，讨论强制性失效后可能出现的有限维核与相容条件。

\ProseParagraphBreak

阅读时宜始终把边界类和测试空间写在方程旁边。同一内部表达式配上零迹、零均值或自由边界测试，会成为不同问题。还应区分三个层次：第一变分为零只是驻点条件，强制型给出能量解，正则性估计才允许逐点读取微分方程。

\ProseParagraphBreak

\chapref{b2:ch:16}末从可测系数标量方程的局部有界性走向 Harnack 比较及 Hölder 连续性。普通阅读可以先接受\cref{b2:thm:16-measurable-harnack}，继续其后的 Hölder 推论；需要完整证明时，先补\secref{b2:sec:1701}中二进立方体与分化的工具，再读\cref{b2:thm:C-john-nirenberg}的指数尾，最后读\secref{b2:sec:H201}的幂迭代、对数测试及\cref{b2:thm:H2-linear-harnack}。掌握本编后，可以继续第五编的实变调和分析，学习极大函数、奇异积分和频率分解怎样为更精细的方程估计提供工具。

% !TeX root = ../../Book2.tex
\chapter{变分法}
\label{b2:ch:14}
\BookChapterNumber{b2:ch:14}
% 本章摘要（不计入编号）
\begin{chapterabstract}
本章用函数空间中的微分研究积分泛函，推导标量与向量值 Euler--Lagrange 方程，并从边界项辨认自然边界条件。Dirichlet 能量说明凸性怎样把驻点变成唯一极小点；波动作用量则说明驻点未必是极小点。对梯度幂能量，直接方法通过强制性、弱闭约束与弱下半连续性构造极小元。若能量无下界，寻找驻点就需要另一条道路：选读部分用 Palais--Smale 条件控制近临界序列，以变形引理和山路定理寻找非极小型临界点，并应用于次临界半线性椭圆方程。
\end{chapterabstract}

% 本章目标（不计入编号）
\begin{learninggoals}
\begin{itemize}
\item\label{b2:item:14-goal-derivative}区分方向变分、Gâteaux 导数与 Fréchet 导数，检查积分下求导所需的控制。
\item\label{b2:item:14-goal-boundary}从允许扰动出发推导内部方程及自然边界条件，核对外法向的符号。
\item\label{b2:item:14-goal-energy}用精确能量差或第二变分判断极小性，区别作用量与守恒能量。
\item\label{b2:item:14-goal-direct}在 Sobolev 边界类中实施直接方法，指出强制性、弱紧性和下半连续性各自的作用。
\item\label{b2:item:14-goal-mountain}理解 Palais--Smale 条件控制什么序列，用下降变形解释山路临界点，并验证次临界椭圆模型中的几何条件与强收敛。
\end{itemize}
\end{learninggoals}

一条连接固定端点的曲线怎样使长度最小？一张固定在边框上的薄膜怎样达到平衡？这类问题比较的是整个函数，而不只是有限个参数。微积分仍然有效，但导数、紧性和边界条件都必须放到合适的函数空间中。候选曲线满足什么方程，最小值又是否真的能够达到，是两个相连的问题：前者由变分导数回答，后者需要研究极小化序列的极限。

\ProseParagraphBreak

极小化序列可能没有逐点收敛的子列，却仍能在弱收敛下保留能量信息。第一卷第九章的凸泛函理论解释这种极限为何有用；第二十一至二十四章的 Sobolev 空间与 Poincaré 不等式，则帮助我们同时控制函数、梯度和边界值。关于变分法的系统叙述，可进一步阅读 Bernard Dacorogna 的《Direct Methods in the Calculus of Variations》\cite{Dacorogna2008Direct}及 Lawrence C. Evans 的《Partial Differential Equations》\cite{Evans2010Partial}。

\ProseParagraphBreak

极小化并不能找到所有平衡态。一个能量可能沿某些方向趋于负无穷，却在离开原点时先经过一道正的能量屏障。连接原点与低能量点的路径必须越过这道屏障；让所有路径上的最高能量尽可能低，便产生山路构造。\secref{b2:sec:1405}研究这个几何想法何时能产生真正的临界点。

% !TeX root = ../../Book2.tex
\section{泛函与变分}
\label{b2:sec:1401}

固定两端的曲线可以改变形状，而端点条件保持不变。若用函数 \(u\) 表示曲线，一次允许的微小改变便写成 \(u+th\)，其中 \(h\) 在端点为零。长度、面积和能量都是作用于整条曲线的函数；考察它们随 \(t\) 的变化，就把极值问题带回了一元微积分。不过，在函数空间中，沿每个方向分别求导与对所有小扰动作统一估计并不是同一件事。本节先辨明这一差别，再说明积分泛函的变分怎样计算。

\subsection{Gâteaux 导数}
\label{b2:subsec:140101}

设 \(X\) 是实 Banach 空间，\(U\subset X\) 是开集，\(J:U\to\mathbb R\)。给定 \(u\in U\) 和 \(h\in X\)，若极限存在，记
\[
 \delta J(u;h)=\lim_{t\to0}\frac{J(u+th)-J(u)}{t}.
\]
这里的参数 \(t\) 取实数；当未知函数取复值而泛函取实值时，也应把函数空间看作实线性空间。比如 \(J(u)=\lVert u\rVert^2\) 的变分含有内积的实部，并不是复线性泛函。

\Needspace{6\baselineskip}
\begin{definition}[Gâteaux 可微]
\label{b2:def:14-gateaux}
设 \(X\) 是实 Banach 空间，\(J:U\to\mathbb R\)，\(U\subset X\) 开、\(u\in U\)；\(\delta J(u;h)=\lim_{t\to0}(J(u+th)-J(u))/t\)。若各方向的变分都存在，且存在连续线性泛函 \(A\in X^*\)，使得 \(\delta J(u;h)=Ah\) 对所有 \(h\in X\) 成立，则称 \(J\) 在 \(u\) 处 \term[gateaux-kewei]{Gâteaux 可微}[Gâteaux differentiable]，并称 \(A\) 为其 \term[gateaux-daoshu]{Gâteaux 导数}[Gâteaux derivative]，记作 \(D_GJ(u)\)。
\end{definition}

定义中的线性性不是由方向极限自动保证的。在 \(\mathbb R^2\) 上令
\[
 F(x,y)=\begin{cases}x^3/(x^2+y^2),&(x,y)\ne(0,0),\\0,&(x,y)=(0,0).\end{cases}
\]
沿方向 \((a,b)\ne(0,0)\) 的导数是 \(a^3/(a^2+b^2)\)。这个表达式关于方向不是线性的，故 \(F\) 虽有所有方向导数，却不满足这里的 Gâteaux 可微定义。不同文献有时把“方向导数均存在”也称为 Gâteaux 可微；使用结论时，应核对是否同时要求导数为连续线性映射。

\Needspace{6\baselineskip}
\begin{proposition}[极小点的必要条件]
\label{b2:prop:14-gateaux-minimum}
设 \(X\) 是实 Banach 空间，\(V\subset X\) 是线性子空间，\(J\) 是定义在 \(u\in X\) 某开邻域内的实值泛函。若 \(u\) 是 \(J\) 在仿射空间 \(u+V\) 上的局部极小点，且 \(\delta J(u;h)\) 对每个 \(h\in V\) 都存在，则
\[\delta J(u;h)=0\qquad(h\in V).\]
\end{proposition}
\begin{proof}
固定 \(h\in V\)。当 \(t\) 的绝对值充分小时，\(J(u+th)-J(u)\ge0\)。除以正的 \(t\) 再令 \(t\downarrow0\)，得变分非负；除以负的 \(t\) 再令 \(t\uparrow0\)，得变分非正。两个单侧极限相同，故变分为零。
\end{proof}

此处真正起作用的是正负两个方向都可行。若要求 \(u\ge0\)，在接触零值的区域只能作某些单侧扰动，必要条件一般变成不等式。类似地，固定端点要求扰动在端点消失；固定积分则要求扰动的积分为零。求第一变分之前，先写清允许的扰动，比形式地求导更为重要。

\subsection{Fréchet 导数}
\label{b2:subsec:140102}

第一章在有限维空间中讨论过 Fréchet 导数；同一定义适用于 Banach 空间。设 \(X\) 为实 Banach 空间，\(U\subset X\) 开，\(J:U\to\mathbb R\)、\(u\in U\)。泛函 \(J\) 在 \(u\) 处 Fréchet 可微，意味着存在 \(A\in X^*\)，使
\[
 J(u+h)=J(u)+Ah+r(h),\qquad
 \frac{|r(h)|}{\lVert h\rVert_X}\longrightarrow0\quad(h\to0).
\]
这里的余项估计对所有小方向统一成立。因此 Fréchet 可微必然蕴含 Gâteaux 可微，且两种导数相同；反过来则需要额外条件。

\Needspace{6\baselineskip}
\begin{proposition}[导数连续时的统一逼近]
\label{b2:prop:14-continuous-gateaux}
设 \(X\) 是实 Banach 空间，实值泛函 \(J\) 在 \(u\) 的一个开邻域内 Gâteaux 可微。若 \(v\mapsto D_GJ(v)\) 在 \(u\) 处关于 \(X^*\) 的算子范数连续，则 \(J\) 在 \(u\) 处 Fréchet 可微，且 \(DJ(u)=D_GJ(u)\)。
\end{proposition}
\begin{proof}
给定 \(\varepsilon>0\)，选取 \(r>0\)，使 \(\lVert v-u\rVert_X<r\) 时有 \(\lVert D_GJ(v)-D_GJ(u)\rVert_{X^*}<\varepsilon\)。对 \(\lVert h\rVert_X<r\)，在 \([0,1]\) 上考察
\[g(t)=J(u+th)-J(u)-tD_GJ(u)h.\]
它在区间上连续、内部可微，且 \(|g'(t)|\le\varepsilon\lVert h\rVert_X\)。一元中值定理给出 \(|g(1)|\le\varepsilon\lVert h\rVert_X\)，正是所需的统一余项估计。
\end{proof}

即便 Gâteaux 导数在一点是零映射，原函数也未必在该点连续。\subsecref{b2:subsec:010101}已经计算过的反例是
\[
 J(x,y)=\begin{cases}\displaystyle\frac{x^3y}{x^6+y^2},&(x,y)\ne(0,0),\\0,&(x,y)=(0,0).\end{cases}
\]
它在原点的所有方向导数均为零，即 \(D_GJ(0)=0\)，但沿 \(y=x^3\)、\(x\ne0\) 有 \(J(x,x^3)=1/2\)，故原点处不连续，也不可能 Fréchet 可微。这里需要保留的新区别是：每个固定方向的极限存在，并不提供对小扰动统一的余项估计；前一命题增加的导数算子范数连续性，才足以把这些逐方向极限接成同一个局部线性逼近。

\subsection{第一变分}
\label{b2:subsec:140103}

对实 Banach 空间 \(X\) 的开集 \(U\) 上的 Fréchet 可微泛函 \(J:U\to\mathbb R\)，\(u\in U\)、\(h\in X\) 时的 \(DJ(u)h\) 即泛函对扰动的一阶线性响应，称为 \term[diyi-bianfen]{第一变分}[first variation]。对积分泛函，常见计算是“把导数移入积分号”。这一计算需要一致的支配或余项估计；下面先在 \(C^1\) 空间内证明一个便于直接使用的版本。

\Needspace{6\baselineskip}
\begin{proposition}[光滑积分泛函的第一变分]
\label{b2:prop:14-smooth-integral-variation}
设整数 \(d\geq1\)，\(\Omega\subset\mathbb R^d\) 是有界开集。实函数 \(F(x,z,p)\) 以及对 \((z,p)\in\mathbb R\times\mathbb R^d\) 的一阶偏导数在 \(\overline\Omega\times\mathbb R\times\mathbb R^d\) 上连续。赋予实函数空间 \(C^1(\overline\Omega;\mathbb R)\) 以函数及一阶导数的一致范数。则
\[J(u)=\int_\Omega F(x,u(x),\nabla u(x))\,\mathrm dx\]
为 Fréchet 可微泛函，且
\[
 DJ(u)h=\int_\Omega\bigl(F_z(x,u,\nabla u)h+
             F_p(x,u,\nabla u)\cdot\nabla h\bigr)\,\mathrm dx.
\]
其导数关于 \(u\) 在算子范数下连续。
\end{proposition}
\begin{proof}
固定 \(u\)。当 \(\lVert h\rVert_{C^1}\le1\) 且 \(0\le t\le1\) 时，\((x,u(x)+th(x),\nabla u(x)+t\nabla h(x))\) 落在一个固定紧集内。在这个紧集上，\((z,p)\) 的一阶偏导数一致连续；以 \(\omega(s)\) 表示相应连续模，其中 \(\omega(s)\to0\)。沿 \((h,\nabla h)\) 应用一元微积分基本定理，得
\[
 \begin{split}
 &F(x,u+h,\nabla u+\nabla h)-F(x,u,\nabla u)\\
 &\quad-F_z(x,u,\nabla u)h-F_p(x,u,\nabla u)\cdot\nabla h\\
 &=\int_0^1\bigl[D_{(z,p)}F(x,u+th,\nabla u+t\nabla h)
              -D_{(z,p)}F(x,u,\nabla u)\bigr]\cdot(h,\nabla h)\,\mathrm dt.
 \end{split}
\]
积分后，余项的绝对值不超过
\(C_d|\Omega|\cdot\omega(C_d\lVert h\rVert_{C^1})\cdot\lVert h\rVert_{C^1}\)。所写的一阶项是有界线性泛函，故它就是 Fréchet 导数。用同一个紧集上的连续模比较 \(DJ(u+k)\) 与 \(DJ(u)\)，再对 \(\lVert h\rVert_{C^1}\le1\) 取上确界，就得到导数的算子范数连续性。
\end{proof}

例如，在有界开集 \(\Omega\) 上取实值 \(f\in C(\overline\Omega)\)，则 \(F(x,z,p)=\tfrac12|p|^2-f(x)z\) 满足前一命题的连续性条件，在 \(C^1(\overline\Omega;\mathbb R)\) 上给出
\[DJ(u)h=\int_\Omega\nabla u\cdot\nabla h\,\mathrm dx-\int_\Omega fh\,\mathrm dx.\]
改在 \(H^1(\Omega)\) 上讨论时，上一命题不能原样使用：\(H^1\) 范数一般不控制函数的一致范数，非线性项的增长速度必须与 Sobolev 嵌入相容。特别是，写下 \(\int |u|^q\) 之前，须先确认每个允许的 \(u\) 确有这一可积性。

二次能量可以避开一致范数，直接在其自然空间中计算。设 \(\Omega\subset\mathbb R^d\) 为开集，只要求实值 \(f\in L^2(\Omega)\)，令
\[
 J_f(u)=\frac12\int_\Omega|\nabla u|^2\,\mathrm{d}x
          -\int_\Omega fu\,\mathrm{d}x,
 \qquad u\in H^1(\Omega;\mathbb R).
\]
对每个 \(u,h\in H^1(\Omega;\mathbb R)\)，候选线性响应
\[
 A_uh=\int_\Omega\nabla u\cdot\nabla h\,\mathrm{d}x
             -\int_\Omega fh\,\mathrm{d}x
\]
满足 \(|A_uh|\leq(\|\nabla u\|_2+\|f\|_2)\|h\|_{H^1}\)。直接展开平方还给出精确余项
\[
 J_f(u+h)-J_f(u)-A_uh=\frac12\|\nabla h\|_2^2
                       \leq\frac12\|h\|_{H^1}^2.
\]
除以 \(\|h\|_{H^1}\) 后趋零，故 \(DJ_f(u)=A_u\)。而 \(\|A_u-A_v\|_{(H^1)^*}\leq\|\nabla(u-v)\|_2\)，导数在算子范数下连续。这里的 \(f\) 不必连续；\(L^2\) 配对的有界性及二次余项已经给出所需的可微性。本章后面将以这个空间中的能量进行极小化。

梯度幂能量可以直接在自然的可积空间中微分。这里不能用逐点求导代替范数余项估计，因为扰动可能在很小的集合上很大。

\Needspace{6\baselineskip}
\begin{proposition}[梯度幂能量的连续微分]
\label{b2:prop:14-power-frechet}
设整数 \(d\geq1\)，\(\Omega\subset\mathbb R^d\) 为开集。若 \(1<p<\infty\)，则 \(P(u)=p^{-1}\int_\Omega|\nabla u|^p\) 在实空间 \(W^{1,p}(\Omega)\) 上连续 Fréchet 可微，且
\[DP(u)h=\int_\Omega |\nabla u|^{p-2}\nabla u\cdot\nabla h.\]
\end{proposition}
\begin{proof}
置 \(A_p(z)=|z|^{p-2}z\)，并规定 \(A_p(0)=0\)；\(p=2\) 时直接有 \(A_2(z)=z\)。先核验使积分余项一致变小的两个向量估计：
\[
 |A_p(a+b)-A_p(a)|\le
 \begin{cases}
 C_p(|a|+|b|)^{p-2}|b|,&p\ge2,\\
 C_p|b|^{p-1},&1<p<2.
 \end{cases}
\]
在非零点，\(DA_p(z)=|z|^{p-2}I+(p-2)|z|^{p-4}z\otimes z\)，其算子范数不超过 \(C_p|z|^{p-2}\)。当 \(p\ge2\) 时，沿 \(a+tb\) 积分即得第一式，经过零点也可分成两段取极限。当 \(p<2\)、\(|b|\ge |a|/2\) 时，分别估计两端的 \(|A_p|\)，得到第二式；当 \(|b|<|a|/2\) 时，整条线段离零点至少 \(|a|/2\)，导数估计给出 \(C_p|a|^{p-2}|b|\le C_p|b|^{p-1}\)。

对 \(v,w\in L^p(\Omega;\mathbb R^d)\)，Hölder 不等式于是给出
\[
 \|A_p(v+w)-A_p(v)\|_{p'}\le
 \begin{cases}
 C_p(\|v\|_p+\|w\|_p)^{p-2}\|w\|_p,&p\ge2,\\
 C_p\|w\|_p^{p-1},&1<p<2,
 \end{cases}
 \qquad p'=\frac p{p-1}.
\]
第一种情形用的指数关系是 \(1/p'=(p-2)/p+1/p\)；第二种情形直接把 \(|w|^{p-1}\) 的 \(p'\) 次方积分。另有 \(\|A_p(v)\|_{p'}=\|v\|_p^{p-1}\)，故候选导数确为有界线性泛函。

沿逐点线段积分，再用 Tonelli 与 Hölder，余项满足
\[
 |P(u+h)-P(u)-DP(u)h|
 \le \|\nabla h\|_p\sup_{0\le t\le1}
 \|A_p(\nabla u+t\nabla h)-A_p(\nabla u)\|_{p'}.
\]
上界在 \(p\ge2\) 时为 \(O((\|\nabla u\|_p+\|\nabla h\|_p)^{p-2}\|\nabla h\|_p^2)\)，在 \(p<2\) 时为 \(O(\|\nabla h\|_p^p)\)，除以 \(\|h\|_{W^{1,p}}\) 后都趋零。同一向量估计又说明 \(DP(u+h)\to DP(u)\) 于对偶算子范数，完成证明。
\end{proof}

这也辨明了第一与第二变分的差别：\(1<p<2\) 时一阶导数连续，但 \(|z|^p/p\) 的二阶导数在 \(z=0\) 附近发散；不能把二次能量的有界 Hessian 论证照搬到所有幂能量。下面的存在性只需一阶变分和凸性，不需这种额外二阶正则性。

\subsection{第二变分}
\label{b2:subsec:140104}

设 \(X\) 为实 Banach 空间，\(U\subset X\) 开，\(J:U\to\mathbb R\)，\(u\in U\)、\(h\in X\)，且 \(J\) 在 \(u\) 的邻域可微。当导数映射本身在 \(u\) 处 Fréchet 可微时，其导数 \(D^2J(u)\) 是有界双线性型。对同一方向两次求导所得的
\[\delta^2J(u;h)=D^2J(u)[h,h]\]
称为 \term[dier-bianfen]{第二变分}[second variation]。若 \(J\) 在 \(u\) 附近二阶连续 Fréchet 可微，则沿线段积分两次得到
\[
 J(u+h)-J(u)-DJ(u)h
 =\int_0^1(1-t)D^2J(u+th)[h,h]\,\mathrm dt.
\]
二阶项反映驻点附近的弯曲方向；它提供必要条件，也在具有统一正下界时提供充分条件。

\Needspace{6\baselineskip}
\begin{theorem}[第二变分判别]
\label{b2:thm:14-second-variation-test}
设 \(X\) 是实 Banach 空间，\(u\in X\)，\(V\subset X\) 是闭线性子空间，实值泛函 \(J\) 在 \(u\) 的开邻域内二阶连续 Fréchet 可微。若 \(u\) 是 \(J\) 在 \(u+V\) 上的局部极小点，则
\[DJ(u)h=0,\qquad D^2J(u)[h,h]\ge0\quad(h\in V).\]
反之，若 \(DJ(u)|_V=0\)，且存在 \(c>0\)，使
\[D^2J(u)[h,h]\ge c\lVert h\rVert_X^2\quad(h\in V),\]
则 \(u\) 是 \(J\) 在 \(u+V\) 上的严格局部极小点。
\end{theorem}
\begin{proof}
必要性由一元函数 \(t\mapsto J(u+th)\) 的一、二阶极值条件得到。证明充分性时，利用 \(D^2J\) 的算子范数连续性，取 \(r>0\)，使 \(\lVert v-u\rVert_X<r\) 时 \(\lVert D^2J(v)-D^2J(u)\rVert<c/2\)。对 \(h\in V\)、\(0<\lVert h\rVert_X<r\)，上面的积分余项公式给出
\[
 J(u+h)-J(u)\ge\int_0^1(1-t)\frac c2\lVert h\rVert_X^2\,\mathrm dt
 =\frac c4\lVert h\rVert_X^2>0.
\]
\end{proof}

仅有非负第二变分并不足以判别极小：\(t^4\) 与 \(-t^4\) 在零点的第一、第二导数都为零，却有相反的极值性质。在无限维空间里，即使对每个非零 \(h\) 都有 \(D^2J(u)[h,h]>0\)，也不能直接得到定理要求的统一常数 \(c\)。后面处理 Dirichlet 能量时，正是 Poincaré 不等式把梯度能量转成整个函数空间范数的控制，从而弥补这一差别。

% !TeX root = ../../Book2.tex
\section{Euler–Lagrange 方程}
\label{b2:sec:1402}

第一变分为零是函数空间中的驻点条件。它通常以积分等式出现；把扰动的导数移到系数上，才能读出微分方程。这个步骤也同时产生边界项，因此方程与边界条件不能分开猜测。本节先在光滑函数类中推导，下一章再把积分等式本身作为解的定义。

\subsection{一维问题}
\label{b2:subsec:140201}

设 \(a<b\)，考虑端点固定为 \(u(a)=A\)、\(u(b)=B\) 的实函数，及泛函
\[J(u)=\int_a^b F(x,u(x),u'(x))\,\mathrm dx.\]
允许扰动满足 \(h(a)=h(b)=0\)。如果 \(u\) 为极小点，\cref{b2:prop:14-gateaux-minimum,b2:prop:14-smooth-integral-variation}给出
\[\int_a^b\bigl(F_z(x,u,u')h+F_p(x,u,u')h'\bigr)\,\mathrm dx=0.\]

\Needspace{6\baselineskip}
\begin{theorem}[一维 Euler--Lagrange 方程]
\label{b2:thm:14-euler-lagrange-one}
设 \(a<b\)，实值 \(F\in C^2([a,b]\times\mathbb R^2)\)，实值 \(u\in C^2([a,b])\) 是 \(J(v)=\int_a^bF(x,v,v')\,\mathrm dx\) 在 \(v(a)=u(a)\)、\(v(b)=u(b)\) 的函数类中在 \(C^1\) 范数下的局部极小点。则
\[F_z(x,u,u')-\frac{\mathrm d}{\mathrm dx}F_p(x,u,u')=0\quad(a<x<b).\]
更一般地，对任意 \(u\in C^2([a,b])\)，该方程与所有零端点扰动的第一变分为零等价。
\end{theorem}
\begin{proof}
分部积分把第一变分写成 \(\int_a^b q(x)h(x)\,\mathrm dx\)，其中 \(q=F_z-\frac{\mathrm d}{\mathrm dx}F_p\) 连续。若某点 \(q(x_0)>0\)，连续性保证一个小区间内 \(q>0\)；在其中选非负且不恒零的光滑紧支集扰动，便使积分为正，矛盾。对 \(q(x_0)<0\) 同理。因此 \(q=0\)。反向只须把方程代回分部积分等式，边界项因扰动端点为零而消失。
\end{proof}

曲线的弧长对应 \(F=\sqrt{1+p^2}\)，于是 \(\frac{\mathrm d}{\mathrm dx}(u'/\sqrt{1+(u')^2})=0\)。由于 \(p\mapsto p/\sqrt{1+p^2}\) 严格单调，\(u'\) 为常数，驻曲线只能是连接两端点的线段。它确为全局极小曲线，可以再由 \(p\mapsto\sqrt{1+p^2}\) 的凸性证明：以线段斜率 \(m\) 处的切线作下界，积分后线性修正项为 \(\frac{m}{\sqrt{1+m^2}}\int_a^b(u'-m)=0\)。这里微分方程找出候选者，凸性才完成极小性的判断。

\ProseParagraphBreak

若 \(F\) 不显含 \(x\)，沿任何 \(C^2\) 驻曲线有
\[
 \frac{\mathrm d}{\mathrm dx}(F-u'F_p)
 =u'\left(F_z-\frac{\mathrm d}{\mathrm dx}F_p\right)=0.
\]
这给出一个一阶守恒量，常比二阶方程更易使用。但若原式包含 \(x\)，右端还要加上 \(F_x\)，不能继续断言守恒。即便有守恒量，也须重新代入端点条件；局部的一阶积分并不保证全局可解。

\subsection{多维问题}
\label{b2:subsec:140202}

多维区域中，导数向量 \(\nabla h\) 的分部积分产生散度。于是梯度型能量自然对应散度型方程，这也是弱解理论优先采用散度形式的原因。

\Needspace{6\baselineskip}
\begin{theorem}[多维 Euler--Lagrange 方程]
\label{b2:thm:14-euler-lagrange-many}
设整数 \(d\geq1\)，\(\Omega\subset\mathbb R^d\) 是有界开集，实值 \(F\in C^2(\overline\Omega\times\mathbb R\times\mathbb R^d)\)，实值 \(u\in C^2(\Omega)\cap C^1(\overline\Omega)\)。则 \(J(u)=\int_\Omega F(x,u,\nabla u)\,\mathrm dx\) 对每个 \(h\in C_c^\infty(\Omega)\) 的第一变分为零，当且仅当
\[F_z(x,u,\nabla u)-\operatorname{div}\bigl(F_p(x,u,\nabla u)\bigr)=0\quad\text{于 }\Omega.\]
特别地，任一允许这些紧支集扰动的光滑局部极小点都满足此方程。
\end{theorem}
\begin{proof}
第一变分公式对紧支集扰动成立。逐坐标分部积分，边界项为零，得到连续函数 \(F_z-\operatorname{div}F_p\) 与任意测试函数的乘积积分为零。若此函数在一点非零，在小球内选同号的非负凸起函数便产生矛盾，故函数处处为零。反向由同一分部积分恒等式立即得到。
\end{proof}

例如图像面积泛函 \(\int_\Omega\sqrt{1+|\nabla u|^2}\,\mathrm dx\) 的驻点满足
\[\operatorname{div}\frac{\nabla u}{\sqrt{1+|\nabla u|^2}}=0.\]
当梯度很小时，分母接近 \(1\)，线性化给出 \(\Delta u=0\)。这解释了调和函数为何可以描述小斜率的平衡曲面，同时也说明二次 Dirichlet 能量只是面积能量在小梯度附近的近似：大梯度下两者的增长速度和紧性要求不同。

\ProseParagraphBreak

还有一种常见误读：把 \(F_p(x,u,\nabla u)\) 看成只依赖 \(x\) 的显式系数，而漏掉链式法则。以 \(F=\tfrac12a(u)|\nabla u|^2\) 为例，方程是
\[\tfrac12a'(u)|\nabla u|^2-\operatorname{div}(a(u)\nabla u)=0,\]
若展开散度，等价于 \(a(u)\Delta u+\tfrac12a'(u)|\nabla u|^2=0\)，绝不是单纯的 \(a(u)\Delta u=0\)。在粗糙系数或弱正则性情形，保留散度形式通常更可靠。

\subsection{自然边界条件}
\label{b2:subsec:140203}

边界值若预先固定，扰动在边界为零，边界项便不参与驻点条件；边界值若可以自由变化，边界项必须自己消失，由此得到 \term[ziran-bianjie-tiaojian]{自然边界条件}[natural boundary condition]。设整数 \(d\geq1\)，\(\Omega\subset\mathbb R^d\) 为非空有界 \(C^2\) 开集，不要求连通；\(u\) 为实值且在 \(\overline\Omega\) 的邻域内为 \(C^2\)。设实值密度 \(F\) 在 \(\overline\Omega\times\mathbb R\times\mathbb R^d\) 的邻域内为 \(C^2\)，\(B,B_z\) 在 \(\partial\Omega\times\mathbb R\) 上连续。考虑边界能量
\[
 J(u)=\int_\Omega F(x,u,\nabla u)\,\mathrm dx
       +\int_{\partial\Omega}B(x,u)\,\mathrm dS.
\]
对任意实值 \(C^1\) 扰动 \(h\)（定义在闭区域的邻域内），散度公式给出
\[
 DJ(u)h=\int_\Omega(F_z-\operatorname{div}F_p)h\,\mathrm dx
       +\int_{\partial\Omega}(F_p\cdot\nu+B_z)h\,\mathrm dS.
\]
先取紧支集扰动得到内部方程，再允许任意光滑边界迹。这一边界迹自由度可具体核验。把一小块边界写成 \(x_d=\gamma(x')\)，取支集紧含于坐标块的边界函数 \(\psi(x')\)，在邻域中定义
\[h(x',x_d)=\psi(x')\eta(x_d-\gamma(x')),\qquad \eta(0)=1,\]
其中 \(\eta\) 是支集足够小的光滑截断。由于 \(\gamma\in C^2\)，此延拓至少为 \(C^1\)，已足以作第一变分；再乘切向截断并在坐标块外置零，不引入拼接跳跃。有限个坐标块的单位分解可延拓全边界上的允许迹。如果连续边界残差在一点为正，就选择该点附近非负且非零的 \(\psi\)，边界积分为正，与驻定矛盾；为负时同理。因此得到
\[F_p(x,u,\nabla u)\cdot\nu+B_z(x,u)=0\quad\text{于 }\partial\Omega.\]
这里 \(\nu\) 始终是外单位法向。

\ProseParagraphBreak

对于二次模型，不必要求密度对 \(x\) 具有上述全部二阶导数。只取实值 \(f\in C(\overline\Omega)\)、\(\beta,g\in C(\partial\Omega)\)，令 \(F=\tfrac12|p|^2-fz\)、\(B=\tfrac12\beta z^2-gz\)。直接展开二次能量并分部积分，就得到光滑驻点的必要条件
\[-\Delta u=f\quad\text{于 }\Omega,\qquad \partial_\nu u+\beta u=g\quad\text{于 }\partial\Omega.\]
当 \(\beta=0\) 时，这是 Neumann 条件。特别地，完全自由的纯 Dirichlet 能量给出零法向导数，而不是零边界值。若只在一部分边界固定函数值，则在该部分扰动为零，在其余部分产生自然条件；交界处的正则性是另一问题，形式推导不自动保证解在那里光滑。

在纯 Neumann 情形，还必须核查数据相容性。设 \(\Omega_j\) 为 \(\Omega\) 的各连通分支，\cref{b2:prop:13-neumann-compatibility}给出逐分支的必要条件
\[
 \int_{\Omega_j}f\,\mathrm{d}x
       +\int_{\partial\Omega_j}g\,\mathrm{d}S=0.
\]
也可以直接从变分看到这个约束：将 \(u\) 在 \(\Omega_j\) 上加一个常数 \(c\)，梯度能量不变，而
\[
 J(u+c\mathbf1_{\Omega_j})-J(u)
 =-c\left(\int_{\Omega_j}f\,\mathrm{d}x
                 +\int_{\partial\Omega_j}g\,\mathrm{d}S\right).
\]
这些逐分支常数方向在当前 \(C^2\) 边界的光滑变分类中可行，在 \(H^1\) 弱变分类中也可行。因此若有一项不为零，泛函沿该方向就没有驻点，且向一个方向无下界；若全部相容，能量反而对每个分支上的常数平移不变。连通区域须固定一个常数，非连通区域须分别处理各分支的常数自由度。自然边界方程描述的是驻点必须满足的条件，不能代替相容性、极小元存在性和归一化后的唯一性证明。

\ProseParagraphBreak

一维的符号尤其值得核对。对 \(\int_a^bF(x,u,u')\,\mathrm dx+B_a(u(a))+B_b(u(b))\)，两端自由时条件分别为
\[-F_p(a,u(a),u'(a))+B_a'(u(a))=0,\qquad F_p(b,u(b),u'(b))+B_b'(u(b))=0.\]
左端的负号来自外法向朝左。若遗漏它，带弹簧或外力的边界模型就会产生相反的力平衡。

\subsection{向量值问题}
\label{b2:subsec:140204}

现在固定整数 \(d,m\geq1\)、有界开集 \(\Omega\subset\mathbb R^d\)，令 \(u=(u^1,\ldots,u^m)\in C^2(\Omega;\mathbb R^m)\cap C^1(\overline\Omega;\mathbb R^m)\)，把梯度视为矩阵 \(P=(p_i^\alpha)\)，其中 \(p_i^\alpha=\partial_i u^\alpha\)。对实值密度 \(F\in C^2(\overline\Omega\times\mathbb R^m\times\mathbb R^{m\times d})\)，紧支集光滑向量扰动的第一变分为
\[
 DJ(u)h=\int_\Omega\sum_{\alpha=1}^m
 \left(F_{z^\alpha}h^\alpha+\sum_{i=1}^dF_{p_i^\alpha}\partial_i h^\alpha\right)\,\mathrm dx.
\]
各分量扰动可以独立选择。因此，若 \(u\) 光滑且对所有紧支集向量扰动驻定，则逐分量的基本引理给出系统
\[
 F_{z^\alpha}(x,u,Du)-\sum_{i=1}^d\partial_iF_{p_i^\alpha}(x,u,Du)=0,
 \qquad \alpha=1,\ldots,m.
\]
推导只使用上一证明的逐分量分部积分，无须另加各分量彼此独立的假设；它们的耦合已经包含在密度 \(F\) 中。

\ProseParagraphBreak

例如 \(F=\tfrac12\sum_\alpha|\nabla u^\alpha|^2+W(u)\) 给出 \(-\Delta u^\alpha+\partial_\alpha W(u)=0\)。若 \(W(z)=\tfrac12z^{\mathsf T}Mz\)，其中 \(M\) 为实对称矩阵，则为 \(-\Delta u+Mu=0\)。若写入一个非对称矩阵 \(M\)，二次型 \(z^{\mathsf T}Mz\) 实际只保留对称部分；相应变分绝不产生整个非对称矩阵。这一简单计算已经预示：一般线性椭圆方程虽可用双线性型求解，却未必是某个标量二次能量的极值方程。

\ProseParagraphBreak

还须区别值域约束。例如要求 \(|u|=1\) 时，可行的一阶扰动必须满足 \(u\cdot h=0\)。此时只能推出 Euler--Lagrange 向量与所有切向扰动正交，故它沿法向 \(u\)；不能把每个分量直接置零。对光滑的球面约束，乘子实际上可以算出。设整数 \(d,m\geq1\)，\(\Omega\subset\mathbb R^d\) 为非空有界开集，\(u\in C^2(\Omega;\mathbb R^m)\cap C^1(\overline\Omega;\mathbb R^m)\)、\(|u|=1\)，考虑半 Dirichlet 能量。对任意紧支集 \(h\) 且 \(u\cdot h=0\)，使用真正留在球面上的曲线
\[u_t=\frac{u+th}{|u+th|}=\frac{u+th}{\sqrt{1+t^2|h|^2}}.\]
它保持支集外的边界数据，\(\partial_tu_t|_{t=0}=h\)，故驻定给出 \(\int Du:Dh=0\)。对任意紧支集向量 \(\varphi\)，可取 \(h=\varphi-(u\cdot\varphi)u\)。分部积分后说明 \(-\Delta u\) 的切向投影为零，即 \(\Delta u=(u\cdot\Delta u)u\)。另一方面对 \(|u|^2=1\) 微分两次得到
\[0=\tfrac12\Delta|u|^2=|Du|^2+u\cdot\Delta u,\]
于是约束方程为 \(\Delta u+|Du|^2u=0\)。在一维，\(u(x)=(\cos kx,\sin kx)\) 满足 \(u''=-k^2u\)、\(|u'|^2=k^2\)，因而是该约束问题的驻曲线；它通常不是无约束 Dirichlet 问题的驻点。切向扰动确定了光滑驻曲线的约束方程；其极小性还取决于能量在其他可行曲线上的比较。

把这些约束并列，便能看清第一变分为何有时等于零、有时只得到一个方向的不等式。\Cref{b2:tab:14-admissible-variations}中的“边界迹自由”保持区域不动，只让函数的边界值变化；它不是移动区域形状的自由边界问题。约束也可叠加：例如固定迹且固定总积分时，须同时要求 \(h\in H_0^1(\Omega)\) 和 \(\int_\Omega h=0\)。

\begin{table}[htbp]
\centering
\small
\setlength{\tabcolsep}{5pt}
\caption[不同约束下的第一变分方向]{不同约束下的第一变分方向。固定非空有界 Lipschitz 开集 \(\Omega\subset\mathbb R^d\)。除球面行外，均在实空间 \(H^1(\Omega)\) 中讨论，固定迹的参照函数 \(g\in H^1(\Omega)\)，\(M\in\mathbb R\)，\(K\subset H^1(\Omega)\) 为非空闭凸集，泛函 \(J\) 在相应函数的开邻域 Fréchet 可微；经典方程仍按本节光滑积分模型理解。球面行另取 \(C^2\) 向量值积分模型、\(u\in C^2(\Omega;\mathbb R^m)\cap C^1(\overline\Omega;\mathbb R^m)\) 及紧支集 \(C^1\) 扰动，约束在每个空间点成立。各行均列必要条件。}
\label{b2:tab:14-admissible-variations}
\begin{tabularx}{\linewidth}{@{}>{\raggedright\arraybackslash}p{3.1cm}>{\raggedright\arraybackslash}X>{\raggedright\arraybackslash}p{5.7cm}@{}}
\toprule
约束 & 可行的一阶方向 & 第一变分条件 \\
\midrule
固定迹：\(u\in g+H_0^1(\Omega)\) & \(h\in H_0^1(\Omega)\)，\(u+th\) 对正负小 \(t\) 均可行 & \(DJ(u)h=0\) 对所有零迹方向成立 \\
固定积分：\(\int_\Omega u=M\) & \(h\in H^1(\Omega)\)、\(\int_\Omega h=0\)；如另固定迹，还须 \(h\in H_0^1\) & 在零积分子空间上 \(DJ(u)h=0\)，方程可留下常数乘子 \\
边界迹自由，区域固定 & \(h\in H^1(\Omega)\)，边界迹可变化 & 所有方向的第一变分为零；光滑积分模型产生自然边界条件 \\
逐点球面值：\(|u(x)|=1\) & \(u(x)\cdot h(x)=0\)；\(h\) 是可行曲线的速度 & Euler--Lagrange 向量的切向分量为零，不要求整个向量为零 \\
闭凸障碍：\(u\in K\) & \(v-u\)、\(v\in K\)，线段 \(u+t(v-u)\) 在 \(0\leq t\leq1\) 可行 & 极小点满足 \(DJ(u)(v-u)\geq0\)，这是单侧变分不等式 \\
\bottomrule
\end{tabularx}
\end{table}

固定积分行的常数乘子来自一个简单的线性事实。记 \(L(h)=\int_\Omega h\)，在所讨论的线性扰动空间内若 \(L\) 非零，可取 \(h_0\) 使 \(L(h_0)=1\)。当 \(DJ(u)\) 在 \(\ker L\) 上为零时，把 \(h\) 分成 \(h-L(h)h_0\) 与 \(L(h)h_0\)，便有 \(DJ(u)h=\lambda L(h)\)，其中 \(\lambda=DJ(u)h_0\)。

球面行的切向条件只限制曲线的初始速度；一般有 \(|u+th|^2=1+t^2|h|^2\)，所以这条直线不会留在球面上。前面使用的 \((u+th)/|u+th|\) 才是具有该速度的真实可行曲线。闭凸障碍则恰好相反：从 \(u\) 指向每个 \(v\in K\) 的单侧线段由凸性保证可行，除以正参数后得到不等式；只有另加 \(J\) 的凸性时，这组不等式才同时足以保证全局极小。

% !TeX root = ../../Book2.tex
\section{能量泛函}
\label{b2:sec:1403}

同一个偏微分方程可以通过能量的驻点条件出现，但“能量”一词并不总表示待极小化的非负量。椭圆方程常对应空间能量的极小值；波动方程对应时空作用量的驻点，同时另有随时间守恒的正能量。本节用这两个模型把几种不同作用区分开。

\subsection{Dirichlet 能量}
\label{b2:subsec:140301}

设 \(\Omega\subset\mathbb R^d\) 是有界开集，记 \(H_0^1(\Omega)=W_0^{1,2}(\Omega)\)。第一卷定义 22.2.1 将它定义为 \(C_c^\infty(\Omega)\) 在 \(H^1\) 中的闭包，这一定义不要求边界光滑。由第一卷定理 24.1.2 的零边界 Poincaré 不等式，\(\lVert v\rVert_2\le C_P\lVert\nabla v\rVert_2\)，所以 \(\lVert\nabla v\rVert_2\) 是这个空间上的等价 Hilbert 范数。

\ProseParagraphBreak

设 \(\Omega\subset\mathbb R^d\) 是有界开集，给定实值 \(f\in L^2(\Omega)\) 和实值 \(g\in H^1(\Omega)\)，把允许函数写为 \(g+H_0^1(\Omega)\)，考虑 \term[dirichlet-nengliang]{Dirichlet 能量}[Dirichlet energy]及线性外力项组成的泛函
\[
 E(u)=\frac12\int_\Omega|\nabla u|^2\,\mathrm dx-\int_\Omega fu\,\mathrm dx.
\]
仿射空间的写法精确表达“\(u\) 与 \(g\) 具有相同边界数据”：即使边界粗糙到尚未定义逐点边界迹，也可以使用它。

\Needspace{6\baselineskip}
\begin{proposition}[二次能量的驻点与极小点]
\label{b2:prop:14-dirichlet-minimum}
设整数 \(d\geq1\)，\(\Omega\subset\mathbb R^d\) 为非空有界开集，\(f\in L^2(\Omega;\mathbb R)\)、\(g\in H^1(\Omega;\mathbb R)\)，\(E(v)=\frac12\int_\Omega|\nabla v|^2-\int_\Omega fv\)。\(u\in g+H_0^1(\Omega;\mathbb R)\) 是 \(E\) 在该仿射空间上的极小点，当且仅当
\[
 \int_\Omega\nabla u\cdot\nabla h\,\mathrm dx=\int_\Omega fh\,\mathrm dx
 \qquad(h\in H_0^1(\Omega)).
\]
这样的极小点至多有一个，且对每个 \(h\in H_0^1(\Omega)\) 都有
\[E(u+h)-E(u)=\tfrac12\lVert\nabla h\rVert_2^2.\]
\end{proposition}
\begin{proof}
Cauchy--Schwarz 不等式保证所有积分有限。直接展开平方得
\[
 E(u+h)-E(u)=\int_\Omega\nabla u\cdot\nabla h-\int_\Omega fh
              +\tfrac12\lVert\nabla h\rVert_2^2.
\]
其中前两项构成连续线性泛函，最后一项是二阶余项，故一阶导数正是前两项。极小点的双侧变分为零，得到所列积分等式。反之，积分等式使能量差等于非负的梯度平方，所以 \(u\) 为全局极小点。若另一点 \(u+h\) 也极小，则 \(\nabla h=0\)；Poincaré 不等式推出 \(h=0\)。
\end{proof}

这一命题暂时没有断言极小点存在；它说明一旦构造出驻点，就无须再排查鞍点。二次展开还给出误差的精确解释：任何满足同一边界数据的近似函数，其能量超额恰为梯度误差平方的一半。与一般的第二变分局部判别相比，这是全局恒等式，原因在于能量严格二次且对允许方向一致凸。

\ProseParagraphBreak

在区间 \((0,1)\) 上取 \(f=1\)、\(g=0\)。函数 \(u(x)=x(1-x)/2\) 满足 \(-u''=1\) 和零端点条件；对光滑紧支集扰动分部积分，再用密度推广到 \(H_0^1\)，即可验证积分等式。因此它是唯一极小点，能量为 \(-1/24\)。负能量并不矛盾：梯度项非负，外力项却可以降低总泛函值；强制性要求的是大范数下的增长，而不是每一点的非负性。

\subsection{波动作用量与守恒能量}
\label{b2:subsec:140302}

设整数 \(d\geq1\)，\(\Omega\subset\mathbb R^d\) 为非空有界 \(C^1\) 开集，\(c,T>0\)，并考虑在 \([0,T]\times\overline\Omega\) 的开邻域内为 \(C^2\)、在侧边界取零值的实函数 \(u(t,x)\)。时空作用量为
\[
 \mathcal A(u)=\frac12\int_0^T\int_\Omega
       \bigl(|\partial_tu|^2-c^2|\nabla_xu|^2\bigr)\,\mathrm dx\,\mathrm dt.
\]
对在初、末时刻和空间边界均为零的扰动 \(h\)，时、空分别分部积分，得到
\[
 D\mathcal A(u)h=\int_0^T\int_\Omega
       (-\partial_t^2u+c^2\Delta_xu)h\,\mathrm dx\,\mathrm dt.
\]
所以驻点满足波动方程 \(\partial_t^2u-c^2\Delta u=0\)。这里固定初末值是推导变分方程时对扰动的限制；实际波动初值问题仍由初位移和初速度确定，不能把两种问题混为一谈。

\ProseParagraphBreak

作用量的二次型有正负两部分，一般并不在解处取得极小值。例如取乘积扰动 \(h(t,x)=a(t)b(x)\)，其中两因子均光滑紧支集，则
\[
 2\mathcal A(h)=\lVert a'\rVert_2^2\lVert b\rVert_2^2
      -c^2\lVert a\rVert_2^2\lVert\nabla b\rVert_2^2.
\]
固定非零 \(b\) 后，取 \(a_N(t)=\eta(t)\sin(Nt)\)，其中 \(0\ne\eta\in C_c^\infty(0,T)\)。展开正弦平方并用\cref{b2:thm:03-riemann-lebesgue}，有 \(\|a_N\|_2^2\to\|\eta\|_2^2/2\) 及 \(N^{-2}\|a_N'\|_2^2\to\|\eta\|_2^2/2\)，所以作用量为正。反之，固定非零 \(a\)，取 \(b_N(x)=\psi(x)\sin(Nx_1)\)，其中 \(0\ne\psi\in C_c^\infty(\Omega)\)。相同的平方展开给出 \(\|b_N\|_2^2\to\|\psi\|_2^2/2\) 及 \(N^{-2}\|\nabla b_N\|_2^2\to\|\psi\|_2^2/2\)，所以作用量为负。若 \(u\) 为驻点，精确二次展开给出 \(\mathcal A(u+th)-\mathcal A(u)=t^2\mathcal A(h)\)。故可把两类扰动都缩得任意小，在每个驻点附近仍有增加和降低作用量的方向。

\Needspace{6\baselineskip}
\begin{proposition}[波动能量守恒]
\label{b2:prop:14-wave-energy}
设整数 \(d\geq1\)，\(c,T>0\)，\(\Omega\subset\mathbb R^d\) 为非空有界 \(C^1\) 开集，实函数 \(u\) 在 \([0,T]\times\overline\Omega\) 的某开邻域内属于 \(C^2\)，满足 \(u_{tt}-c^2\Delta u=0\) 于 \((0,T)\times\Omega\)，且 \(u(t,y)=0\) 对 \(0\leq t\leq T\)、\(y\in\partial\Omega\) 成立。则
\[
 \mathcal E(t)=\frac12\int_\Omega
          \bigl(|u_t(t,x)|^2+c^2|\nabla u(t,x)|^2\bigr)\,\mathrm dx
\]
与时间无关。
\end{proposition}
\begin{proof}
对时间求导并对空间分部积分，得
\[
 \mathcal E'(t)=\int_\Omega u_t(u_{tt}-c^2\Delta u)\,\mathrm dx
           +c^2\int_{\partial\Omega}u_t\partial_\nu u\,\mathrm dS=0.
\]
内部项由方程为零，边界项由对恒定零边界条件求时间导数所得 \(u_t=0\) 而消失。
\end{proof}

守恒能量中空间项带正号，作用量中带负号。前者衡量一个时刻的状态，后者比较整段时空运动。对同样的区域、实值 \(C^2\) 解及恒定零边界，若方程改为 \(u_{tt}-c^2\Delta u+\gamma u_t=f\)，其中 \(\gamma\geq0\)、\(f\) 连续，则同一计算给出
\[
 \mathcal E'(t)=\int_\Omega f(t,x)u_t(t,x)\,\mathrm{d}x
                    -\gamma\|u_t(t)\|_2^2.
\]
这些等式将外界做功和耗散分开记录。\cref{b2:tab:14-action-energy}集中比较作用量与能量；表中的守恒结论始终包含零边界且无外力、无阻尼的条件。

\begin{table}[htbp]
\centering
\caption[波动作用量与守恒能量的区别]{波动作用量与守恒能量的区别。取整数 \(d\geq1\)、\(c,T>0\)、非空有界 \(C^1\) 开集 \(\Omega\subset\mathbb R^d\)，函数实值且光滑到足以完成正文的微分与分部积分，侧边界为恒定零值。作用量列的变分还固定初、末时刻的函数值；能量列讨论满足波动方程的实际演化。}
\label{b2:tab:14-action-energy}
\small
\begin{tabularx}{\textwidth}{@{}p{5.2em}XX@{}}
\toprule
比较项目 & 时空作用量 \(\mathcal A\) & 守恒能量 \(\mathcal E(t)\)\\
\midrule
积分范围 & 整段 \((0,T)\times\Omega\) & 固定时刻的 \(\Omega\)\\
空间梯度项 & \(-c^2|\nabla u|^2/2\) & \(+c^2|\nabla u|^2/2\)\\
数学作用 & 对允许扰动驻定，推出运动方程 & 控制速度与梯度，进而估计解的差\\
极值或演化 & 驻点一般为鞍点 & 无外力、无阻尼时随时间守恒\\
\bottomrule
\end{tabularx}
\end{table}

\subsection{PDE 与极值问题}
\label{b2:subsec:140303}

从微分方程寻找泛函时，最容易检验的是对称性。设 \(H\) 是实 Hilbert 空间，\(a:H\times H\to\mathbb R\) 是有界双线性型，\(\ell\in H^*\)。令
\[J(u)=\tfrac12a(u,u)-\ell(u).\]
直接展开可得
\[DJ(u)h=\tfrac12\bigl(a(u,h)+a(h,u)\bigr)-\ell(h).\]
因此仅当 \(a\) 对称时，驻点条件才是 \(a(u,h)=\ell(h)\)。反对称部分在 \(a(u,u)\) 中完全消失，却可能实质影响方程的解。

\ProseParagraphBreak

例如 \(H=\mathbb R^2\)，令 \(Au=(u_1-u_2,u_1+u_2)\)，\(a(u,h)=(Au)\cdot h\)。尽管 \(a(u,u)=|u|^2\) 严格为正，极小化 \(\tfrac12|u|^2-f\cdot u\) 得到的是 \(u=f\)，而原方程是 \(Au=f\)。两者一般不同。下一章的 Lax--Milgram 定理直接求解有界强制双线性型，无须先把它变成能量；这正是从变分极小化过渡到一般弱解方法时必须增加的能力。

\ProseParagraphBreak

另一个区别是驻点与动力学。若 \(E(u)=\tfrac12\int|\nabla u|^2\)，形式上的 \(L^2\) 梯度为 \(-\Delta u\)。沿最陡下降运动 \(u_t=-\operatorname{grad}E(u)\)，得到热方程 \(u_t=\Delta u\)，而平衡点满足 Laplace 方程。对光滑零边界解，分部积分给出
\[\frac{\mathrm d}{\mathrm dt}E(u(t))=-\int_\Omega|u_t|^2\,\mathrm dx\le0.\]
这项恒等式说明热演化趋向降低能量，却不在这里自动证明长时间收敛；还需要轨道紧性、极限点的识别等信息。变分结构可以提供估计和候选极限，但不能代替这些额外步骤。

% !TeX root = ../../Book2.tex
\section{直接方法}
\label{b2:sec:1404}

先推导 Euler--Lagrange 方程再求解，有时会把存在性问题变得更难。直接方法反过来考察能量本身：取能量趋近下确界的序列，控制它的大小，抽取极限，再证明极限没有丢失边界条件和能量下界。第一卷定理 9.5.9 已在自反 Banach 空间中证明这一抽象原则。对积分能量而言，关键是把能量有界转为函数空间中的有界，再判断弱极限能否保留约束并降低能量。

\subsection{强制性}
\label{b2:subsec:140401}

在 Banach 空间 \(X\) 的约束集 \(K\) 上，泛函的强制性指
\[J(v)\longrightarrow+\infty\qquad(v\in K,\ \lVert v\rVert_X\to\infty).\]
它保证能量有界的候选者不能逃到范数无穷远处。先取整数 \(d\geq1\)、非空有界开集 \(\Omega\subset\mathbb R^d\) 及实值 \(f\in L^2(\Omega)\)。第一卷定理24.1.2的零边界 Poincaré 不等式说明：只要开集包含于有限宽条带，就有 \(\|v\|_2\leq C_P\|\nabla v\|_2\) 对全部 \(v\in H_0^1(\Omega)\) 成立。有界开集满足这个几何条件，不需要连通性或边界光滑性。因此 Cauchy--Schwarz 不等式给出
\[
 \frac12\lVert\nabla v\rVert_2^2-\int_\Omega fv
 \ge\frac12\lVert\nabla v\rVert_2^2-C_P\lVert f\rVert_2\lVert\nabla v\rVert_2.
\]
右端是开口向上的二次式；又因梯度范数在 \(H_0^1\) 上与完整的 \(H^1\) 范数等价，故这项能量强制。若固定实值 \(g\in H^1(\Omega)\) 并允许边界类 \(g+H_0^1(\Omega)\)，写成 \(u=g+v\)，交叉项 \(\int\nabla g\cdot\nabla v\) 仍只按 \(\lVert\nabla v\rVert_2\) 线性增长，因此同样得到强制性。

\ProseParagraphBreak

边界限制不能从这个论证中删除。在同一有界开集的实空间 \(H^1(\Omega)\) 中，常函数的梯度能量为零，故单靠梯度无法控制范数。若 \(\int_\Omega f\ne0\)，沿适当符号的常函数甚至能让总能量趋向 \(-\infty\)。齐次 Neumann 问题因而至少要求 \(\int_\Omega f=0\)。对有界 \(C^1\) 开集、连续源项及连续边界通量的经典问题，\cref{b2:prop:13-neumann-compatibility}还证明了总源量加总外通量为零的必要条件。相容条件排除能量沿常数方向下降，却还没有固定常数自由度。

\ProseParagraphBreak

对于非空有界连通 Lipschitz 区域 \(\Omega\subset\mathbb R^d\)，第一卷定理24.1.3已证明 Poincaré--Wirtinger 不等式；在本处的 \(p=2\) 情形，它写成
\[
 u_\Omega=\frac1{|\Omega|}\int_\Omega u,\qquad
 \|u-u_\Omega\|_{L^2(\Omega)}
 \leq C_\Omega\|\nabla u\|_{L^2(\Omega)}
 \quad(u\in H^1(\Omega)).
\]
所以在零均值子空间中，梯度可以控制完整范数。连通性在这里不可删除：在 \((0,1)\cup(2,3)\) 上分别取常数 \(1\)、\(-1\)，所得函数总均值为零、弱梯度为零，\(L^2\) 范数却为 \(\sqrt2\)。不连通时应在每个连通分支分别处理相容条件、减去该分支的均值；应用上述估计还要求各分支满足相应的几何条件，若分支无限多，不能未经说明就把各自常数当成统一常数。即使开集连通，也不能只凭有界性调用这一版本。Neumann 零均值估计使用上述 Lipschitz 与连通条件；零边界 \(W_0^{1,p}\) 的存在性则适用于任意非空有界开集。

\subsection{极小化序列}
\label{b2:subsec:140402}

设 \(K\) 为非空集合、\(J:K\to(-\infty,+\infty]\)。若 \(m=\inf_KJ\) 是有限实数，可以按下确界的定义选取 \(u_n\in K\)，满足 \(m\le J(u_n)<m+1/n\)。这样的序列称为极小化序列。选取过程既不要求驻点已存在，也不要求每个 \(u_n\) 近似满足 Euler--Lagrange 方程。

\ProseParagraphBreak

控制能量后，强制性使序列有界，但有界不等于强收敛。无限维 Hilbert 空间中的标准正交序列有界，却没有强收敛子列；它仅弱收敛到零。这就是直接方法引入弱拓扑的实质理由。另一方面，若约束集不是弱闭的，弱极限可能不再可行；若能量不是弱下半连续的，极限能量可能突然升高。序列有界、约束保留与能量比较分别承担不同任务，不能用一句“由紧性”全部略去。

\ProseParagraphBreak

二次能量还有一个可直接看见的特殊现象。取非空有界开集 \(\Omega\subset\mathbb R^d\)、实值 \(f\in L^2(\Omega)\) 与 \(g\in H^1(\Omega)\)，记 \(E(u)=\tfrac12\int_\Omega|\nabla u|^2-\int_\Omega fu\)、\(m=\inf_{g+H_0^1(\Omega)}E\)。若 \(u_n,u_k\in g+H_0^1(\Omega)\)，则
\[
 E\left(\frac{u_n+u_k}{2}\right)
 =\frac{E(u_n)+E(u_k)}2-\frac18\lVert\nabla(u_n-u_k)\rVert_2^2.
\]
左端不小于下确界 \(m\)。因此极小化序列事实上在梯度范数下为 Cauchy 序列，借完备性可以直接得到强极限。一般非二次能量不一定有这个精确恒等式，弱紧性的方法因而更具普遍性。

\subsection{弱紧性}
\label{b2:subsec:140403}

设整数 \(d\geq1\)、\(\Omega\subset\mathbb R^d\) 为非空开集、\(1<p<\infty\)，函数取实值。映射
\[
 T:W^{1,p}(\Omega)\longrightarrow(L^p(\Omega))^{d+1},\qquad
 Tu=(u,\partial_1u,\ldots,\partial_du)
\]
把 Sobolev 空间识别为一个闭子空间。这里要区分空间结构与范数的数值。第一卷采用逐分量范数
\[
 N_p(u)=\left(\|u\|_p^p+\sum_{i=1}^d\|\partial_i u\|_p^p\right)^{1/p};
\]
它正好与乘积空间的 \(\ell^p\) 范数对应。能量计算常使用欧氏梯度，得到另一个范数
\[
 G_p(u)=\left(\|u\|_p^p+\int_\Omega|\nabla u|_2^p\right)^{1/p}.
\]
令 \(\gamma=|1/p-1/2|\)，有限维向量的范数比较逐点给出
\[
 d^{-\gamma}N_p(u)\leq G_p(u)\leq d^\gamma N_p(u).
\]
因此两种范数等价；\(p=2\) 时它们相同，一般 \(p\) 时不能把欧氏梯度能量范数与逐分量乘积范数默认当作等距。本处的自反性论证只需要这种等价性。

\ProseParagraphBreak

像的闭性也包含一个需要保留的导数关系。若 \(Tu_n\to(v_0,v_1,\ldots,v_d)\) 于乘积范数，则对每个 \(\varphi\in C_c^\infty(\Omega)\)，Hölder 不等式允许在
\(\int u_n\partial_i\varphi=-\int(\partial_i u_n)\varphi\) 中取极限，得到
\(\int v_0\partial_i\varphi=-\int v_i\varphi\)。所以 \(v_i=\partial_i v_0\) 是弱导数，极限仍属于像。有限个自反 \(L^p\) 空间的乘积及其闭子空间仍自反，故第一卷定理9.4.7给出：每个有界 \(W^{1,p}(\Omega)\) 序列都有在该空间弱收敛的子列。沿同一子列，连续的分量映射给出
\[
 u_n\rightharpoonup u\ \text{于 }L^p(\Omega),\qquad
 \partial_i u_n\rightharpoonup\partial_i u\ \text{于 }L^p(\Omega)
 \quad(1\leq i\leq d).
\]
反过来，有限个分量的弱收敛也能识别 \(W^{1,p}\) 中的弱收敛：闭子空间上的有界线性泛函可由 Hahn--Banach 定理延拓到乘积空间，再分解为对这些分量的有限次测试。

\ProseParagraphBreak

对非空开集 \(\Omega\subset\mathbb R^d\)、\(1<p<\infty\) 与实值 \(g\in W^{1,p}(\Omega)\)，仿射边界类 \(g+W_0^{1,p}(\Omega)\) 是强闭凸集，因而弱闭。这一弱闭性来自第一卷定理9.5.2的闭凸集结论。若另给实值可测障碍 \(\psi\)，约束 \(u\ge\psi\) 几乎处处也确定一个强闭凸集：强收敛序列可抽取进一步逐点几乎处处收敛的子列，从而保留不等式；闭凸性再给出弱闭性。该障碍类还须非空，才能用来寻找极小元。

\ProseParagraphBreak

端点 \(p=1\) 与 \(p=\infty\) 不在上述自反性结论内。例如在 \((0,1)\) 上，\(u_n(x)=\min\{nx,1\}\) 的 \(W^{1,1}\) 范数有界，但导数 \(n\mathbf1_{(0,1/n)}\) 集中到越来越短的边界层中。若某个子列弱收敛到 \(v\in L^1\)，对任意 \(0<\delta<1\) 用 \(\mathbf1_{(\delta,1)}\) 测试，得到 \(\int_\delta^1v=0\)，从而 \(v=0\) 几乎处处；用常函数 \(1\) 测试却应有 \(\int_0^1v=1\)，矛盾。因此导数没有 \(L^1\) 弱收敛子列。函数本身在内部趋于常数 \(1\)，左端数据已丢失。处理线性增长能量时往往需要有界变差函数和测度导数；不能把 \(p>1\) 的论证直接搬到这个端点。

\subsection{弱下半连续性}
\label{b2:subsec:140404}

第一卷定理9.5.8说明：赋范空间上凸泛函的范数下半连续性等价于弱下半连续性。对积分泛函，可以先证强收敛下的下半连续，再利用凸性改变拓扑。不过在取积分之前，还须确认被积函数可测。

\ProseParagraphBreak

具体地，设整数 \(d\geq1\)、\(\Omega\subset\mathbb R^d\) 为非空开集、\(1\leq p<\infty\)，并设
\[
 F:\Omega\times\mathbb R\times\mathbb R^d\longrightarrow[0,\infty)
\]
对每个固定 \((z,\xi)\)，函数 \(x\mapsto F(x,z,\xi)\) 为 Lebesgue 可测；对每个固定 \(x\)，函数 \((z,\xi)\mapsto F(x,z,\xi)\) 连续且凸。取实值 \(u\in W^{1,p}(\Omega)\)，为 \(u\) 及全部一阶弱导数选定有限值可测代表，在共同的零测集上补为零，记
\(V(x)=(u(x),\partial_1u(x),\ldots,\partial_du(x))\)。将 \(V\) 的每个分量先截断到 \([-n,n]\)，再向下取最近的 \(2^{-n}\) 有理网格点，得到只取有限个有理向量值的简单函数 \(V_n\)。每个取值集合都由可测分量的区间条件构成，故可写成
\[
 V_n=\sum_{q\in\mathcal Q_n}q\,\mathbf1_{A_{n,q}},\qquad
 \mathcal Q_n\subset\mathbb Q^{d+1}\ \text{有限},
\]
其中 \((A_{n,q})_{q\in\mathcal Q_n}\) 是 \(\Omega\) 的可测划分，且 \(V_n(x)\to V(x)\)。于是
\[
 F(x,V_n(x))=\sum_{q\in\mathcal Q_n}
                  \mathbf1_{A_{n,q}}(x)F(x,q)
\]
可测；再由后几个变量的连续性取逐点极限，得到 \(F(x,u(x),\nabla u(x))\) 可测。更换代表只改变一个零测集上的值，故
\[
 I(u)=\int_\Omega F(x,u,\nabla u)\,\mathrm dx\in[0,+\infty]
\]
是良定义的泛函。

\ProseParagraphBreak

下面核验下半连续性。若 \(u_n\to u\) 于 \(W^{1,p}(\Omega)\)，当 \(\liminf I(u_n)=+\infty\) 时不等式自动成立；否则先抽取实现该有限下极限的子列，再从中抽取函数与全部一阶导数同时几乎处处收敛的子列。这里导数只有有限个，可以逐次抽取而使用同一个最终子列。连续性给出被积函数几乎处处收敛，非负性使 Fatou 引理可用，因此
\[
 I(u)\leq\liminf_{n\to\infty}I(u_n).
\]
这证明强下半连续性。对任意 \(0<t<1\)，密度的凸性逐点给出
\(I((1-t)u+tv)\leq(1-t)I(u)+tI(v)\)，所以 \(I\) 凸；第一卷定理9.5.8遂给出弱下半连续性。这一段的下半连续性核验也适用于 \(p=1\)，而取得弱子列的自反性步骤另外要求 \(1<p<\infty\)。

\ProseParagraphBreak

上述密度的非负性可放宽为 \(F(x,z,\xi)\geq-a(x)\)，其中 \(a\in L^1(\Omega)\) 与 \((z,\xi)\) 无关；对 \(F+a\) 应用刚才的证明，再减去有限常数 \(\int a\) 即可。若 \(\ell\) 是同一实 Sobolev 空间上的有界线性泛函，线性项在弱收敛下连续，所以 \(I-\ell\) 仍弱下半连续。这个证明覆盖了 \(\tfrac1p\int|\nabla u|^p\) 及许多非负凸势能，却没有覆盖任意非凸密度。特别是向量值问题中，关于梯度的普通凸性虽为充分条件，却远非最一般的下半连续性条件。

\ProseParagraphBreak

以 \(\ell^2\) 上的 \(J(x)=-\lVert x\rVert^2\) 为例，\(J\) 范数连续，但对标准正交序列 \(e_n\rightharpoonup0\) 有 \(J(0)=0>-1=\lim J(e_n)\)。所以范数连续本身不能代替弱下半连续；前面的凸性在改变拓扑时起了决定作用。

\subsection{极小元的存在性}
\label{b2:subsec:140405}

下面把这些核验汇集到一个非线性模型。它既包含 Dirichlet 能量，也说明直接方法并不依赖显式解或线性叠加。证明的次序见\cref{b2:fig:14-direct-method}：强制性先控制极小化序列，自反性提供弱子列；约束弱闭性使极限可行，弱下半连续性使极限达到下确界。这四项解决存在性。严格凸性在取得极小元之后另行排除第二个极小元，不能代替其中任何一步。

% 本书原创：直接方法的存在性链条与严格凸性的独立职责。
\begin{figure}[htbp]
\centering
\begin{tikzpicture}[>=stealth,line width=.7pt,
 every node/.style={font=\small},
 directstep/.style={draw=AnalysisBlue,fill=AnalysisBlue!6,rounded corners=2pt,
              align=center,text width=6cm,inner sep=4pt,minimum height=1.05cm}]
\node[directstep] (sequence) at (0,0)
 {极小化序列\\\(u_n\in K,\quad J(u_n)\to m\)};
\node[directstep] (bounded) at (0,-1.5)
 {强制性\\\((u_n)\) 在 \(X\) 中有界};
\node[directstep] (weak) at (0,-3)
 {自反性与弱紧性\\\(u_{n_k}\rightharpoonup u\)};
\node[directstep] (limit) at (0,-4.8)
 {约束弱闭与能量弱下半连续\\\(u\in K\)\\\(J(u)\leq\liminf_kJ(u_{n_k})=m\)};
\node[directstep] (exists) at (0,-6.6)
 {极小元存在\\\(J(u)=m\)};
\node[directstep,text width=3.1cm] (unique) at (5.5,-6.6)
 {另加严格凸性\\极小元唯一};
\draw[->] (sequence)--(bounded);
\draw[->] (bounded)--(weak);
\draw[->] (weak)--(limit);
\draw[->] (limit)--(exists);
\draw[->] (exists)--(unique);
\end{tikzpicture}
\caption[直接方法的存在性步骤与唯一性条件]{直接方法的存在性步骤与唯一性条件。图中 \(X\) 为实自反 Banach 空间，\(K\subset X\) 非空且弱闭，\(J:K\to(-\infty,+\infty]\) 强制、弱下半连续且至少在一点有限，\(m=\inf_KJ\in\mathbb R\)。左列得到极小元；右下分支另要求 \(K\) 凸、\(J\) 在其有限能量点之间严格凸，才得到唯一性。}
\label{b2:fig:14-direct-method}
\end{figure}


\Needspace{6\baselineskip}
\begin{theorem}[梯度幂能量的极小元]
\label{b2:thm:14-p-energy-minimum}
设整数 \(d\geq1\)，\(\Omega\subset\mathbb R^d\) 是非空有界开集，\(1<p<\infty\)，\(g\in W^{1,p}(\Omega;\mathbb R)\)，\(f\in L^{p'}(\Omega;\mathbb R)\)，其中 \(p'=p/(p-1)\)。则
\[
 J(u)=\frac1p\int_\Omega|\nabla u|^p\,\mathrm dx-\int_\Omega fu\,\mathrm dx
\]
在 \(g+W_0^{1,p}(\Omega;\mathbb R)\) 上有唯一极小元 \(u\)。这里 \(W_0^{1,p}\) 指 \(C_c^\infty(\Omega;\mathbb R)\) 在 \(W^{1,p}\) 范数下的闭包，不要求边界存在经典迹。极小元满足
\[
 \int_\Omega|\nabla u|^{p-2}\nabla u\cdot\nabla h\,\mathrm dx
     =\int_\Omega fh\,\mathrm dx\qquad(h\in W_0^{1,p}(\Omega;\mathbb R)),
\]
其中在 \(\nabla u=0\) 处约定 \(|\nabla u|^{p-2}\nabla u=0\)。
\end{theorem}
\begin{proof}
令 \(u=g+v\)。第一卷定理24.1.2适用于有界开集上的 \(v\in W_0^{1,p}(\Omega)\)，故 \(\|v\|_p\leq C_P\|\nabla v\|_p\)。由三角不等式，\(\lVert\nabla v\rVert_p\le\lVert\nabla u\rVert_p+\lVert\nabla g\rVert_p\)，且
\[\left|\int_\Omega fu\right|\le\lVert f\rVert_{p'}\bigl(\lVert g\rVert_p+C_P\lVert\nabla v\rVert_p\bigr).\]
记 \(t=\|\nabla u\|_p\)，\(A=C_P\|f\|_{p'}\)、\(B=\|f\|_{p'}(\|g\|_p+C_P\|\nabla g\|_p)\)，则
\[
 J(g+v)\geq\frac1p t^p-At-B.
\]
右端在 \(t\geq0\) 上有有限下界，且因 \(p>1\) 而趋向 \(+\infty\)。若 \(\|v\|_{W^{1,p}}\to\infty\)，Poincaré 不等式先给出 \(\|\nabla v\|_p\to\infty\)，再由三角不等式得到 \(t\to\infty\)。所以 \(J\) 在边界类上强制，而且 \(J(g)\) 有限，下确界 \(m\) 是有限实数。

取极小化序列 \(u_n\in g+W_0^{1,p}(\Omega)\)。由强制性，\((u_n)\) 在 \(W^{1,p}\) 中有界；自反性给出弱子列 \(u_{n_k}\rightharpoonup u\)。可行集是非空强闭凸集，因而弱闭，所以 \(u-g\in W_0^{1,p}\)。上一小节以密度 \(F(x,z,\xi)=|\xi|^p/p\) 核验了梯度幂积分的弱下半连续性，而 \(f\) 的配对是有界线性泛函，故弱连续。于是
\[
 m\leq J(u)\leq\liminf_{k\to\infty}J(u_{n_k})=m.
\]
这就得到极小元，也具体实现了第一卷定理9.5.9的直接法。

再证唯一性。\(\xi\mapsto|\xi|^p\) 严格凸：若两个向量长度不同，三角不等式结合 \(s\mapsto s^p\) 的严格凸性给出中点严格不等式；若长度相同而向量不同，则欧氏三角不等式本身严格。若 \(u_1,u_2\) 都极小，而其梯度在正测度集合上不同，将这一不等式积分，中点的梯度能量就严格更小；外力项是线性的，中点仍可行，故与极小性矛盾。所以梯度几乎处处相等；差属于 \(W_0^{1,p}\)，零边界 Poincaré 不等式推出差为零。

最后固定 \(h\in W_0^{1,p}(\Omega;\mathbb R)\)。有限维函数 \(\xi\mapsto|\xi|^p/p\) 在 \(p>1\) 时为 \(C^1\)，梯度是在原点取零的 \(|\xi|^{p-2}\xi\)。沿线段 \(u+th\) 求导时，\(|t|\le1\) 的导数被
\[C_p\bigl(|\nabla u|^{p-1}+|\nabla h|^{p-1}\bigr)|\nabla h|\]
支配；该函数由 Hölder 不等式可积。支配收敛允许交换求导与积分，得第一变分公式。极小点的正负方向都可行，故第一变分为零，即所列等式。
\end{proof}

当 \(p=2\) 时，这就是 Poisson 方程的弱形式。一般 \(p\) 则对应 \(-\operatorname{div}(|\nabla u|^{p-2}\nabla u)=f\)，通常称为 \(p\)-Laplace 方程。所得极小元首先属于能量空间；若要把弱方程逐点展开，还需进一步取得经典导数。

\ProseParagraphBreak

对非空有界开集 \(\Omega\subset\mathbb R^d\)、\(1<p<\infty\)、实值 \(g\in W^{1,p}(\Omega)\)、\(f\in L^{p'}(\Omega)\)，若再给定连续凸函数 \(\Psi:\mathbb R\to[0,\infty)\)，则可在同一 \(g+W_0^{1,p}\) 类中加入势能 \(\int_\Omega\Psi(u)\)。只要某个可行函数的总能量有限，同样的直接方法仍给出极小元；非负势能保留强制性，前述可测性与 Fatou 论证给出弱下半连续性，严格凸的梯度项保留唯一性。但要写出含 \(\Psi'(u)\) 的弱方程，还须补充可微性和增长条件，确保扰动可积并允许积分下求导。存在性所需的下半连续性通常比推导经典微分方程所需的光滑性弱得多，这也是直接方法的重要优势。


非凸密度的困难可以在区间上看见。令 \(w_n\) 在每个长度 \(1/n\) 的小区间上先以斜率 \(1\) 上升半个区间，再以斜率 \(-1\) 下降，且 \(w_n(0)=0\)。则 \(w_n\in W_0^{1,4}(0,1)\)、\(\|w_n\|_\infty\le1/(2n)\)，而 \(|w_n'|=1\) 几乎处处。对 \(C^1\) 测试函数 \(\psi\)，在各小区间减去左端的 \(\psi\) 值，利用 \(w_n'\) 的零均值，得到
\[\left|\int_0^1w_n'\psi\right|\le\frac{\|\psi'\|_\infty}{n}.\]
\(\|w_n'\|_4=1\) 与 \(C^1\) 在 \(L^{4/3}\) 中的稠密性把结论推广到全部对偶测试，所以 \(w_n\rightharpoonup0\) 于 \(W^{1,4}\)。但密度 \(F(p)=(p^2-1)^2\) 给出
\[\int_0^1F(w_n')=0,\qquad \int_0^1F(0)=1.\]
该泛函因而不弱下半连续。范数有界、边界值保留，甚至函数一致收敛，都不足以保留这项能量；弱极限把快速交替的两种斜率平均成零。先以斜率一上升、再以斜率负一下降的帐篷函数仍能取得零能量，因此这个泛函有极小元；失败的是极小化序列的弱极限保持最小能量这一步。前面凸密度的不等式恰好阻止这种能量降落。

% !TeX root = ../../Book2.tex
\OptionalSection{山路定理与 Palais–Smale 条件}{b2:sec:1405}

直接方法寻找能量最低的状态。然而，一项能量可能沿某些方向趋于负无穷，同时在零点附近先升高；这时全局极小元不存在，方程却仍可能有非零解。一个一维图像是
\[
 J(t)=\frac12t^2-\frac14t^4.
\]
零点是局部极小点，\(J(2)<0\)，从 \(0\) 连续走到 \(2\) 必须经过 \(t=1\)，在那里能量达到 \(1/4\)，导数又恰好为零。\cref{b2:fig:14-mountain-profile}画出了曲线先上升、再下降到零以下的一段。函数空间中可以沿很多不同路径绕行；我们希望把各条路径的最高能量尽量压低，找到仍然无法绕开的临界高度。

\begin{figure}[htbp]
\centering
\begin{tikzpicture}[x=4cm,y=4.2cm,font=\small]
 \draw[->] (-0.08,0)--(1.76,0) node[right]{\(t\)};
 \draw[->] (0,-0.42)--(0,0.40) node[above]{\(J(t)\)};
 \draw[AnalysisBlue,thick,domain=0:1.61,samples=90]
   plot (\x,{0.5*\x*\x-0.25*\x*\x*\x*\x});
 \draw[densely dashed,gray] (0,0.25)--(1,0.25)--(1,0);
 \fill[AnalysisBlue] (1,0.25) circle (1.5pt);
 \node[left] at (0,0.25) {\(\frac14\)};
 \node[below] at (1,0) {\(1\)};
 \node[below left] at (0,0) {\(0\)};
 \node[above right] at (1,0.25) {必经的高度};
 \node[above] at (1.4142,0) {\(\sqrt2\)};
\end{tikzpicture}
\caption[一维能量的山口图像]{一维能量 \(J(t)=t^2/2-t^4/4\) 的一段曲线。从 \(0\) 到大于 \(\sqrt2\) 的正数，任何连续路径都经过 \(t=1\)，因而不能始终低于能量 \(1/4\)。}
\label{b2:fig:14-mountain-profile}
\end{figure}

这个想法有两个困难。首先，越来越接近最优高度的路径未必收敛；其次，路径上的高点也未必有收敛子列。Palais–Smale 条件解决第二个困难，下降流则把“不是临界点”转化为“还能继续降低路径”。以下采用实 Hilbert 空间中梯度局部 Lipschitz 的版本，并在一个次临界幂非线性模型上逐项验证这些条件。

\subsection{Palais–Smale 条件与紧性}
\label{b2:subsec:140501}

设 \(H\) 为实 Hilbert 空间，内积记为 \((\cdot,\cdot)_H\)。对 \(J\in C^1(H;\mathbb R)\)，第一卷定理8.4.7的 Riesz 表示把连续线性泛函 \(DJ(u)\) 写成
\[
 DJ(u)h=(\nabla J(u),h)_H,
 \quad \|DJ(u)\|_{H^*}=\|\nabla J(u)\|_H.
\]
这里的梯度由所选 Hilbert 内积确定。例如，能量空间 \(H_0^1\) 的梯度通常不是逐点计算的偏微分表达式 \(-\Delta u-g(u)\)；后者首先是对偶空间中的对象，要经过 Riesz 对应才成为 \(H_0^1\) 中的向量。\(DJ(u)=0\) 的点仍称为驻点，也称临界点；其函数值称为临界值。

\begin{definition}[Palais–Smale 条件]\label{b2:def:14-palais-smale}
设 \(H\) 为实 Hilbert 空间，\(J\in C^1(H;\mathbb R)\)，\(c\in\mathbb R\)。若每个满足
\[
 J(u_n)\longrightarrow c,
 \quad \|DJ(u_n)\|_{H^*}\longrightarrow0
\]
的序列都具有在 \(H\) 范数中收敛的子列，就称 \(J\) 在高度 \(c\) 满足\term[palais-smale-tiaojian]{Palais–Smale 条件}[Palais--Smale condition]，记为 \((\mathrm{PS})_c\)。若只要求 \((J(u_n))\) 有界、\(\|DJ(u_n)\|_{H^*}\to0\) 就能得到范数收敛子列，则称 \(J\) 满足 Palais–Smale 条件，简称 \((\mathrm{PS})\)。满足前面两个极限的序列称为高度 \(c\) 的 Palais–Smale 序列。
\end{definition}

条件同时记录“高度有限”和“各方向斜率趋零”。它并不要求每个能量有界序列都收敛，也不把弱收敛当成结论。若 \(u_{n_k}\to u\) 于 \(H\)，连续性给出 \(J(u)=c\)、\(DJ(u)=0\)，所以这个紧性条件确实把近似驻点变为驻点。全局 \((\mathrm{PS})\) 等价于每个有限 \(c\) 上的 \((\mathrm{PS})_c\)：一方面直接包含，另一方面先从有界实数列 \(J(u_n)\) 抽出一个收敛子列，再在所得高度应用条件。

\ProseParagraphBreak

无限维中，“序列有界”还不足以验证 \((\mathrm{PS})\)。在实空间 \(\ell^2\) 上令
\[
 J(x)=\frac12\sum_{j=1}^{\infty}\frac{x_j^2}{j}.
\]
标准单位向量 \(e_n\) 满足 \(J(e_n)=1/(2n)\to0\)、\(\|\nabla J(e_n)\|=1/n\to0\)，但不同两项距离为 \(\sqrt2\)。因此 \((\mathrm{PS})_0\) 失败，尽管这组近似驻点始终位于单位球上。后面的椭圆模型将先取得有界性，再借紧嵌入和方程残差得到能量范数中的强收敛。

\begin{lemma}[非临界高度附近的梯度下界]\label{b2:lem:14-ps-gradient-gap}
设 \(H\) 为实 Hilbert 空间，\(J\in C^1(H;\mathbb R)\)，\(c\in\mathbb R\)。假定 \(J\) 满足 \((\mathrm{PS})_c\)，且不存在同时满足 \(J(u)=c\)、\(DJ(u)=0\) 的 \(u\in H\)。则存在 \(a>0\)、\(\delta>0\)，使
\[
 |J(u)-c|\le a\quad\Longrightarrow\quad
 \|\nabla J(u)\|_H\ge\delta.
\]
\end{lemma}
\begin{proof}
若不存在这样的 \(a,\delta\)，则对每个正整数 \(n\) 可以找到 \(u_n\)，使 \(|J(u_n)-c|\le1/n\)、\(\|DJ(u_n)\|_{H^*}<1/n\)。由 \((\mathrm{PS})_c\) 抽出范数收敛子列，极限 \(u\) 满足 \(J(u)=c\)、\(DJ(u)=0\)，与假设矛盾。
\end{proof}

在有限维紧集上，连续且处处为正的梯度长度有正下界；这里的能量带通常不紧。Palais–Smale 条件提供了另一种办法：假如统一下界失效，就会产生一组近似驻点，而它们的极限正是被排除的临界点。

\subsection{变形引理}
\label{b2:subsec:140502}

当 \(\nabla J(u)\ne0\) 时，方向 \(-\nabla J(u)\) 降低能量。为了把所有点在同一段时间内降低相近的高度，我们改用
\[
 -\frac{\nabla J(u)}{\|\nabla J(u)\|_H^2}.
\]
沿这个向量场，能量的时间导数恰为 \(-1\)。分母也说明了为什么需要上一引理：若梯度过小，下降速度在空间中可能失去控制。我们还要在目标高度之外关掉运动，使路径的低能量端点保持原位。

\begin{lemma}[有界局部 Lipschitz 向量场的流]\label{b2:lem:14-hilbert-flow}
设 \(H\) 为实 Hilbert 空间，\(F:H\to H\) 局部 Lipschitz，且存在 \(M<\infty\) 使 \(\|F(u)\|_H\le M\) 对全部 \(u\in H\) 成立。对每个 \(u\in H\)，初值问题
\[
 \dot z(t)=F(z(t)),\quad z(0)=u
\]
在全部 \(t\ge0\) 上有唯一 \(C^1\) 解。记解为 \(\eta(t,u)\)，则 \(\eta:[0,\infty)\times H\to H\) 连续，且 \(\|\eta(t,u)-u\|_H\le Mt\)。若 \(F(u)=0\)，则 \(\eta(t,u)=u\)。
\end{lemma}
\begin{proof}
在初值附近选取半径 \(r>0\) 的闭球，使 \(F\) 在稍大的邻域内有 Lipschitz 常数 \(L\)。对足够小的 \(T>0\)，令 \(MT<r\)、\(LT<1\)。完备空间 \(C([0,T];H)\) 中取闭集
\(\{z:\sup_t\|z(t)-u\|_H\le r\}\)，映射
\[
 (\mathcal Tz)(t)=u+\int_0^tF(z(s))\,\mathrm ds
\]
保持该集，并且 Lipschitz 常数不超过 \(q=LT<1\)。这里连续的 \(H\) 值函数的积分可由 Riemann 和在完备空间中取极限定义。从常值曲线 \(z_0(t)=u\) 开始迭代 \(z_{k+1}=\mathcal Tz_k\)，则
\[
 \|z_{k+1}-z_k\|_\infty\le q^k\|z_1-z_0\|_\infty.
\]
几何级数估计使 \((z_k)\) 成为 Cauchy 列；其极限 \(z\) 由 \(\mathcal T\) 连续性满足 \(z=\mathcal Tz\)。两个不动点的距离不超过自身的 \(q\) 倍，故只能相同。积分基本定理再给出 \(C^1\) 解。两个局部解在共同存在区间相同，因而可以逐段拼接。

若最大正向存在时间为有限的 \(T_*\)，由 \(\|\dot z\|_H\le M\) 有
\(\|z(t)-z(s)\|_H\le M|t-s|\)。于是当 \(t\uparrow T_*\) 时，\(z(t)\) 在 \(H\) 中有极限。从该极限重新解局部初值问题，就把原解延长到 \(T_*\) 以后，矛盾。因此解全局存在，积分同一估计即得位移界。

连续依赖也来自局部 Lipschitz 估计。在一个共同的 Lipschitz 邻域内，两条轨道满足
\[
 \|z(t)-\widetilde z(t)\|_H
 \le \|z(0)-\widetilde z(0)\|_H
      +L\int_0^t\|z(s)-\widetilde z(s)\|_H\,\mathrm ds,
\]
故由\cref{b2:lem:ode-gronwall}不超过 \(e^{Lt}\|z(0)-\widetilde z(0)\|_H\)。该引理用于此处连续非负的距离函数，其积分常数正是 \(L\)。固定轨道在有限时间区间上的像是紧集，可用有限个这样的邻域覆盖，再把时间分成有限段依次使用该估计。初值足够接近时，各段轨道仍停留在所选邻域内，于是得到初值的连续依赖；与 \(\|\dot z\|_H\le M\) 合并，得到时间与初值的联合连续性。最后，若 \(F(u)=0\)，常值轨道已经满足初值问题，唯一性给出结论。
\end{proof}

\begin{lemma}[变形引理]\label{b2:lem:14-deformation}
设 \(H\) 为实 Hilbert 空间，\(J\in C^1(H;\mathbb R)\)，且 \(\nabla J:H\to H\) 局部 Lipschitz。设 \(c\in\mathbb R\)、\(\varepsilon>0\)、\(\delta>0\)，并且
\[
 |J(u)-c|\le3\varepsilon
 \quad\Longrightarrow\quad\|\nabla J(u)\|_H\ge\delta.
\]
则存在连续映射 \(\Phi:H\to H\)，满足
\begin{enumerate}
\item\label{b2:item:14-deform-fixed}
若 \(|J(u)-c|\ge2\varepsilon\)，则 \(\Phi(u)=u\)；
\item\label{b2:item:14-deform-decrease}
对全部 \(u\in H\)，有 \(J(\Phi(u))\le J(u)\)；
\item\label{b2:item:14-deform-band}
若 \(J(u)\le c+\varepsilon\)，则 \(J(\Phi(u))\le c-\varepsilon\)。
\end{enumerate}
\end{lemma}
\begin{proof}
取 Lipschitz 函数 \(\chi:\mathbb R\to[0,1]\)，在 \([c-\varepsilon,c+\varepsilon]\) 上为一，在 \(( -\infty,c-2\varepsilon]\cup[c+2\varepsilon,\infty)\) 上为零，中间两段线性连接。在开带 \(|J(u)-c|<3\varepsilon\) 上令
\[
 F(u)=-\chi(J(u))\frac{\nabla J(u)}{\|\nabla J(u)\|_H^2},
\]
在其余点令 \(F(u)=0\)。开带内分母有正下界，所以 \(F\) 局部 Lipschitz。开带边缘附近的 \(\chi(J(u))\) 恒为零，两种定义相容，故 \(F\) 在整个 \(H\) 上局部 Lipschitz，并有 \(\|F\|_H\le1/\delta\)。

由\cref{b2:lem:14-hilbert-flow}得到流 \(\eta(t,u)\)。链式法则给出
\[
 \frac{\mathrm d}{\mathrm dt}J(\eta(t,u))
 =-\chi(J(\eta(t,u)))\le0.
\]
定义 \(\Phi(u)=\eta(2\varepsilon,u)\)。连续性以及能量不增已经成立；在 \(|J(u)-c|\ge2\varepsilon\) 的点，\(F(u)=0\)，所以这些点固定。

最后设 \(J(u)\le c+\varepsilon\)。若起点已经不高于 \(c-\varepsilon\)，能量不增即可。否则，只要轨道能量还高于 \(c-\varepsilon\)，它就处于 \((c-\varepsilon,c+\varepsilon]\)，在那里导数恒为 \(-1\)。它因此在不超过 \(2\varepsilon\) 的时间内降到 \(c-\varepsilon\)，以后仍不增加。这证明最后一项。
\end{proof}

上述结论称为\term[bianxing-yinli]{变形引理}[deformation lemma]。它改变的是整条路径，而不是只挑出路径的一个最大点；这使变形后的对象仍是一条连续路径。局部 Lipschitz 梯度保证了流的唯一性和对初值的连续依赖。一般 \(C^1\) 泛函的梯度只保证连续，不能直接套用这个流引理；更一般的山路理论用局部 Lipschitz 的伪梯度向量场替代原梯度。Antonio Ambrosetti 的《Critical points and nonlinear variational problems》\cite[Sections 2--3]{Ambrosetti1992Critical}讨论了相应的紧性、变形和极小极大方法。

\subsection{山路定理}
\label{b2:subsec:140503}

设零点周围有一圈能量至少为 \(\alpha>0\) 的球面，而球面外的某个点 \(e\) 能量不高于零。每条连接 \(0\) 与 \(e\) 的路径都必须穿过这圈球面；我们能够改变穿越位置，却不能让整条路径都处在 \(\alpha\) 以下。\term[shanlu-dingli]{山路定理}[mountain pass theorem]把这一几何障碍转化为临界点。

\begin{theorem}[山路定理，局部 Lipschitz 梯度版本]\label{b2:thm:14-mountain-pass}
设 \(H\) 为实 Hilbert 空间，\(J\in C^1(H;\mathbb R)\)，\(\nabla J:H\to H\) 局部 Lipschitz，且 \(J(0)=0\)。假定存在 \(\rho>0\)、\(\alpha>0\)、\(e\in H\)，满足
\[
 J(u)\ge\alpha\quad(\|u\|_H=\rho),
 \quad \|e\|_H>\rho,
 \quad J(e)\le0.
\]
令
\[
 \Gamma=\{\gamma\in C([0,1];H):\gamma(0)=0,\ \gamma(1)=e\},
 \quad c=\inf_{\gamma\in\Gamma}\max_{0\le t\le1}J(\gamma(t)).
\]
若 \(J\) 在这个 \(c\) 上满足 \((\mathrm{PS})_c\)，则存在 \(u_*\in H\) 使
\[
 DJ(u_*)=0,\quad J(u_*)=c\ge\alpha>0.
\]
特别地，\(u_*\ne0\)，且 \(u_*\ne e\)。
\end{theorem}
\begin{proof}
线段 \(\gamma(t)=te\) 属于 \(\Gamma\)，而 \(J\circ\gamma\) 在紧区间上有有限最大值，因此 \(c<\infty\)。任意 \(\gamma\in\Gamma\) 的范数从 \(0\) 连续变到大于 \(\rho\)，介值定理给出某个 \(t\) 使 \(\|\gamma(t)\|_H=\rho\)。故每条路径的最大能量至少为 \(\alpha\)，从而 \(c\ge\alpha>0\)。

反设高度 \(c\) 上没有临界点。\cref{b2:lem:14-ps-gradient-gap}给出能量带宽 \(a>0\) 与梯度下界 \(\delta>0\)。取 \(\varepsilon>0\) 使 \(3\varepsilon<a\)、\(2\varepsilon<c\)，应用\cref{b2:lem:14-deformation}得到 \(\Phi\)。两个端点的能量都不高于零，因而小于 \(c-2\varepsilon\)，所以 \(\Phi(0)=0\)、\(\Phi(e)=e\)。

按下确界的定义，可取 \(\gamma\in\Gamma\) 满足 \(\max_tJ(\gamma(t))<c+\varepsilon\)。复合路径 \(\widetilde\gamma=\Phi\circ\gamma\) 连续且端点不变，因此仍属于 \(\Gamma\)。但变形引理给出
\[
 \max_{0\le t\le1}J(\widetilde\gamma(t))\le c-\varepsilon<c,
\]
与 \(c\) 是全部此类路径最大值的下确界矛盾。于是高度 \(c\) 上有临界点。其能量严格高于两个端点，故与它们不同。
\end{proof}

证明没有要求存在一条“最佳路径”。我们只取最大能量小于 \(c+\varepsilon\) 的近似路径；若高度 \(c\) 处没有驻点，就把整条近似路径向下推，造成矛盾。在开头的一维例子中，取 \(e=2\)，每条路径都经过 \(1\)，线段路径又实现最大值 \(1/4\)，故 \(c=1/4\)。无限维的结论保留这种过山口的含义，但通常不能据此断言所得临界点唯一，也不能只凭名称判定它的二阶导数有几个负方向。

\subsection{次临界半线性椭圆模型}
\label{b2:subsec:140504}

现在考虑
\[
 -\Delta u=|u|^{p-2}u\quad\text{于 }\Omega,
 \quad u\in H_0^1(\Omega;\mathbb R).
\]
这里的右端与线性 Poisson 方程相反，来自能量中的负项：
\[
 J(u)=\frac12\int_\Omega|\nabla u|^2\,\mathrm dx
          -\frac1p\int_\Omega|u|^p\,\mathrm dx.
\]
当 \(p>2\) 时，小振幅的二次梯度能量占优势，大振幅的负 \(p\) 次项占优势。这给出山路几何；严格低于 Sobolev 临界指数则提供紧性。

\begin{theorem}[次临界幂方程的非零弱解]\label{b2:thm:14-subcritical-mountain-model}
设整数 \(d\ge1\)，\(\Omega\subset\mathbb R^d\) 为非空有界开集，并令
\[
 2^*=\begin{cases}2d/(d-2),&d\ge3,\\+\infty,&d=1,2.\end{cases}
\]
若 \(2<p<2^*\)，则存在非零的实函数 \(u\in H_0^1(\Omega)\)，使
\[
 \int_\Omega\nabla u\cdot\nabla v\,\mathrm dx
   =\int_\Omega|u|^{p-2}uv\,\mathrm dx
 \quad(v\in H_0^1(\Omega;\mathbb R)).
\]
零边界条件按 \(H_0^1\) 的闭包定义理解。该解可取为能量
\(J(u)=\tfrac12\|\nabla u\|_2^2-\tfrac1p\|u\|_p^p\)
的一个正能量山路临界点。
\end{theorem}
\begin{proof}
在 \(H=H_0^1(\Omega;\mathbb R)\) 上取
\((u,v)_H=\int_\Omega\nabla u\cdot\nabla v\)。第一卷定理24.1.2的零边界 Poincaré 不等式说明梯度范数与完整 \(H^1\) 范数等价，所以这是 Hilbert 空间。第一卷推论25.2.6还说明：在任意有界开集上，零边界嵌入 \(H\to L^p(\Omega)\) 在当前严格次临界指数下是紧的。特别有
\[
 \|u\|_p\le C\|u\|_H.
\]
这个零边界版本适用于这里的有界开集，无需定义逐点边界值。

先验证泛函及其导数。记 \(N(s)=|s|^{p-2}s\)。由一维微积分，
\[
 |N(a)-N(b)|\le(p-1)(|a|+|b|)^{p-2}|a-b|
 \quad(a,b\in\mathbb R).
\]
对函数使用 Hölder 不等式，得到
\begin{equation}\label{b2:eq:14-power-lipschitz}
 \|N(u)-N(w)\|_{p'}
 \le(p-1)(\|u\|_p+\|w\|_p)^{p-2}\|u-w\|_p,
 \quad p'=\frac p{p-1}.
\end{equation}
例如其中的三个指数满足 \((p-2)/p+1/p=1/p'\)。沿 \(u+th\) 积分一阶导数，\eqref{b2:eq:14-power-lipschitz}使余项在有界球内不超过常数乘 \(\|h\|_p^2\)。因此 \(J\) Fréchet 可微，且
\[
 DJ(u)v=(u,v)_H-\int_\Omega N(u)v,
 \quad
 \|DJ(u)-DJ(w)\|_{H^*}\le C_R\|u-w\|_H
\]
在 \(\|u\|_H,\|w\|_H\le R\) 时成立。Riesz 对应等距，故梯度在每个有界球上 Lipschitz。

其次验证山路几何。嵌入估计给出
\[
 J(u)\ge\frac12\|u\|_H^2-\frac{C^p}{p}\|u\|_H^p.
\]
因 \(p>2\)，可选小的 \(\rho>0\)，使右端在 \(\|u\|_H=\rho\) 时至少为 \(\alpha=\rho^2/4>0\)。再选非零 \(\varphi\in C_c^\infty(\Omega;\mathbb R)\)。沿射线有
\[
 J(t\varphi)=\frac{t^2}2\|\varphi\|_H^2
                  -\frac{t^p}p\|\varphi\|_p^p\longrightarrow-\infty
 \quad(t\to+\infty).
\]
故可取足够大的 \(T\)，使 \(e=T\varphi\) 满足 \(\|e\|_H>\rho\)、\(J(e)<0\)。

最后证明 \((\mathrm{PS})\)。设 \((J(u_n))\) 有界、\(\|DJ(u_n)\|_{H^*}\to0\)。消去非线性项得到关键恒等式
\[
 J(u_n)-\frac1pDJ(u_n)u_n
       =\left(\frac12-\frac1p\right)\|u_n\|_H^2.
\]
因此存在 \(C<\infty\) 与 \(\varepsilon_n\to0\)，使
\[
 \left(\frac12-\frac1p\right)\|u_n\|_H^2
       \le C+\varepsilon_n\|u_n\|_H.
\]
左边的系数为正，二次式估计说明 \((u_n)\) 有界。由第一卷定理9.4.7的自反空间弱子列结论，抽取 \(u_n\rightharpoonup u\) 于 \(H\)；再由上述紧嵌入，沿进一步子列有 \(u_n\to u\) 于 \(L^p\)。强极限与弱极限相同，因为连续嵌入保留弱收敛。

现在以 \(u_n-u\) 代入近似方程。其 \(H\) 范数有界，所以
\(DJ(u_n)(u_n-u)\to0\)。另外，Hölder 不等式给出
\[
 \left|\int_\Omega N(u_n)(u_n-u)\right|
 \le\|u_n\|_p^{p-1}\|u_n-u\|_p\longrightarrow0.
\]
于是 \((u_n,u_n-u)_H\to0\)。弱收敛又给出 \((u,u_n-u)_H\to0\)，两式相减便是
\(\|u_n-u\|_H^2\to0\)。这证明了所需的范数紧性。

应用\cref{b2:thm:14-mountain-pass}得到 \(DJ(u)=0\)、\(J(u)>0\)。前面的导数公式给出全部 \(H_0^1\) 测试的弱等式，正能量排除了 \(u=0\)。
\end{proof}

次临界性在证明中有一个精确的使用位置：把有界弱收敛序列变成 \(L^p\) 强收敛序列，从而控制非线性项。当 \(d\ge3\)、\(p=2^*\) 时，连续嵌入仍然成立，紧嵌入却不成立。取内部点 \(x_0\) 和固定非零 \(\varphi\in C_c^\infty(\mathbb R^d)\)，对足够小的 \(\varepsilon\) 令
\[
 u_\varepsilon(x)=\varepsilon^{-(d-2)/2}
               \varphi\left(\frac{x-x_0}{\varepsilon}\right).
\]
其支集落在 \(\Omega\)，而 \(\|\nabla u_\varepsilon\|_2\)、\(\|u_\varepsilon\|_{2^*}\) 都保持不变，函数却在 \(x\ne x_0\) 处最终为零。这组集中函数没有 \(L^{2^*}\) 强收敛子列。它解释了紧性步骤的障碍；要判断临界方程的存在性，还须结合能量高度、区域几何或另外的紧性机制。

\ProseParagraphBreak

所得结论是自然能量空间中的非零弱解。它没有自动给出符号、唯一性或经典二阶正则性；这些问题分别需要正部截断、比较原理或椭圆估计等进一步分析。对 \(p>2\) 而言，直接极小化当前能量会沿射线降到负无穷，山路方法却从零点附近的正能垒中找到了另一种驻点。这是极小极大方法与极小化方法在这个模型中的实际分工。

% 本章小结（不计入编号）
\begin{chaptersummary}
第一变分把极小问题转成对全部允许方向成立的积分等式；分部积分把该等式拆成内部 Euler--Lagrange 方程和自然边界条件。允许的方向由约束决定，故固定边界、自由边界及单侧约束产生不同形式。第二变分非负只是必要条件，一致正下界才给出局部极小性的稳定判据。对凸二次能量，精确能量差进一步把驻点与全局唯一极小点联系起来。

\ProseParagraphBreak

直接方法从极小化序列出发。强制性控制序列，Sobolev 空间的自反性提供弱极限，闭凸约束保留边界数据，弱下半连续性完成能量比较。所获极小元首先属于弱函数空间；它是否更光滑，需要另外的方程估计。

\ProseParagraphBreak

山路方法比较的是路径上的最高能量。若这一高度附近没有临界点，沿负梯度方向作适当截断的下降流，就能把整条路径压低，违背高度的极小性。Palais--Smale 条件把能量受控且导数趋零的序列转成真正的临界点；次临界 Sobolev 紧嵌入则在半线性模型中提供所需的强收敛。下一章将以弱形式直接求解一般双线性问题，尤其处理不能来自对称能量的算子。
\end{chaptersummary}


\BookExercises

以下习题均为本书自拟。题目通过改变空间、边界或能量符号，检验变分结论的适用范围；最后两题计算山路高度及其随方程参数的变化。一维 Sobolev 函数的端点值按第一卷命题21.4.1取闭区间上的绝对连续代表，因此 \(u(1)-u(0)=\int_0^1u'\)。

% 自拟：无限维正定第二变分不足以保证极小。
\Needspace{6\baselineskip}
\begin{exercise}\label{b2:ex:14-positive-hessian}
在实空间 \(\ell^2\) 上令 \(J(x)=\sum_{n\ge1}x_n^2/n-\lVert x\rVert^4\)。证明 \(J\) 二阶连续 Fréchet 可微，零点为驻点且 \(D^2J(0)[h,h]>0\) 对每个 \(h\ne0\) 成立，但零点不是局部极小点。指出第二变分判别的哪项条件不成立。
\end{exercise}
\begin{solution}
有界自伴算子 \(Ax=(x_n/n)\) 给出 \(J(x)=\langle Ax,x\rangle-\langle x,x\rangle^2\)，所以它是连续多项式，特别是二阶连续可微。\(DJ(0)=0\)，且 \(D^2J(0)[h,h]=2\sum h_n^2/n>0\)。令 \(x^{(n)}=\sqrt{2/n}\,e_n\)，则 \(\lVert x^{(n)}\rVert\to0\)，而 \(J(x^{(n)})=2/n^2-4/n^2=-2/n^2<0=J(0)\)。在单位向量 \(e_n\) 上第二变分等于 \(2/n\to0\)，故不存在统一正下界 \(c\lVert h\rVert^2\)。
\end{solution}

% 自拟：带权端点问题，应用驻点和严格凸性。
\Needspace{6\baselineskip}
\begin{exercise}\label{b2:ex:14-weighted-string}
设 \(a_0>0\)、\(a\in C^1([0,1];\mathbb R)\)、\(a(x)\geq a_0\)。求 \(\tfrac12\int_0^1a(x)|u'(x)|^2\,\mathrm dx\) 在实函数 \(u\in H^1(0,1)\)、\(u(0)=0,u(1)=1\) 中的极小元和极小值，并用能量差证明唯一性。
\end{exercise}
\begin{solution}
记 \(L=\int_0^1a(s)^{-1}\,\mathrm ds\)，令 \(u(x)=L^{-1}\int_0^x a(s)^{-1}\,\mathrm ds\)，则 \(au'=L^{-1}\)。任一零端点扰动 \(h\in H_0^1(0,1)\) 满足 \(\int au'h'=L^{-1}(h(1)-h(0))=0\)，此等式也可由光滑逼近得到。因此能量差为 \(\tfrac12\int a|h'|^2\ge0\)，且等号迫使 \(h'=0\)、\(h=0\)。极小值是 \((2L^2)^{-1}\int a^{-1}=1/(2L)\)。
\end{solution}

% 自拟：边界弹簧与外法向符号。
\Needspace{6\baselineskip}
\begin{exercise}\label{b2:ex:14-boundary-spring}
给定 \(\beta>0\)、\(q\in\mathbb R\)，在实函数 \(u\in H^1(0,1)\)、\(u(0)=0\) 中极小化
\[J(u)=\tfrac12\int_0^1|u'|^2\,\mathrm dx+\tfrac\beta2u(1)^2-qu(1).\]
求极小元，并核验自然边界条件与全局极小性。
\end{exercise}
\begin{solution}
第一变分为 \(\int u'h'+(\beta u(1)-q)h(1)\)，其中 \(h(0)=0\)。内部方程是 \(u''=0\)，右端条件是 \(u'(1)+\beta u(1)=q\)，所以 \(u(x)=qx/(1+\beta)\)。对任意允许扰动，交叉项消失，能量差为 \(\tfrac12\int|h'|^2+\tfrac\beta2h(1)^2\)。它非负，等号时 \(h=0\)，故所得解是唯一全局极小元。
\end{solution}

% 自拟：积分约束产生常数乘子，含显式解。
\Needspace{6\baselineskip}
\begin{exercise}\label{b2:ex:14-mass-constraint}
给定 \(M\in\mathbb R\)，在实空间 \(H_0^1(0,1;\mathbb R)\) 中，附加 \(\int_0^1u=M\)，极小化 \(\tfrac12\int_0^1|u'|^2\)。不用抽象乘子定理，求出唯一极小元和极小值。
\end{exercise}
\begin{solution}
取 \(u(x)=6Mx(1-x)\)，则端点为零且积分为 \(M\)。允许扰动 \(h\) 满足零端点和零积分；分部积分给出 \(\int u'h'=-\int u''h=12M\int h=0\)。能量差为 \(\tfrac12\lVert h'\rVert_2^2\)，故它是唯一极小元。又 \(\int_0^1(1-2x)^2\,\mathrm dx=1/3\)，极小值为 \(6M^2\)。常数 \(-u''=12M\) 正是积分约束所允许留下的乘子；若错误地令所有紧支集方向都可行，就会误得 \(u''=0\)。
\end{solution}

% 自拟：临界负势能与最优Poincare常数。
\Needspace{6\baselineskip}
\begin{exercise}\label{b2:ex:14-negative-potential}
设 \(\lambda\in\mathbb R\)，在实空间 \(H_0^1(0,\pi;\mathbb R)\) 上考察 \(J_\lambda(u)=\tfrac12\int_0^\pi(|u'|^2-\lambda|u|^2)\)。使用区间正弦展开所得 \(\lVert u\rVert_2\le\lVert u'\rVert_2\)，分别讨论 \(\lambda<1\)、\(\lambda=1\)、\(\lambda>1\) 的极小元与强制性。
\end{exercise}
\begin{solution}
若 \(\lambda<1\)，当 \(0\le\lambda<1\) 时 \(J_\lambda\ge(1-\lambda)\lVert u'\rVert_2^2/2\)，当 \(\lambda<0\) 时更有 \(J_\lambda\ge\lVert u'\rVert_2^2/2\)，故强制且唯一极小元为零。若 \(\lambda=1\)，正弦展开给出 \(2J_1=\frac\pi2\sum_{n\ge1}(n^2-1)|a_n|^2\)，所以极小元恰为 \(c\sin x\)，泛函不强制。若 \(\lambda>1\)，沿 \(u=c\sin x\) 有 \(J_\lambda=\pi(1-\lambda)c^2/4\to-\infty\)，故没有有限极小值。展开公式先对有限正弦和成立，再由 \(H_0^1\) 逼近推广。
\end{solution}

% 自拟：非对称双线性型的有限维反例。
\Needspace{6\baselineskip}
\begin{exercise}\label{b2:ex:14-skew-part}
令 \(A=\left(\begin{smallmatrix}2&-3\\3&2\end{smallmatrix}\right)\)、\(f=(1,0)\)。对 \(u\in\mathbb R^2\)，分别求 \(Au=f\) 的解及 \(\tfrac12u^{\mathsf T}Au-f\cdot u\) 的极小元。证明 \((Au)\cdot u\ge2|u|^2\)，并解释这一强制性为何不足以使两个问题相同。
\end{exercise}
\begin{solution}
\(A^{-1}=13^{-1}\left(\begin{smallmatrix}2&3\\-3&2\end{smallmatrix}\right)\)，故方程解为 \((2/13,-3/13)\)。由于 \(u^{\mathsf T}Au=2|u|^2\)，泛函是 \(|u|^2-u_1\)，唯一极小元为 \((1/2,0)\)。强制性等式确实成立，但二次型只保留 \((A+A^{\mathsf T})/2=2I\)；驻点方程为 \(2u=f\)，反对称部分没有进入能量。
\end{solution}

% 自拟：障碍约束产生变分不等式。
\Needspace{6\baselineskip}
\begin{exercise}\label{b2:ex:14-obstacle}
设整数 \(d\geq1\)，\(\Omega\subset\mathbb R^d\) 为非空有界开集，\(K\) 是实空间 \(H_0^1(\Omega;\mathbb R)\) 的非空范数闭凸集，\(f\in L^2(\Omega;\mathbb R)\)。令 \(E(v)=\frac12\int_\Omega|\nabla v|^2-\int_\Omega fv\)。证明 \(E\) 在 \(K\) 上有唯一极小元 \(u\)，并证明它等价于
\[\int_\Omega\nabla u\cdot\nabla(v-u)\ge\int_\Omega f(v-u)\qquad(v\in K).\]
\end{exercise}
\begin{solution}
Poincaré 不等式给出强制性；闭凸集弱闭，能量凸且范数连续，故弱下半连续。第一卷定理 9.5.9 给出存在性，梯度平方的严格凸性给出唯一性。若 \(u\) 极小，对 \(u+t(v-u)\in K\)、\(0<t\le1\) 展开能量，除以 \(t\) 令 \(t\downarrow0\)，即得不等式。反之，取 \(h=v-u\) 于精确二次展开，线性项由不等式非负，平方项也非负，故 \(E(v)\ge E(u)\)。凸约束只提供单侧可行线段，因而得到不等式。
\end{solution}

% 自拟：受迫阻尼波的能量估计。
\Needspace{6\baselineskip}
\begin{exercise}\label{b2:ex:14-damped-wave}
设整数 \(d\geq1\)，\(c,T,\gamma>0\)，\(\Omega\subset\mathbb R^d\) 为非空有界 \(C^1\) 开集，实函数 \(u\) 在 \([0,T]\times\overline\Omega\) 的某开邻域内属于 \(C^2\)。设 \(f\in C([0,T]\times\overline\Omega;\mathbb R)\)，且 \(u_{tt}-c^2\Delta u+\gamma u_t=f\) 于 \((0,T)\times\Omega\)，\(u(t,y)=0\) 对 \(0\leq t\leq T\)、\(y\in\partial\Omega\) 成立。记
\[
 \mathcal E(t)=\frac12\int_\Omega
            (|u_t(t,x)|^2+c^2|\nabla u(t,x)|^2)\,\mathrm{d}x.
\]
证明对 \(0\le t\le T\)，
\[\mathcal E(t)+\frac\gamma2\int_0^t\lVert u_s\rVert_2^2\,\mathrm ds
 \le\mathcal E(0)+\frac1{2\gamma}\int_0^t\lVert f(s)\rVert_2^2\,\mathrm ds.\]
说明该估计怎样给出相同初边值数据的唯一性。
\end{exercise}
\begin{solution}
乘方程以 \(u_t\) 并分部积分，得 \(\mathcal E'+\gamma\lVert u_t\rVert_2^2=\int fu_t\)。由 Young 不等式，右端不超过 \(\lVert f\rVert_2^2/(2\gamma)+\gamma\lVert u_t\rVert_2^2/2\)。移项并对时间积分，得到所需估计。对两个解之差应用，右端外力与初始能量都为零，故差的时间导数和梯度均为零；零边界或零初位移固定剩余常数，差恒为零。
\end{solution}

% 自拟：一维p能量、端点约束与非线性通量；常外力部分增加实质训练，
% 正文存在性、可微性和严格凸性的论证均不依赖本题。
\Needspace{6\baselineskip}
\begin{exercise}\label{b2:ex:14-p-laplace-interval}
设 \(1<p<\infty\)，\(p'=p/(p-1)\)，
\[
 \mathcal C=\{u\in W^{1,p}(0,1;\mathbb R):u(0)=0,\ u(1)=1\},
 \qquad J_p(u)=\frac1p\int_0^1|u'|^p\,\mathrm{d}x.
\]
端点值取绝对连续代表。令 \(A_p(s)=|s|^{p-2}s\)（\(s\ne0\)），并规定 \(A_p(0)=0\)。
\begin{enumerate}
\item\label{b2:item:14-p-interval-unforced}证明 \(J_p\) 在 \(\mathcal C\) 上的唯一极小元为 \(u(x)=x\)，并求极小值。由第一变分写出
\(\int_0^1 A_p(u')h'=0\)（\(h\in W_0^{1,p}(0,1;\mathbb R)\)），
解释这就是分布方程 \((A_p(u'))'=0\)，并指出严格凸性如何保证唯一性。
\item\label{b2:item:14-p-interval-forced}给定 \(\lambda\in\mathbb R\)，在同一类 \(\mathcal C\) 上极小化
\[
 J_{p,\lambda}(u)=J_p(u)-\lambda\int_0^1u(x)\,\mathrm{d}x.
\]
写出以一个实常数 \(C\) 为参数的候选解，并证明端点条件唯一确定 \(C\)，所得候选确为唯一全局极小元。所用方程应写成非线性通量的弱形式，不预设 \(u\) 具有经典二阶导数。
\end{enumerate}
\end{exercise}
\begin{solution}
对 \(u\in\mathcal C\)，绝对连续表示给出 \(\int_0^1u'=1\)。
函数 \(s\mapsto|s|^p\) 严格凸，Jensen 不等式因而给出
\[
 J_p(u)\geq\frac1p\left|\int_0^1u'\right|^p=\frac1p.
\]
等号成立当且仅当 \(u'\) 几乎处处为常数；该常数只能为 \(1\)，再由 \(u(0)=0\)
得到 \(u(x)=x\)。这同时给出极小值与唯一性，而非只验证一个驻点。

允许方向恰为 \(h\in W_0^{1,p}(0,1;\mathbb R)\)。
由\cref{b2:prop:14-power-frechet}，极小点满足
\[
 DJ_p(u)h=\int_0^1A_p(u')h'\,\mathrm{d}x=0.
\]
特别取 \(h\in C_c^\infty(0,1)\)，便是所列分布方程。
反过来，若通量 \(A_p(u')\in L^{p'}(0,1)\) 的分布导数为零，由一维绝对连续代表的积分表示，它几乎处处为常数；
\(A_p\) 严格递增且可逆，故 \(u'\) 也为常数，端点再将它固定为 \(1\)。
这里凸性把驻点变成全局极小点，严格凸性再排除不同的极小点。

对常外力问题，第一变分为
\[
 \int_0^1 A_p(u')h'\,\mathrm{d}x
       =\lambda\int_0^1h\,\mathrm{d}x
 \qquad(h\in W_0^{1,p}(0,1;\mathbb R)),
\]
即 \(-(A_p(u'))'=\lambda\)。
记 \(B_p(q)=\operatorname{sgn}(q)|q|^{1/(p-1)}\)，其中 \(B_p(0)=0\)；
这是 \(A_p\) 的连续严格递增逆映射。积分通量方程得到候选
\[
 u_C(x)=\int_0^x B_p(C-\lambda s)\,\mathrm{d}s.
\]
它属于 \(W^{1,p}(0,1)\)，满足左端值；右端值要求
\[
 F_\lambda(C)\coloneqq\int_0^1 B_p(C-\lambda s)\,\mathrm{d}s=1.
\]
在任何紧的 \(C\) 区间内，被积函数统一有界且连续，故 \(F_\lambda\) 连续。
\(C\) 严格增大时每个被积值都严格增大，所以 \(F_\lambda\) 严格递增。
又 \(C-|\lambda|\leq C-\lambda s\leq C+|\lambda|\)，
由 \(B_p\) 的单调性可知 \(F_\lambda(C)\to\pm\infty\) 当 \(C\to\pm\infty\)。
介值定理于是给出唯一的 \(C\)。

所构造的 \(u=u_C\) 满足 \(A_p(u')=C-\lambda x\)。
对绝对连续且零端点的 \(h\)，分部积分直接给出上述弱等式；
这里微分的是通量，不需要先求 \(u''\)。
最后取任意 \(v\in\mathcal C\)，由凸函数的切线不等式及弱等式，
\[
 \begin{aligned}
 J_{p,\lambda}(v)-J_{p,\lambda}(u)
 &\geq\int_0^1A_p(u')(v'-u')\,\mathrm{d}x
          -\lambda\int_0^1(v-u)\,\mathrm{d}x\\
 &=0.
 \end{aligned}
\]
所以候选为全局极小元。严格凸性使等号必须满足 \(v'=u'\) 几乎处处，
再由同一端点值得 \(v=u\)。当 \(\lambda=0\) 时，参数为 \(C=1\)，公式退回 \(u(x)=x\)。
\end{solution}

% 自拟选读题：二维临界点与路径障碍；增加临界点类型和极小极大值的直接计算。
\Needspace{6\baselineskip}
\begin{optionalexercise}[一个可以直接计算的山口]\label{b2:ex:14-mountain-finite}
在实空间 \(\mathbb R^2\) 上令
\[
 J(x,y)=\frac{x^2+y^2}{2}-\frac{x^4}{4},
 \quad e=(2,0),
 \quad\Gamma=\{\gamma\in C([0,1];\mathbb R^2):\gamma(0)=(0,0),\ \gamma(1)=e\}.
\]
不用抽象山路定理，计算 \(c=\inf_{\gamma\in\Gamma}\max_tJ(\gamma(t))\)。求出全部临界点，利用二阶导数区分它们的类型，并证明 \(J\) 满足 Palais–Smale 条件。
\end{optionalexercise}
\begin{solution}
每条路径的第一坐标都从 \(0\) 走到 \(2\)，因而在某点等于 \(1\)。该点的能量为 \(1/4+y^2/2\ge1/4\)，故 \(c\ge1/4\)。线段路径 \(\gamma(t)=(2t,0)\) 的最大能量等于 \(1/4\)，所以 \(c=1/4\)。

梯度为 \((x-x^3,y)\)，临界点恰为 \((0,0)\)、\((1,0)\)、\((-1,0)\)。Hessian 为 \(\operatorname{diag}(1-3x^2,1)\)。它在原点正定，给出严格局部极小；在另外两点有一个负特征值和一个正特征值，给出鞍点。若 \(\nabla J(x_n,y_n)\to0\)，则 \(y_n\to0\) 且 \(x_n-x_n^3\to0\)。后者迫使 \((x_n)\) 有界，否则三次项的绝对值趋于无穷。因此序列在有限维空间中有收敛子列，特别满足 Palais–Smale 条件。对这个多项式，梯度趋零本身就迫使序列有界。
\end{solution}

% 自拟选读题：参数缩放和路径族双射；不复述正文PS验证。
\Needspace{6\baselineskip}
\begin{optionalexercise}[非线性强度与山路高度]\label{b2:ex:14-mountain-scaling}
设整数 \(d\ge1\)，\(\Omega\subset\mathbb R^d\) 为非空有界开集，\(2<p<2^*\)，其中 \(2^*=2d/(d-2)\) 当 \(d\ge3\)，\(2^*=+\infty\) 当 \(d=1,2\)。在实空间 \(H=H_0^1(\Omega)\) 上令
\[
 J_\lambda(u)=\frac12\|\nabla u\|_2^2-\frac\lambda p\|u\|_p^p
 \quad(\lambda>0).
\]
固定 \(e_1\in H\) 使 \(J_1(e_1)<0\)，令 \(e_\lambda=\lambda^{-1/(p-2)}e_1\)，并以从 \(0\) 到 \(e_\lambda\) 的所有连续路径定义山路高度 \(c_\lambda\)。证明
\[
 c_\lambda=\lambda^{-2/(p-2)}c_1.
\]
若 \(u_1\ne0\) 是 \(-\Delta u=|u|^{p-2}u\) 的零边界弱解，验证 \(u_\lambda=\lambda^{-1/(p-2)}u_1\) 是 \(-\Delta u=\lambda|u|^{p-2}u\) 的非零弱解。解释为什么比较山路高度时，端点也需要随参数缩放。
\end{optionalexercise}
\begin{solution}
设 \(s=\lambda^{-1/(p-2)}>0\)。则 \(\lambda s^p=s^2\)，所以对每个 \(v\in H\) 有 \(J_\lambda(sv)=s^2J_1(v)\)。映射 \(\gamma\mapsto s\gamma\) 将从 \(0\) 到 \(e_1\) 的路径族双射到从 \(0\) 到 \(e_\lambda\) 的路径族，其逆为除以 \(s\)。在每条路径上先取最大值，再对双射的两个路径族取下确界，即得 \(c_\lambda=s^2c_1\)。

对任意 \(v\in H\)，原弱方程给出
\[
 \int\nabla(su_1)\cdot\nabla v
 =s\int|u_1|^{p-2}u_1v
 =\lambda s^{p-1}\int|u_1|^{p-2}u_1v,
\]
最后一步使用 \(\lambda s^{p-2}=1\)，这正是所求弱方程。因 \(s>0\)，非零性保留。若强行保持同一端点，缩放后的路径通常不属于待比较的路径族，上面的双射就不成立；因此能量缩放公式本身不足以推出那个不同路径族的高度公式。
\end{solution}

% !TeX root = ../../Book2.tex
\chapter{弱解与 Lax–Milgram 方法}
\label{b2:ch:15}
\BookChapterNumber{b2:ch:15}
% 本章摘要（不计入编号）
\begin{chapterabstract}
本章把散度型方程写成 Hilbert 空间中的测试等式。先明确 Sobolev 边界条件与负阶源项，再完整证明 Lax--Milgram 定理，并以 Poisson 方程说明存在性、唯一性、稳定性和有限维近似。一般有界可测系数无需具有经典导数；一致椭圆性控制主部，Gårding 不等式处理低阶项。最后借紧嵌入与第一卷的 Fredholm 理论，刻画原方程可能出现的有限维核及伴随相容条件。
\end{chapterabstract}

% 本章目标（不计入编号）
\begin{learninggoals}
\begin{itemize}
\item\label{b2:item:15-goal-weak}辨认测试空间所编码的内部方程与边界条件，正确使用 \(H^{-1}\) 数据。
\item\label{b2:item:15-goal-lax}分别核验型的有界性、强制性与空间完备性，掌握 Lax--Milgram 的闭值域证明。
\item\label{b2:item:15-goal-stability}用能量等式证明唯一性、比较和数据稳定性，区分 Galerkin 最佳与准最佳逼近，并处理非齐次边界延拓。
\item\label{b2:item:15-goal-general}说明有界可测系数已足以给出弱存在性，区别原问题与移位问题。
\item\label{b2:item:15-goal-fredholm}在选读范围内，用伴随核写出 Fredholm 相容条件。
\end{itemize}
\end{learninggoals}

设一根杆的导热系数在接头处突然变化。平衡温度不应跳跃，穿过接头的热通量也应保持一致，但温度斜率可以不同。若只寻找处处二次可微的解，便会把这样的基本模型排除在外。弱形式保留可积梯度和通量平衡，把经典导数无法直接表达的问题放进一个完备的函数空间。

\ProseParagraphBreak

上一章用极小化构造对称能量的解；本章直接求解双线性等式，从而允许非对称系数。Hilbert 空间表示与紧算子的抽象结论采用第一卷第八章和附录 C 的证明，Sobolev 边界与紧嵌入采用第一卷第二十二至二十五章。有关多维 Sobolev 空间与椭圆边值问题的变分表述，可参阅 Haim Brezis 的《Functional Analysis, Sobolev Spaces and Partial Differential Equations》\cite[Chapter 9]{Brezis2011Functional}；Lax--Milgram、能量估计与 Fredholm 择一的进一步阅读见 Lawrence C. Evans 的《Partial Differential Equations》\cite[Section 6.2]{Evans2010Partial}。本章将先完成这些结果在所用空间中的证明，再讨论有限维近似与粗糙系数；进一步阅读不承担正文缺失的证明步骤。

% !TeX root = ../../Book2.tex
\section{弱形式}
\label{b2:sec:1501}

分段材料中的温度可以连续而斜率发生跳跃；此时逐点写二阶导数会掩盖界面上的通量条件。弱形式把微分移到测试函数上，使未知函数只需具有有限能量，同时仍保留方程及界面平衡。本节将第一章的局部弱解概念与上一章的 Sobolev 边界类结合起来，并为比 \(L^2\) 更粗糙的源项确定合适空间。

\subsection{测试函数}
\label{b2:subsec:150101}

在开集 \(\Omega\subset\mathbb R^d\) 中，测试函数取自 \(C_c^\infty(\Omega)\)。它们的支集离边界有正距离，因此分部积分不会引入尚未定义的边界值。若 \(u\in H_{\mathrm{loc}}^1(\Omega)\)、\(A\in L_{\mathrm{loc}}^\infty(\Omega;\mathbb R^{d\times d})\)，则
\[\langle-\operatorname{div}(A\nabla u),\varphi\rangle
 =\int_\Omega A\nabla u\cdot\nabla\varphi\,\mathrm dx\]
定义了一个分布。矩阵约定为 \((A\nabla u)_i=\sum_j a_{ij}\partial_j u\)，因而散度作用于通量的第一个指标。

\ProseParagraphBreak

局部测试函数足以辨认内部方程，但能量估计常须取 \(\varphi=u\)、截断后的 \(u\) 或差商。它们通常不光滑，也不紧支集。合法性依靠的是密度与连续延拓：若左右两边对 \(H_0^1\) 范数连续，而 \(C_c^\infty\) 在该空间稠密，则测试等式自动推广到整个 \(H_0^1\)。若只知道局部 \(H^1\)，便先乘紧支集截断函数，在较小区域完成同样步骤。不能在未建立这些连续性之前直接“取测试函数为解”。

\subsection{弱解}
\label{b2:subsec:150102}

以下主要讨论实函数；实系数方程的复值版本可分别作用于实部与虚部。设整数 \(d\geq1\)，\(\Omega\subset\mathbb R^d\) 为开集，\(A\in L^\infty(\Omega;\mathbb R^{d\times d})\)，\(f\in\mathcal D'(\Omega;\mathbb R)\)。本章主要研究有限全局能量的弱解，因此称 \(u\in H^1(\Omega;\mathbb R)\) 满足散度方程的弱形式，是指
\[\int_\Omega A\nabla u\cdot\nabla\varphi=\langle f,\varphi\rangle
 \qquad(\varphi\in C_c^\infty(\Omega)).\]
第一章的局部弱解只要求 \(u\in H_{\mathrm{loc}}^1(\Omega)\)，使用的是同一个测试恒等式；本处的 \(H^1(\Omega)\) 要求更强，还控制整个区域上的函数与梯度平方积分。全局能量弱解必为局部弱解，反向一般不成立。两者都尚未包含边界条件；零 Dirichlet 条件还要在下一小节另行加入。对于 \(A\in C^1(\Omega;\mathbb R^{d\times d})\)、\(u\in C^2(\Omega)\) 及连续函数 \(f\)，分部积分和测试函数基本引理使该恒等式与经典方程等价；粗糙情形则以积分等式本身为准。

\ProseParagraphBreak

这个空间差别已经出现在三维的去心球上。取
\[
 \Omega=B(0,1)\setminus\{0\}\subset\mathbb R^3,\qquad u(x)=|x|^{-1}.
\]
函数在 \(\Omega\) 内光滑；径向公式给出
\(\Delta u=u''(r)+2u'(r)/r=2r^{-3}-2r^{-3}=0\)。
每个紧含于 \(\Omega\) 的子集都离原点有正距离，所以 \(u\in H_{\mathrm{loc}}^1(\Omega)\)，并对所有紧支集测试函数满足 \(-\Delta u=0\)。但球坐标计算给出
\[
 \begin{aligned}
 \int_\Omega |u|^2\,\mathrm dx&=4\pi\int_0^1\mathrm dr=4\pi,\\
 \int_{\{\varepsilon<|x|<1\}}|\nabla u|^2\,\mathrm dx
 &=4\pi\int_\varepsilon^1r^{-2}\,\mathrm dr\\
 &=4\pi(\varepsilon^{-1}-1)\longrightarrow+\infty.
 \end{aligned}
\]
故它是局部弱解，却不属于本章的全局能量解类。这里原点已从区域删去，测试函数接触不到原点；不能把在整个球内的点源方程与去心球上的齐次方程混同。

\Needspace{6\baselineskip}
\begin{proposition}[一维界面的通量条件]
\label{b2:prop:15-interface-flux}
设 \(-1<s<1\)，实系数 \(a\) 在 \((-1,s)\)、\((s,1)\) 上分别为正常数，实函数 \(u\) 在两闭子区间上分别为 \(C^2\)，且在 \(s\) 处连续。若 \(-(au')'=f\) 在两侧经典成立，其中实值 \(f\in L^2(-1,1)\)，则 \(u\in H^1(-1,1)\) 在整个区间满足 \(-(au')'=f\) 的弱形式，当且仅当
\[a(s-)u'(s-)=a(s+)u'(s+).\]
\end{proposition}
\begin{proof}
分段光滑且连续的 \(u\) 属于 \(H^1\)，其弱导数就是分段的一阶导数。在两段分别积分，外端因测试函数紧支集而消失，内端得到
\[\int_{-1}^1 au'\varphi'=\int_{-1}^1 f\varphi
  +\bigl(a(s-)u'(s-)-a(s+)u'(s+)\bigr)\varphi(s).\]
测试函数在 \(s\) 处可取任意值，因此等式对全部测试函数成立恰当且仅当括号为零。
\end{proof}

例如 \(a=1\) 于左半区间、\(a=2\) 于右半区间，令 \(u(x)=x\)（\(x\le0\)）、\(u(x)=x/2\)（\(x\ge0\)）。其导数跳跃，但通量恒为 \(1\)，所以 \(-(au')'=0\) 弱成立。若误要求 \(u'\) 连续，就会排除这个正确的平衡解；若允许 \(u\) 本身跳跃，又会产生分布导数中的点质量，通常已不属于所选能量空间。

\subsection{Sobolev 边界条件}
\label{b2:subsec:150103}

对任意有界开集，齐次 Dirichlet 条件定义为 \(u\in H_0^1(\Omega)\)；非齐次数据由一个延拓 \(g\in H^1(\Omega)\) 表示，条件为 \(u-g\in H_0^1(\Omega)\)。第一卷定理 23.1.6 说明，在有界 Lipschitz 区域上，这恰好等价于 Sobolev 迹满足 \(\operatorname{Tr}u=\operatorname{Tr}g\)。粗糙区域上仍可使用闭包定义，但不能未经证明把它等同于某个朴素逐点边界条件。

\ProseParagraphBreak

自然边界条件的写法不同。以非空有界 \(C^1\) 开集 \(\Omega\subset\mathbb R^d\) 为例，若实函数 \(u\) 在 \(\overline\Omega\) 的某开邻域内为 \(C^2\)，记 \(f=-\Delta u\)、\(g_N=\partial_\nu u\)，则在 \(H^1(\Omega;\mathbb R)\) 的全部方向测试时，Green 公式是
\[\int_\Omega\nabla u\cdot\nabla v
  =\int_\Omega fv+\int_{\partial\Omega}g_N\operatorname{Tr}v\,\mathrm dS.\]
这条公式先对光滑 \(v\) 成立，再由 \(H^1\) 光滑逼近和迹算子的连续性推广到全部 \(H^1\) 测试函数。它编码 \(\partial_\nu u=g_N\)，而不是要求 \(u\) 在边界为零。取 \(v=1\) 得到必要相容条件 \(\int_\Omega f+\int_{\partial\Omega}g_N=0\)；若区域不连通，还须在每个连通分支上分别成立，见\cref{b2:prop:13-neumann-compatibility}。测试空间不同，编码的边界问题便不同；仅看内部方程无法辨认二者。

\ProseParagraphBreak

即使区域光滑，也不宜把任意边界函数都当作合法 \(H^1\) Dirichlet 数据。数据应当处于迹空间，并存在 \(H^1\) 延拓；第一卷第二十三章在 Lipschitz 图中用分数阶范数刻画这种要求。给定一个 \(H^1\) 延拓 \(g\)，就可以把满足该边界数据的函数写为 \(g+H_0^1(\Omega)\) 中的元素。

\subsection{负阶 Sobolev 空间初步}
\label{b2:subsec:150104}

设整数 \(d\geq1\)，\(\Omega\subset\mathbb R^d\) 是非空有界开集，以梯度内积
\(\langle u,v\rangle_{H_0^1}=\int_\Omega\nabla u\cdot\nabla v\)
赋予实空间 \(H_0^1(\Omega;\mathbb R)\) Hilbert 结构。第一卷定理24.1.2的零边界 Poincaré 不等式给出
\(\|v\|_2\leq C_P\|\nabla v\|_2\)，所以梯度范数与完整 \(H^1\) 范数等价，这个空间仍完备。定义
\[H^{-1}(\Omega)=(H_0^1(\Omega))^*,\qquad
 \lVert f\rVert_{H^{-1}}=\sup_{\lVert\nabla v\rVert_2\le1}|\langle f,v\rangle|.
\]
这是本章使用的 \term[fujie-sobolev-kongjian]{负阶 Sobolev 空间}[negative-order Sobolev space]的第一种情形。换成全 \(H^1\) 范数会得到等价的对偶范数；连续对偶中的元素不变，但下面的精确范数等式使用的是梯度范数约定。每个对偶元素 \(f\in H^{-1}(\Omega)\) 限制到 \(C_c^\infty(\Omega)\) 都给出分布：对支集落在固定紧集 \(K\Subset\Omega\) 内的 \(\varphi\)，有
\[
 |\langle f,\varphi\rangle|
 \leq\|f\|_{H^{-1}}\,|K|^{1/2}\sup_K|\nabla\varphi|,
\]
所以配对在测试函数拓扑下连续。密度保证这个限制是单射，因此也可以把 \(H^{-1}\) 看作一类分布。

实值函数 \(f\in L^2(\Omega)\) 通过 \(\langle f,v\rangle=\int_\Omega fv\) 进入这个对偶，因为
\[
 |\langle f,v\rangle|\leq \|f\|_2\|v\|_2
                 \leq C_P\|f\|_2\|\nabla v\|_2.
\]
更一般地，若 \(F\in L^2(\Omega;\mathbb R^d)\)，则
\(\langle-\operatorname{div}F,v\rangle=\int_\Omega F\cdot\nabla v\)
在 \(H_0^1\) 上连续，范数不超过 \(\|F\|_2\)。后一个例子并不要求散度本身是一项 \(L^2\) 函数；它只要求给出的测试作用有能量范数上界。

\Needspace{6\baselineskip}
\begin{proposition}[负阶数据的通量表示]
\label{b2:prop:15-negative-representation}
设整数 \(d\geq1\)，\(\Omega\subset\mathbb R^d\) 为非空有界开集，\(H_0^1(\Omega;\mathbb R)\) 取梯度内积，\(H^{-1}(\Omega)=(H_0^1(\Omega;\mathbb R))^*\) 取对偶范数
\(\|f\|_{H^{-1}}=\sup_{\|\nabla v\|_2\leq1}|\langle f,v\rangle|\)。
对每个 \(f\in H^{-1}(\Omega)\)，存在唯一 \(w\in H_0^1(\Omega;\mathbb R)\) 使
\[
 \langle f,v\rangle=\int_\Omega\nabla w\cdot\nabla v\,\mathrm dx
 \qquad(v\in H_0^1(\Omega;\mathbb R)).
\]
因而 \(F_0=\nabla w\in L^2(\Omega;\mathbb R^d)\) 是 \(f=-\operatorname{div}F_0\) 的通量表示，且
\(\|F_0\|_2=\|f\|_{H^{-1}}\)。
任一其他通量表示恰为 \(F=F_0+G\)，其中 \(G\in L^2(\Omega;\mathbb R^d)\)、\(\operatorname{div}G=0\) 于 \(\mathcal D'(\Omega)\)，并满足
\[
 \|F\|_2^2=\|F_0\|_2^2+\|G\|_2^2.
\]
特别地，\(F_0\) 是所有通量表示中唯一的最小 \(L^2\) 范数表示。
\end{proposition}
\begin{proof}
Poincaré 不等式使梯度内积给出完备的 Hilbert 空间。直接在这个空间上应用第一卷定理8.4.7的 Riesz 表示，得到唯一 \(w\) 及所列配对，且
\[
 \|f\|_{H^{-1}}=\|w\|_{H_0^1}=\|\nabla w\|_2.
\]
取 \(F_0=\nabla w\) 就得到通量表示。

若 \(F\) 也表示 \(f\)，则 \(G=F-F_0\) 满足
\(\int_\Omega G\cdot\nabla v=0\) 对全部 \(v\in H_0^1\) 成立。特别取紧支集光滑 \(v\)，便知 \(\operatorname{div}G=0\) 于分布意义。反过来，若 \(G\in L^2\) 的分布散度为零，等式先对紧支集测试函数成立；由于其左端对梯度范数连续，密度使它对全部 \(H_0^1\) 成立，所以 \(F_0+G\) 仍表示 \(f\)。最后取 \(v=w\)，有 \(\int_\Omega G\cdot F_0=0\)，展开平方即得范数恒等式；等号 \(\|F\|_2=\|F_0\|_2\) 当且仅当 \(G=0\)，证明最小范数表示的唯一性。
\end{proof}

因此负阶数据确定的是通量的散度；最小范数准则才从所有通量中选出梯度通量。可以自由叠加的 \(G\) 与全部零边界梯度正交，但这个条件只等价于分布散度为零，不能在未定义法向迹时再解释为逐点零法向通量。

这个空间也确实容纳比函数更粗糙的数据，但是否容纳点源取决于维数。以下直接检查点值泛函；所需的不是“点质量是分布”，而是它对能量范数连续。

\Needspace{6\baselineskip}
\begin{proposition}[点质量与能量对偶]\label{b2:prop:15-point-mass-dimension}
设整数 \(d\geq1\)，\(\Omega\subset\mathbb R^d\) 为非空有界开集，\(H^{-1}(\Omega)\) 表示实空间 \(H_0^1(\Omega;\mathbb R)\) 在梯度范数下的连续对偶。对 \(x_0\in\Omega\)，点质量分布定义为 \(\langle\delta_{x_0},\varphi\rangle=\varphi(x_0)\)、\(\varphi\in C_c^\infty(\Omega)\)。若 \(d=1\)、\(\Omega=I=(a,b)\) 为有限开区间，则 \(\delta_{x_0}\in H^{-1}(I)\)；若 \(d\geq2\)，则 \(\delta_{x_0}\notin H^{-1}(\Omega)\)。
\end{proposition}
\begin{proof}
在 \(I=(a,b)\) 上，对光滑零边界函数由
\(\varphi(x_0)=\int_a^{x_0}\varphi'\) 得
\(|\varphi(x_0)|\leq\sqrt{b-a}\|\varphi'\|_2\)。密度把点值泛函延拓到 \(H_0^1(I)\)，其值与绝对连续代表的点值一致：第一卷命题21.4.1给出一维 \(H^1\) 函数的绝对连续代表，而上述增量估计保证零边界光滑逼近的点值趋于该代表的点值。

高维先平移并缩放一个紧含于 \(\Omega\) 的球，使观察点为零、球半径为一；这个固定换元只改变范数常数。若 \(d>2\)，取 \(\eta\in C_c^\infty(B(0,1))\)、\(\eta(0)=1\)，令 \(\eta_\varepsilon(x)=\eta(x/\varepsilon)\)。其点值为一，梯度范数平方却是 \(\varepsilon^{d-2}\|\nabla\eta\|_2^2\to0\)。不存在同时控制这些点值的能量常数。

二维的简单缩放使梯度范数不变，须改用较缓的对数过渡。取 \(0<\varepsilon<1/2\)，令径向函数 \(w_\varepsilon\) 在 \(r\leq\varepsilon^2\) 等于一，在 \(\varepsilon^2<r<\varepsilon\) 等于 \(\log(\varepsilon/r)/\log(1/\varepsilon)\)，在 \(r\geq\varepsilon\) 等于零。它连续、分段光滑，且
\[
 \|\nabla w_\varepsilon\|_2^2
 =\frac{2\pi}{\log(1/\varepsilon)^2}
   \int_{\varepsilon^2}^{\varepsilon}\frac{\mathrm dr}{r}
 =\frac{2\pi}{\log(1/\varepsilon)}\longrightarrow0.
\]
以非负单位质量磨光核在半径小于 \(\varepsilon^2/2\) 的尺度卷积，得到光滑紧支集函数：中心邻域仍恒为一，支集仍在单位球内，Young 不等式使梯度范数不增。因此光滑测试函数同样违反 \(|\varphi(0)|\leq C\|\nabla\varphi\|_2\)。这正是否定点质量属于对偶所需的结论。
\end{proof}
一维点源的精确范数将在习题中进一步计算；高维的替代训练则计算小球平均源的能量如何发散。\Cref{b2:tab:15-source-classes}把本节已证的源项范围放在一起，便于应用后面的存在性定理前先核对数据。

\begin{table}[htbp]
\centering
\small
\caption[四类源项与能量对偶]{四类源项与能量对偶。固定整数 \(d\geq1\) 和非空有界开集 \(\Omega\subset\mathbb R^d\)，所有空间取实数域，\(H^{-1}(\Omega)=(H_0^1(\Omega;\mathbb R))^*\) 使用梯度对偶范数，\(C_P\) 为 \(\Omega\) 的零边界 Poincaré 常数。点质量两行另按所列维数和位置取数据；结论均已在正文证明。}
\label{b2:tab:15-source-classes}
\begin{tabularx}{\linewidth}{@{}>{\raggedright\arraybackslash}p{3.8cm}>{\centering\arraybackslash}p{1.5cm}>{\raggedright\arraybackslash}X@{}}
\toprule
数据形式及条件 & 属于 \(H^{-1}\) & 连续性估计或障碍 \\
\midrule
\(f\in L^2(\Omega)\) & 是 & \(\|f\|_{H^{-1}}\leq C_P\|f\|_2\)，使用零边界 Poincaré 不等式 \\
\(-\operatorname{div}F\)，\(F\in L^2(\Omega;\mathbb R^d)\) & 是 & \(\|-\operatorname{div}F\|_{H^{-1}}\leq\|F\|_2\)；不要求散度是普通函数 \\
\(d=1\)，\(\Omega=(a,b)\)，\(a<s<b\) 时的 \(\delta_s\) & 是 & \(|v(s)|\leq\sqrt{b-a}\|v'\|_2\)，点值取绝对连续代表 \\
\(d\geq2\)，\(x_0\in\Omega\) 时的 \(\delta_{x_0}\) & 否 & 测试函数点值固定为一，梯度能量可趋零；高于二维用缩放，二维用对数截断 \\
\bottomrule
\end{tabularx}
\end{table}

后面的零 Dirichlet 弱存在性定理只适用于已核验属于 \(H^{-1}\) 的数据，不能把任意点源自动代入；点质量的分布意义本身并不提供能量对偶中的连续性。

% !TeX root = ../../Book2.tex
\section{弱 \texorpdfstring{\(L^p\)}{Lᵖ} 空间}
\label{b2:sec:1702}

在端点问题中，函数超过高度 \(t\) 的区域往往比函数范数更容易估计。本节把高度与区域大小之间的关系写成函数空间语言。底空间可以是一般测度空间，重排部分以无原子的 Lebesgue 空间为例。

\subsection{分布函数}
\label{b2:subsec:170201}

设 \(f\) 可测，其\term[fenbu-hanshu]{分布函数}[distribution function]定义为 \(d_f(t)\coloneqq\mu\{|f|>t\}\)，\(t>0\)。它单调不增；对递减高度 \(t_n\downarrow t\)，集合 \(\{|f|>t_n\}\) 递增到 \(\{|f|>t\}\)，所以 \(d_f(t_n)\to d_f(t)\)。严格不等号保证这里的右连续性。
\symidx{\(d_f(t)\)：函数的分布函数}
\begin{proposition}[层蛋糕公式]\label{b2:thm:17-layer-cake}
对 \(0<p<\infty\) 和可测 \(f\)，有允许无穷值的恒等式
\[
 \int|f|^p\,\mathrm d\mu=p\int_0^\infty t^{p-1}d_f(t)\,\mathrm dt.
\]
\end{proposition}
\begin{proof}
逐点使用 \(|f(x)|^p=\int_0^\infty p t^{p-1}\mathbf1_{\{t<|f(x)|\}}\,\mathrm dt\)，再用 Tonelli 换序。非负性使此式在两端都无穷时仍有明确意义。
\end{proof}
在有限测度空间上，常函数的分布函数是一个阶梯函数；在 \(\mathbb R^d\) 上，\(|x|^{-d/p}\mathbf1_{\{|x|<1\}}\) 的高水平尾部则正好具有 \(t^{-p}\) 的尺度。这两类对象分别说明有界函数与临界奇性的分布图形。

\subsection{弱型 \texorpdfstring{\(L^p\)}{Lp}}
\label{b2:subsec:170202}

对 \(0<p<\infty\)，若
\[
 \|f\|_{p,\infty}\coloneqq\sup_{t>0}t\cdot d_f(t)^{1/p}<\infty,
\]
则称 \(f\) 属于\term[ruo-lp-kongjian]{弱 \(L^p\) 空间}[weak Lp space]，记为 \(L^{p,\infty}\)。这里给出的是拟范数：由 \(\{|f+g|>t\}\subset\{|f|>t/2\}\cup\{|g|>t/2\}\) 得到带常数的三角不等式，不宣称常数为一。
\begin{proposition}\label{b2:thm:17-weak-inclusions}
有 \(\|f\|_{p,\infty}\leq\|f\|_p\)。若 \(p>1\)、\(\mu(E)<\infty\)，则
\[
 \int_E|f|\,\mathrm d\mu\leq\dfrac{p}{p-1}\|f\|_{p,\infty}\cdot\mu(E)^{1-1/p}.
\]
反之，在半有限测度空间中，若所有有限测度 \(E\) 都满足右式中某个固定常数的估计，则 \(f\in L^{p,\infty}\)。
\end{proposition}
\begin{proof}
第一项由 Chebyshev 得到。记 \(A=\|f\|_{p,\infty}\)、\(m=\mu(E)\)。在 \(m=0\) 或 \(A=0\) 时结论直接成立；否则按层蛋糕公式，积分不超过 \(\int_0^\infty\min\{m,A^pt^{-p}\}\,\mathrm dt\)。在 \(t=A m^{-1/p}\) 分开积分即得。反向在 \(E=\{|f|>t\}\) 上使用 \(t\mu(E)\leq\int_E|f|\)，若集合测度无穷，先取其中有限测度子集；在半有限空间上便推出 \(t\mu(E)^{1/p}\leq C\)。本卷的 Lebesgue 空间满足这一条件。
\end{proof}
积分集合刻画解释了为什么 \(p>1\) 时弱空间中的函数仍局部可积，而 \(p=1\) 时不能作同样断言。

\ProseParagraphBreak

若算子 \(T\) 满足 \(\|Tf\|_{q,\infty}\leq A\|f\|_p\)，称它为\term[ruoxing-guji]{弱型估计}[weak type estimate]的 \((p,q)\) 版本；强型把目标换为 \(L^q\)。当 \(q=\infty\) 时这里直接使用本性上确界，不另定义分布函数拟范数。

% 扩写定位：b2:subsec:170202
在 Lebesgue 空间上，集合积分刻画还给出一个真正的范数。设 \(1<p<\infty\)，定义
\[
 N_p(f)\coloneqq\sup_{0<|E|<\infty}|E|^{-1/p'}\int_E|f|.
\]
这里可以只取有界可测集合：一般集合与球相交后逐步逼近即可。由上一命题，
\[
 \|f\|_{p,\infty}\leq N_p(f)\leq p'\|f\|_{p,\infty}.
\]
左侧来自超水平集合及其有限测度子集。绝对齐次性、三角不等式和 \(N_p(f)=0\) 蕴含 \(f=0\) 几乎处都直接由定义得到。因而在这个指数范围内，弱空间可以赋予与原拟范数等价的 Banach 范数。
\begin{proposition}\label{b2:prop:17-weak-complete}
\(L^{p,\infty}(\mathbb R^d)\) 关于 \(N_p\) 完备。若 \(k\in L^1\)，则卷积具有连续延拓，并满足
\[
 N_p(k*f)\leq\|k\|_1N_p(f).
\]
\end{proposition}
\begin{proof}
设 \(f_n\) 在 \(N_p\) 中 Cauchy。在每个有限球 \(B\) 上，
\(\|f_n-f_m\|_{L^1(B)}\leq |B|^{1/p'}N_p(f_n-f_m)\)，故得到相容的局部 \(L^1\) 极限 \(f\)。取沿所有整数球几乎处收敛的对角子列。对任意有限测度集合 \(E\)，Fatou 给出
\(\int_E|f_n-f|\leq\liminf_m\int_E|f_n-f_m|\)；先固定足够大的 \(n\)，再对 \(E\) 取上确界，得到 \(N_p(f_n-f)\to0\)。这同时证明 \(f\) 的范数有限。

对卷积，Tonelli 和平移不变性给出
\[
 \int_E\int |k(y)|\,|f(x-y)|\,\mathrm dy\,\mathrm dx
 \leq |E|^{1/p'}\|k\|_1N_p(f).
\]
右端有限，故积分在每个有限球上几乎处绝对收敛；其绝对值控制 \(|k*f|\)，再取上确界便得结论。
\end{proof}
例如，临界势核的弱范数可以在与一个可积平滑核卷积后保持同一指数。这里并不需要卷积核正定，也不要求 \(f\) 本身可积；用局部集合积分而不是全空间 \(L^1\) 积分，正好避免了这一障碍。

\ProseParagraphBreak

完备性不能误读成平滑紧支集函数在弱范数下稠密。令 \(f(x)=|x|^{-d/p}\mathbf1_{\{|x|<1\}}\)，并令 \(g\) 有界。在足够小的球上，\(|f-g|\geq |f|/2\)；高水平集合的尺度于是给出 \(\|f-g\|_{p,\infty}\geq c_{d,p}>0\)。有界截断 \(f\mathbf1_{\{|f|\leq N\}}\) 因此也不收敛到 \(f\) 的弱范数。处理算子的弱空间延拓时，需要使用已证明的积分定义，或者指定额外的收敛方式，不能照搬强 \(L^p\) 的稠密性论证。

\subsection{Lorentz 空间 初步}
\label{b2:subsec:170203}

定义递减重排 \(f^*(s)\coloneqq\inf\{t\geq0:d_f(t)\leq s\}\)，\(s>0\)，下确界为空时取无穷。它把函数值按大小重新排列，而不保留位置。\term[lorentz-kongjian]{Lorentz 空间}[Lorentz space]通过
\[
 \|f\|_{L^{p,q}}\coloneqq
 \left(\int_0^\infty \bigl(s^{1/p}f^*(s)\bigr)^q\dfrac{\mathrm ds}{s}\right)^{1/q}
 \quad(0<q<\infty)
\]
衡量整个高度分布，\(q=\infty\) 时改为上确界。这里仍用拟范数约定。重排与原函数具有相同分布，故层蛋糕公式给出 \(\|f\|_{L^{p,p}}=\|f\|_p\)，而 \(L^{p,\infty}\) 与上段定义一致。
\symidx{\(f^*(s)\)：递减重排}

\ProseParagraphBreak

对 \(1<p<\infty\)，若 \(f\in L^{p,1}\)、\(g\in L^{p',\infty}\)，则
\[
 \int|fg|\leq\|f\|_{L^{p,1}}\cdot\|g\|_{L^{p',\infty}}.
\]
确实，写 \(f=\int_0^\infty\mathbf1_{\{|f|>t\}}\,\mathrm dt\)，集合积分的重排界为 \(\int_E|g|\leq\int_0^{|E|}g^*(s)\,\mathrm ds\)。这个界由层蛋糕和 \( |E\cap\{|g|>u\}|\leq\min\{|E|,d_g(u)\}\) 得到。换序后得到 \(\int|fg|\leq\int_0^\infty f^*(s)g^*(s)\,\mathrm ds\)，再代入 \(g^*(s)\leq\|g\|_{p',\infty}s^{-1/p'}\)。这是把一个临界弱估计与更细的输入空间配对的基本方式。

\ProseParagraphBreak

例如令 \(f^*(s)\) 在充分小的 \(s>0\) 上等于 \(s^{-1/p}\bigl(\log(e/s)\bigr)^{-\beta}\)，其余部分取单调有限支集延拓。其 \(L^{p,q}\) 积分在 \(s=0\) 收敛当且仅当 \(\beta q>1\)，弱 \(L^p\) 只要求 \(\beta\geq0\)。第二个指数记录了强空间 \(L^p\) 所不能区分的端点细节。

% 扩写定位：b2:subsec:170203
几种分布尾部可以帮助区分这些空间。首先，函数在有限测度集合上有界，并不能说明其重排在大 \(s\) 上怎样；当函数支集也有限时，重排在支集测度之后为零。其次，一个幂次奇性在小球上产生高水平的幂次尾部，决定局部的可积临界；同一幂次放在球外，则产生低水平的大测度尾部，决定远场的可积临界。两种限制可以落在相反方向，因此全空间上的空间包含关系不能只看一个局部模型。

\ProseParagraphBreak

具体地，若 \(f(x)=|x|^{-a}\mathbf1_{\{|x|<1\}}\)，径向积分说明 \(f\in L^p\) 当且仅当 \(ap<d\)，临界 \(ap=d\) 只有弱 \(L^p\)。若 \(g(x)=|x|^{-a}\mathbf1_{\{|x|>1\}}\)，则 \(g\in L^p\) 当且仅当 \(ap>d\)，临界同样只有弱空间。把两部分加在一起得到纯幂次 \(|x|^{-a}\)；它属于临界弱空间，却不属于任何全空间强 \(L^p\)。这不是弱空间定义的反常现象，而是它准确记录了在每个高度上超水平球的尺度。

\ProseParagraphBreak

有时输入的第二指数需要比弱空间更严格。设 \(f^*(s)\) 在 \(s\downarrow0\) 时具有下列量级：
\[
 f^*(s)\asymp s^{-1/p}(\log(e/s))^{-\beta}.
\]
在其余部分作单调的有限支集延拓。代入 Lorentz 定义，换元 \(v=\log(e/s)\)，所得尾积分为 \(\int^\infty v^{-\beta q}\,\mathrm dv\)。因此固定 \(p\) 后，\(q\) 越小要求越强；例如在 \(\beta=1/p\) 时恰好不属于 \(L^p\)，但可能属于具有更大第二指数的空间。这里说的是原点附近的量级，若 \(\beta\) 大，原显示公式在整个区间上未必递减，应以这种局部模型加单调延拓来实现重排。

\ProseParagraphBreak

前述端点配对中，\(L^{p,1}\) 一般不能放宽为 \(L^p\)。取临界重排
\[
 g^*(s)=s^{-1/p'}\mathbf1_{\{0<s<1\}}.
\]
它属于弱 \(L^{p'}\)。若 \(f^*(s)=s^{-1/p}(\log(e/s))^{-\beta}\) 且 \(1/p<\beta\leq1\)，则 \(f\in L^p\)，但同位置递减模型的乘积积分归结为
\[
 \int_0^1s^{-1}(\log(e/s))^{-\beta}\,\mathrm ds=\infty.
\]
取 \(\beta>1\) 则 \(f\in L^{p,1}\)，前面证明的配对界保证收敛。于是把 \(L^{p,1}\) 换成 \(L^p\) 一般会破坏结论，这个反例直接指出第二指数在临界积分中的作用。

\ProseParagraphBreak

在有限测度区域 \(E\) 上，若 \(0<q<p\)，同样按层蛋糕把高度在 \(\|f\|_{p,\infty}|E|^{-1/p}\) 处分开，可得
\[
 \|f\|_{L^q(E)}
 \leq\left(\frac p{p-q}\right)^{1/q}|E|^{1/q-1/p}\|f\|_{p,\infty}.
\]
这不是弱 \(L^p\) 的强端点；它降低了目标指数，同时用到了区域测度。若 \(q=p\)，高度积分再次成为对数发散；若 \(|E|=\infty\)，所显示的区域因子不提供有限上界。因而使用这一局部提升时，指数严格性与区域有限性都必须保留。

\ProseParagraphBreak

还有一种较弱的输入端点，只先在集合指示函数上检验算子。若线性算子 \(T\) 在有限支集简单函数上定义，并满足
\(\|T\mathbf1_E\|_{q,\infty}\leq A|E|^{1/p}\)、\(1\leq p<\infty\)、\(1<q<\infty\)，称这是一项限制弱型估计。它并不直接给出全部 \(L^p\) 输入的弱型界，但足以推出
\[
 N_q(Tf)\leq C_{p,q}A\|f\|_{L^{p,1}}.
\]
对非负简单函数，层蛋糕是有限个集合指示函数的正线性组合。范数三角不等式与弱范数等价给出
\[
 N_q(Tf)\leq q'A\int_0^\infty d_f(t)^{1/p}\,\mathrm dt
 =\frac{q'A}{p}\|f\|_{L^{p,1}}.
\]
最后的等式由 Tonelli 对重排的层蛋糕积分得到。一般复简单函数分解为实部与虚部的正负部分，常数只增加一个固定因子。

\ProseParagraphBreak

有限支集简单函数在 \(L^{p,1}\) 中稠密：先截去位置尾部、大值尾部和极小值尾部，重排积分的可积性保证误差趋零，再在剩余有界有限测度区域作简单逼近。因 \(N_q\) 完备，以上估计产生唯一连续延拓。这个结果说明限制弱型端点自然对应更细的 \(L^{p,1}\) 输入，而不能在没有新论证的情况下换成所有 \(L^p\)。当目标 \(q=1\) 时，本段使用的等价 Banach 范数已不存在，仍需另一种处理。

% !TeX root = ../../Book2.tex
\section{Poisson 方程}
\label{b2:sec:1503}

对 Poisson 方程，把二阶导数经分部积分转到测试函数上，就得到关于梯度的双线性等式。Lax--Milgram 定理可以在有界区域的能量空间中求解这一等式，并允许源项属于 \(H^{-1}\)。这种写法将存在性、稳定性与边界数据放在同一个框架中。

本节使用\subsecref{b2:subsec:150104}的实空间约定：\(H^{-1}(\Omega)\) 是带梯度范数的 \(H_0^1(\Omega;\mathbb R)\) 的连续对偶，\(\langle f,v\rangle\) 表示对偶配对。下面常数为一的估计与这一梯度范数约定相配；换用完整 \(H^1\) 范数时，对偶空间不变，数值常数则须作相应比较。

\subsection{弱解存在性}
\label{b2:subsec:150301}

\Needspace{6\baselineskip}
\begin{theorem}[Dirichlet 问题的弱可解性]
\label{b2:thm:15-poisson-existence}
设整数 \(d\geq1\)，\(\Omega\subset\mathbb R^d\) 为非空有界开集，空间取实数域。以梯度范数赋予 \(H_0^1(\Omega)\) 范数，令 \(H^{-1}(\Omega)=(H_0^1(\Omega))^*\)，其对偶范数为 \(\|f\|_{H^{-1}}=\sup_{v\in H_0^1,\,\|\nabla v\|_2\leq1}|\langle f,v\rangle|\)。对每个 \(f\in H^{-1}(\Omega)\)，存在唯一 \(u\in H_0^1(\Omega)\)，使
\[\int_\Omega\nabla u\cdot\nabla v\,\mathrm dx=\langle f,v\rangle
 \qquad(v\in H_0^1(\Omega)).\]
并且 \(\lVert\nabla u\rVert_2\le\lVert f\rVert_{H^{-1}}\)。
\end{theorem}
\begin{proof}
取 \(H=H_0^1(\Omega)\)，内积为 \(\langle u,v\rangle_H=\int\nabla u\cdot\nabla v\)。第一卷定理24.1.2的零边界 Poincaré 不等式适用于任意包含于有限宽条带的开集，故适用于这里的有界 \(\Omega\)，不要求边界正则或连通。它使梯度范数与完整 Sobolev 范数等价，故 \(H\) 完备。型 \(a(u,v)=\langle u,v\rangle_H\) 的有界常数与强制常数均为 \(1\)，右端按 \(H^{-1}\) 定义连续。\cref{b2:thm:15-lax-milgram} 因而给出存在、唯一性及估计。
\end{proof}

本例中 Lax--Milgram 退化为 Riesz 表示；一般系数时，同一个定理才显出更大的作用。该弱解也恰为 \(\tfrac12\lVert\nabla v\rVert_2^2-\langle f,v\rangle\) 的唯一极小元：只要将上一章的二次展开中的 \(L^2\) 配对换成对偶配对，证明逐字成立。两种构造在对称情形相遇，但弱型方法还可以继续处理非对称问题。

\ProseParagraphBreak

对 \(\Omega=(0,1)\)、\(f=\delta_s\)、\(0<s<1\)，解为
\[
 u(x)=\begin{cases}(1-s)x,&0\le x\le s,\\s(1-x),&s\le x\le1.\end{cases}
\]
它连续、分段仿射，导数在 \(s\) 处向下跳 \(1\)。两段分部积分给出 \(\int_0^1u'v'=v(s)\)，先对光滑 \(v\) 成立，再由点值的连续性推广。因此点源产生一个合法弱解，却不产生经典 \(C^2\) 解；这不是存在性定理的缺陷，而是数据本身的粗糙性。

\subsection{唯一性}
\label{b2:subsec:150302}

唯一性也可以直接从能量等式看出。若 \(u_1,u_2\) 具有同一数据，差 \(w\in H_0^1\) 满足 \(\int\nabla w\cdot\nabla v=0\)。取 \(v=w\)，得到 \(\nabla w=0\)；Poincaré 不等式推出 \(w=0\)。这个结论对不连通区域也成立，因为零边界约束在每个分支上排除了常数自由度。

\ProseParagraphBreak

内部方程则没有唯一性。例如任一调和函数都可加到 Poisson 方程的一个局部解上，而不改变源项。纯 Neumann 问题即使规定通量，仍允许每个连通分支上的加法常数。因而“取差再测试”最后得出梯度为零之后，还必须指出究竟由哪项边界条件或归一化排除了常数，不能直接省略最后一步。

\Needspace{6\baselineskip}
\begin{proposition}[弱比较原理]
\label{b2:prop:15-weak-comparison}
设整数 \(d\geq1\)，\(\Omega\subset\mathbb R^d\) 为非空有界开集，空间取实数域，\(H^{-1}(\Omega)=(H_0^1(\Omega))^*\)。设 \(f_1,f_2\in H^{-1}(\Omega)\)，\(u_1,u_2\in H_0^1(\Omega)\) 满足
\[
 \int_\Omega\nabla u_i\cdot\nabla v=\langle f_i,v\rangle
 \quad(v\in H_0^1(\Omega),\ i=1,2).
\]
若 \(\langle f_1-f_2,v\rangle\le0\) 对每个几乎处处非负的 \(v\in H_0^1(\Omega)\) 成立，则 \(u_1\le u_2\) 几乎处处。
\end{proposition}
\begin{proof}
令 \(w=u_1-u_2\)。第一卷命题22.2.3保证取正部保持实空间 \(H_0^1\)，命题21.3.5给出相应弱梯度公式。因此 \(w_+=\max\{w,0\}\in H_0^1\)，且 \(\nabla w_+=\mathbf1_{\{w>0\}}\nabla w\)。以它测试差方程，得
\[\lVert\nabla w_+\rVert_2^2=\langle f_1-f_2,w_+\rangle\le0.\]
零边界 Poincaré 不等式推出 \(w_+=0\)，即所需结论。所引截断结果的关键是先光滑逼近正部函数，在零水平集上用弱梯度为零处理折点，再利用正部在 \(H^1\) 中的强连续性传递零边界逼近；没有要求解逐点可微。
\end{proof}

这里的非负性还须在数据空间中解释。对非空有界开集 \(\Omega\)，实泛函 \(f\in H^{-1}(\Omega)\) 满足 \(f\geq0\)，是指
\[
 \langle f,v\rangle\geq0
 \quad\text{对所有几乎处处非负的 }v\in H_0^1(\Omega;\mathbb R).
\]
它是对测试函数作用的正性，不是给一般分布指定逐点值。若 \(f\in L^2(\Omega;\mathbb R)\) 且 \(f\geq0\) 几乎处处，已建立的 \(L^2\hookrightarrow H^{-1}\) 将配对写成可积的 \(\int fv\)，因而符合这个定义。不能只写 \(f\in L^1\) 就断言它属于 \(H^{-1}\)，也不能默认 \(fv\) 对全部 \(H_0^1\) 测试函数可积。将弱比较原理用于零源项与非负源项，得到非负数据的零边界弱解非负，与\chapref{b2:ch:13}经典最大值原理的符号一致。比较结论是几乎处处关系；尚未得到连续代表元时，不应改说“每个点都非负”。

\subsection{能量估计与 Galerkin 近似}
\label{b2:subsec:150303}

取解自身作测试，Cauchy--Schwarz 的对偶形式给出
\[\lVert\nabla u\rVert_2^2=\langle f,u\rangle
 \le\lVert f\rVert_{H^{-1}}\lVert\nabla u\rVert_2.\]
若梯度非零可约去一因子；若为零估计显然成立。因此解映射 \(H^{-1}\to H_0^1\) 连续。对于两组数据，同一计算给出
\[\lVert\nabla(u_1-u_2)\rVert_2\le\lVert f_1-f_2\rVert_{H^{-1}}.\]
特别地，平滑数据的近似会在能量范数中产生相应解的收敛，这为从光滑近似过渡到粗糙数据提供了稳定性依据。

\ProseParagraphBreak

若 \(f\in L^2\)，进一步有 \(\lVert\nabla u\rVert_2\le C_P\lVert f\rVert_2\)、\(\lVert u\rVert_2\le C_P^2\lVert f\rVert_2\)。注意区域放大时 \(C_P\) 也会变化：在 \(r\Omega\) 上换元表明 \(C_P\) 按长度尺度增长。这与二阶方程解相对于源项多出两个长度因子的事实一致，估计中的几何依赖不能无故删除。

\Needspace{6\baselineskip}
\begin{proposition}[有限维近似的最佳误差]
\label{b2:prop:15-galerkin-poisson}
设整数 \(d\geq1\)，\(\Omega\subset\mathbb R^d\) 为非空有界开集，空间取实数域，\(H^{-1}(\Omega)=(H_0^1(\Omega))^*\)。给定 \(f\in H^{-1}(\Omega)\) 与有限维线性子空间 \(V\subset H_0^1(\Omega)\)，设 \(u\in H_0^1(\Omega)\) 满足 \(\int_\Omega\nabla u\cdot\nabla v=\langle f,v\rangle\) 对全部 \(v\in H_0^1(\Omega)\) 成立。则存在唯一 \(u_V\in V\) 满足 \(\int_\Omega\nabla u_V\cdot\nabla v=\langle f,v\rangle\) 对所有 \(v\in V\) 成立，而且
\[\lVert\nabla(u-u_V)\rVert_2=\inf_{v\in V}\lVert\nabla(u-v)\rVert_2.\]
\end{proposition}
\begin{proof}
限制到 \(V\) 后型仍有界、强制，故有限维方程唯一可解。两式相减得 \(\int\nabla(u-u_V)\cdot\nabla v=0\) 对所有 \(v\in V\) 成立。对任意 \(v\in V\)，分解 \(u-v=(u-u_V)+(u_V-v)\)，两项在梯度内积下正交，故误差平方等于两项范数平方之和。最小值恰在 \(v=u_V\) 达到。
\end{proof}

近似解的误差受该有限维空间中的最佳逼近误差控制；因此，选取怎样的近似空间至关重要。若一列子空间递增且其并在 \(H_0^1\) 中稠密，最佳误差便趋零。递增性可改成对每个固定函数都能沿整个序列逼近的条件，但只有并集稠密不够，例如在良好空间之间反复插入零空间便会破坏收敛。存在性、离散方程的可解性与近似空间的逼近能力仍是三个可以分别核验的事实。

\ProseParagraphBreak

Poisson 型中，误差与近似空间真正正交，因而得到常数为一的最佳逼近。一般非对称型仍有从两个弱方程相减得到的 Galerkin 正交性，但它不能随意交换两个变量。下面的\term[cea-yinli]{Céa 引理}[Céa's lemma]用有界性和强制性代替勾股恒等式，给出准最佳逼近。

\Needspace{6\baselineskip}
\begin{proposition}[Céa 引理]\label{b2:prop:15-cea}
设 \(H\) 为实 Hilbert 空间，\(a:H\times H\to\mathbb R\) 为双线性型，存在 \(0\leq M<\infty\)、\(\alpha>0\)，使
\[
 |a(x,y)|\leq M\|x\|_H\cdot\|y\|_H,\qquad
 a(x,x)\geq\alpha\|x\|_H^2\quad(x,y\in H).
\]
给定 \(\ell\in H^*\) 及闭线性子空间 \(V\subset H\)，令 \(u\in H\)、\(u_V\in V\) 分别满足
\[
 a(u,y)=\ell(y)\quad(y\in H),\qquad
 a(u_V,v)=\ell(v)\quad(v\in V).
\]
则两解都存在且唯一，并且
\[
 \|u-u_V\|_H\leq\frac M\alpha\inf_{v\in V}\|u-v\|_H.
\]
特别地，每个有限维线性子空间均满足这里的闭性要求。
\end{proposition}
\begin{proof}
闭子空间 \(V\) 以继承的内积仍为 Hilbert 空间，限制后的型保留有界性与同一个强制常数，\(\ell|_V\) 连续。分别在 \(H\) 与 \(V\) 上使用\cref{b2:thm:15-lax-milgram}，得到两解。令 \(e=u-u_V\)，相减得到 \(a(e,v)=0\) 对所有 \(v\in V\) 成立。对任意 \(v\in V\)，因 \(v-u_V\in V\)，有
\[
 \alpha\|e\|_H^2\leq a(e,e)
 =a(e,u-v)+a(e,v-u_V)=a(e,u-v)
 \leq M\|e\|_H\cdot\|u-v\|_H.
\]
若 \(e=0\)，结论成立；否则约去 \(\|e\|_H\)，再对 \(v\in V\) 取下确界。误差关系 \(a(e,v)=0\) 对第二变量中的 \(v\in V\) 成立，这就足以得到同样适用于非对称型的估计。
\end{proof}

因此，强制型的离散解能否沿一列空间收敛，仍由这些空间能否逼近真解决定；有界性与强制性只增加了准最佳常数 \(M/\alpha\)。下面的分段仿射计算采用 Poisson 对称型，可以继续使用前一命题中常数为一的精确最佳逼近。

\ProseParagraphBreak

在区间 \((0,1)\) 上，这个结论可以算得很具体。取整数 \(N\geq2\)、均匀节点 \(x_j=jh\)（\(0\leq j\leq N\)）、\(h=1/N\)，令 \(V_h\) 为节点间分段仿射、两端为零的连续函数空间。基函数 \(\phi_j\)（\(1\leq j\leq N-1\)）在 \(x_j\) 取值一，在其余节点取值零。每个基函数只有相邻两段上的斜率非零，直接积分得
\[\int_0^1\phi_j'^2=\frac2h,\qquad
 \int_0^1\phi_j'\phi_{j+1}'=-\frac1h,\qquad
 \int_0^1\phi_j'\phi_k'=0\quad(|j-k|>1).\]
因此弱解的有限维近似由三对角线性方程组确定。矩阵的正定性并不需要另作行列式计算：任意系数向量对应的二次型就是分段仿射函数的梯度平方积分，零值只可能对应零函数。

\ProseParagraphBreak

例如取 \(f=1\)，精确弱解为 \(u(x)=x(1-x)/2\)。它在各小区间的导数平均等于两端函数值之差除以 \(h\)，而任意 \(v_h\in V_h\) 的导数在该段为常数。故 \(u\) 的节点线性插值 \(I_hu\) 满足 \(\int(u'- (I_hu)')v_h'=0\)，由唯一性即知 \(u_h=I_hu\)。在长度 \(h\) 的一段上，导数误差是以中点为零的线性函数，其平方积分为 \(h^3/12\)；函数误差则为 \((x-x_j)(x_{j+1}-x)/2\)。逐段积分并相加得到
\[\lVert u'-u_h'\rVert_2=\frac h{\sqrt{12}},\qquad
 \lVert u-u_h\rVert_2=\frac{h^2}{\sqrt{120}}.\]
这里第二个收敛速度比第一个快，来自二次函数的具体结构，不能仅从抽象存在定理读出。一般数据下的误差阶还需要解的正则性和近似空间的逼近估计；下一章提供其中的正则性工具。

\ProseParagraphBreak

弱解的误差也可以用残差直接测量。设整数 \(d\geq1\)、\(\Omega\subset\mathbb R^d\) 为非空有界开集，实值 \(f\in H^{-1}(\Omega)\)，\(u\in H_0^1(\Omega)\) 满足 \(\int_\Omega\nabla u\cdot\nabla v=\langle f,v\rangle\) 对所有实值 \(v\in H_0^1(\Omega)\) 成立。对任意实值 \(w\in H_0^1(\Omega)\) 定义 \(r=f+\Delta w\in H^{-1}\)，则 \(\langle r,v\rangle=\int\nabla(u-w)\cdot\nabla v\)。以 \(H_0^1\) 的梯度范数取对偶范数，对单位梯度范数的测试函数取上确界，Cauchy--Schwarz 给出上界，而在误差非零时取 \(v=(u-w)/\lVert\nabla(u-w)\rVert_2\) 达到该界。因此 \(\lVert r\rVert_{H^{-1}}=\lVert\nabla(u-w)\rVert_2\)。对分段仿射近似，残差可能含节点通量跳跃，不能只计算每段内为零的二阶导数；负范数恰好保留这些分布贡献。

\subsection{非齐次 Dirichlet 数据与 Neumann 问题}
\label{b2:subsec:150304}

取整数 \(d\geq1\)、非空有界开集 \(\Omega\subset\mathbb R^d\)、实值 \(f\in H^{-1}(\Omega)\) 与 \(g\in H^1(\Omega;\mathbb R)\)。令 \(u=g+w\)，其中 \(w\in H_0^1(\Omega;\mathbb R)\)。方程等价于
\[\int_\Omega\nabla w\cdot\nabla v
  =\langle f,v\rangle-\int_\Omega\nabla g\cdot\nabla v\qquad(v\in H_0^1).\]
右端属于 \(H^{-1}\)，范数不超过 \(\lVert f\rVert_{H^{-1}}+\lVert\nabla g\rVert_2\)。因此齐次理论给出唯一 \(w\)，从而给出唯一满足 \(u-g\in H_0^1\) 的解；并有
\[\lVert\nabla w\rVert_2\le\lVert f\rVert_{H^{-1}}+\lVert\nabla g\rVert_2.\]
完整 \(H^1\) 估计还要加上 \(g\) 的 \(L^2\) 范数。

\ProseParagraphBreak

若另一个延拓 \(\widetilde g\) 与 \(g\) 相差 \(H_0^1\) 函数，两者给出的仿射可行集完全相同；由唯一性得到同一个 \(u\)。所以解依赖边界类，而不依赖人为选择的延拓。若要比较两组边界数据，先选择它们差的一个延拓，再把上面的稳定性估计应用于差方程；优化延拓的范数便自然引出商空间形式的迹范数。

\ProseParagraphBreak

例如在 \(\mathbb R^d\) 的单位球上给定仿射边界数据 \(g(x)=a\cdot x+b\)，其中 \(a\in\mathbb R^d\)、\(b\in\mathbb R\)，且 \(f=0\)。由于 \(\int a\cdot\nabla v=0\) 对 \(v\in H_0^1\) 成立，齐次修正为零，故 \(u=g\)。对一般边界数据，弱方法也能构造能量解；在球上还可借\chapref{b2:ch:13}的表示公式识别这个解。不同构造的解一旦处于同一能量边界类，就由唯一性自动一致。


\ProseParagraphBreak

\noindent\textbf{Neumann 问题。}区域保持固定而边界值自由时，应换一个 Hilbert 空间，不能沿用零 Dirichlet 条件排除常数的唯一性论证。

\Needspace{6\baselineskip}
\begin{proposition}[Neumann 数据与归一化]
\label{b2:prop:15-neumann-normalized}
设整数 \(d\geq1\)，\(\Omega\subset\mathbb R^d\) 为非空有界连通 Lipschitz 区域，\(f\in L^2(\Omega;\mathbb R)\)、\(g_N\in L^2(\partial\Omega;\mathbb R)\)，边界积分取 \((d-1)\) 维面积测度（\(d=1\) 时为端点计数测度），\(\operatorname{Tr}\) 为 Sobolev 迹，求解类为 \(u\in H^1(\Omega;\mathbb R)\)。方程
\[\int_\Omega\nabla u\cdot\nabla v
 =\int_\Omega fv+\int_{\partial\Omega}g_N\operatorname{Tr}v
 \quad(v\in H^1(\Omega))\]
可解当且仅当 \(\int_\Omega f+\int_{\partial\Omega}g_N=0\)。指定 \(\int_\Omega u=0\) 后解唯一，且
\[\|u\|_{H^1}\le C_\Omega(\|f\|_2+\|g_N\|_{L^2(\partial\Omega)}).\]
\end{proposition}
\begin{proof}
测试 \(v=1\) 给出必要条件。令 \(H=\{v\in H^1(\Omega;\mathbb R):\int_\Omega v=0\}\)。积分泛函因 \(\Omega\) 测度有限而连续，故 \(H\) 是闭子空间；第一卷定理24.1.3的 Poincaré--Wirtinger 不等式适用于这里的有界连通 Lipschitz 区域，使 \(\|\nabla v\|_2\) 成为 \(H\) 上等价的完备范数。右端有界性来自
\[
 |\ell(v)|\le\|f\|_2\|v\|_2+
 \|g_N\|_{L^2(\partial\Omega)}\|\operatorname{Tr}v\|_{L^2(\partial\Omega)}
 \le C_\Omega(\|f\|_2+\|g_N\|_2)\|\nabla v\|_2,
\]
这里第一卷定理23.1.3给出有界 Lipschitz 区域上的 \(H^1\to L^2(\partial\Omega)\) 迹估计，刚才的零均值不等式再将 \(H^1\) 范数换成梯度范数。Lax--Milgram 在 \(H\) 上给出唯一解及估计。对任意 \(v\in H^1(\Omega)\)，分解 \(v=v_0+v_\Omega\)，其中 \(v_\Omega=|\Omega|^{-1}\int_\Omega v\)、\(v_0\in H\)。左端常数项为零，右端常数项因相容条件为零，所以受限测试解确实解全部测试方程。两解之差测试自身，得到梯度为零；连通性说明差为一个常数，归一化排除它。
\end{proof}

在 \((0,1)\) 上取 \(f=1\)、左端外向通量 \(g_N(0)=-1\)、右端通量 \(g_N(1)=0\)。总数据为 \(1-1+0=0\)。积分 \(-u''=1\) 得 \(u'=-x+C\)；左端条件 \(-u'(0)=-1\) 给 \(C=1\)，右端也满足。减去均值后 \(u=x-x^2/2-1/3\)。这个例子同时检验了相容条件中的加号和左端外法向的负号。若区域有多个连通分支，必须在每个分支分别相容、分别归一化，单个全局零均值无法排除所有常数自由度。

% !TeX root = ../../Book2.tex
\section{Riesz–Thorin 插值}
\label{b2:sec:1704}

线性算子可沿复参数改变输入和输出的指数，从而得到具有精确几何平均常数的插值估计。为使用这项方法，需要复变函数的最大模原理；本节把所需的条带估计单独证明。

\subsection{解析族}
\label{b2:subsec:170401}

下面采用\term[hadamard-sanxian-yinli]{Hadamard 三线引理}[Hadamard three-lines lemma]的下述版本。

\begin{lemma}[Hadamard 三线估计]\label{b2:thm:17-three-lines}
设 \(F\) 在 \(0<\operatorname{Re}z<1\) 全纯，在闭条带连续且有界，边界满足 \( |F(\mathrm it)|\leq M_0\)、\( |F(1+it)|\leq M_1\)。则 \(0<\theta<1\) 时
\[
 |F(\theta)|\leq M_0^{1-\theta}M_1^\theta.
\]
\end{lemma}
\begin{proof}
先设 \(M_0,M_1>0\)，令 \(G(z)=F(z)M_0^{z-1}M_1^{-z}\)，幂以实对数定义。对 \(\varepsilon>0\)，函数 \(G(z)e^{\varepsilon(z^2-1)}\) 在两条竖边的模不超过一；在上下水平边 \(\operatorname{Im}z=\pm R\) 上，指数的模不超过 \(e^{-\varepsilon R^2}\)，故当 \(R\) 足够大时也不超过一。对矩形应用最大模原理，再令 \(R\to\infty\)、\(\varepsilon\downarrow0\)。某个 \(M_i=0\) 时，先以 \(M_i+\delta\) 替代，令 \(\delta\downarrow0\)。
\end{proof}
一个\term[jiexi-zu]{解析族}[analytic family]是复参数下的算子族，其对测试输入、测试输出的矩阵元为全纯函数。使用三线估计时必须检查边界范数和条带内增长；“公式含有 \(z\)”不是解析族的充分条件。下面先用固定算子和变化的测试函数，避免讨论一般算子的定义域变化。

\subsection{Riesz–Thorin 定理}
\label{b2:subsec:170402}

\begin{theorem}[Riesz--Thorin]\label{b2:thm:17-riesz-thorin}
设 \(T\) 是定义在两个端点 \(L^{p_0}(X)+L^{p_1}(X)\) 上的兼容线性算子，底空间为 \(\sigma\)-有限测度空间。设 \(1\leq p_i,q_i\leq\infty\)，且 \(\|Tf\|_{q_i}\leq A_i\|f\|_{p_i}\)。若
\[
 \dfrac1p=\dfrac{1-\theta}{p_0}+\dfrac\theta{p_1},\qquad
 \dfrac1q=\dfrac{1-\theta}{q_0}+\dfrac\theta{q_1},
\]
则 \(\|Tf\|_q\leq A_0^{1-\theta}A_1^\theta\|f\|_p\)。
\end{theorem}
\begin{proof}
先取有限测度支集上的简单函数 \(f,g\)，分别归一化为 \(\|f\|_p=\|g\|_{q'}=1\)，并先讨论 \(p,q'<\infty\)。在非零值处写
\[
 f_z=\dfrac f{|f|}\,|f|^{p((1-z)/p_0+z/p_1)},\quad
 g_z=\dfrac g{|g|}\,|g|^{q'((1-z)/q_0'+z/q_1')}.
\]
零值处定义为零。它们是有限个固定简单函数的解析系数组合。于是 \(F(z)=\int(Tf_z)g_z\) 全纯；在竖边 \(i=0,1\) 上有 \(\|f_{i+it}\|_{p_i}=\|g_{i+it}\|_{q_i'}=1\)，无穷指数时对应模恒为一。端点 Hölder 得到边界常数 \(A_i\)。条带内各系数的模只依赖实部，因此 \(F\) 有界。三线估计在 \(z=\theta\) 给出双线性配对的上界。

对所有简单 \(g\) 取上确界，\(L^q\) 对偶范数公式得到 \(\|Tf\|_q\) 的估计；该公式也适用于 \(q=1,\infty\)，只需用简单函数近似符号或在接近本性上确界的有限测度集合上测试。若 \(p=\infty\)，两个源端点都是无穷，直接对固定 \(f_z=f\) 作同样论证，或者用有限端点目标指数的对偶插值。一般有限 \(p\) 输入由简单函数稠密性延拓；原端点定义与延拓在公共部分相同，因而没有另造一个不兼容的算子。
\end{proof}
对解析算子族 \(T_z\)，在相同构造中把 \(T\) 换为 \(T_z\)，只要所有测试矩阵元在条带中满足同样的增长条件，结论相同。若增长至多 \(e^{c|\operatorname{Im}z|}\)，先乘 \(e^{\delta(z-\theta)^2}\) 抑制增长，得到带有限常数的估计。这一版本将在曲面限制的解析插值证明中使用。

\subsection{Hausdorff–Young 再论}
\label{b2:subsec:170403}

\begin{corollary}[Hausdorff--Young]\label{b2:thm:17-hausdorff-young}
在本卷 \(e^{-2\pi\mathrm ix\cdot\xi}\) 的约定下，\(1\leq p\leq2\) 时
\[
 \|\widehat f\|_{p'}\leq\|f\|_p.
\]
\end{corollary}
\begin{proof}
Fourier 变换从 \(L^1\) 到 \(L^\infty\) 的范数至多一；\chapref{b2:ch:04}的 Plancherel 等式给出 \(L^2\) 等距性。以这两个兼容定义应用 Riesz--Thorin。对 \(1<p\leq2\)，变换先在 \(L^1\cap L^2\cap L^p\) 上定义，再依估计延拓；不宣称任意 \(L^p\) 函数的绝对积分逐点存在。
\end{proof}
插值给出常数一，却没有证明这在中间指数最优。Gaussian 的可计算范数已经提示最优常数问题比有界性更细。逆变换具有相同的两个端点，因而同样得到 \(L^p\) 到 \(L^{p'}\) 的估计。

\subsection{插值思想}
\label{b2:subsec:170404}

两种插值处理的是不同信息。Marcinkiewicz 使用大小分解和分布函数，适用于极大函数这样的次线性算子，代价是通常不精确的常数。Riesz--Thorin 使用线性叠加及复参数，只需强型端点，常数按对数凸性组合。对奇异积分，常用弱 \((1,1)\) 和强 \((2,2)\) 先得到 \(1<p<2\)，再以伴随的相同估计补齐 \(p>2\)。

\ProseParagraphBreak

若一个估计含尺度参数 \(R\)，端点常数写成 \(A_i(R)\)，插值后的幂为 \(A_0(R)^{1-\theta}A_1(R)^\theta\)。例如两个端点分别增长为 \(R^a\) 和 \(R^b\)，中间指数的增长为 \(R^{(1-\theta)a+\theta b}\)。这项尺度核算必须在对 \(R\) 求和以前完成；一个在每个尺度上成立但常数不一致的估计，不能直接推出全算子有界。

% 扩写定位：b2:subsec:170404
为了看清“把输入分配到两个端点”与函数空间之间的联系，可以引入一个直接可计算的量。对一对嵌入同一线性空间的 Banach 空间 \(X_0,X_1\)，以及 \(f\in X_0+X_1\)，定义
\[
 K(t,f;X_0,X_1)\coloneqq
 \inf_{f=f_0+f_1}\bigl(\|f_0\|_{X_0}+t\|f_1\|_{X_1}\bigr),\qquad t>0.
\]
这个下确界允许分解随 \(t\) 改变。仅取固定分解通常丢失端点附近的信息。在 Lebesgue 空间的一对 \(L^1,L^\infty\) 中，它具有完整的重排描述。
\begin{proposition}\label{b2:prop:17-k-functional}
对 \(f\in L^1+L^\infty\)，有
\[
 K(t,f;L^1,L^\infty)=\int_0^t f^*(s)\,\mathrm ds.
\]
若线性算子 \(T:X_0+X_1\to Y_0+Y_1\) 在两端点的范数分别至多 \(A_0,A_1>0\)，则
\[
 K(t,Tf;Y_0,Y_1)\leq A_0K(tA_1/A_0,f;X_0,X_1).
\]
\end{proposition}
\begin{proof}
设 \(f=f_0+f_1\)。任意 \(|E|\leq t\) 都满足
\(\int_E|f|\leq\|f_0\|_1+t\|f_1\|_\infty\)。
在无原子的 Lebesgue 空间，取超水平集合，并在必要时取水平集合的一部分，便有
\(\sup_{|E|\leq t}\int_E|f|=\int_0^t f^*\)。当超水平集合或水平值不能恰好取到最优时，用逼近集合取极限即可。因此右侧不超过每个分解的代价。

反向令 \(a=f^*(t)\)，以保持相位的截断
\[
 f_1(x)=
 \begin{cases}f(x),&|f(x)|\leq a,\\a f(x)/|f(x)|,&|f(x)|>a,\end{cases}
 \qquad f_0=f-f_1.
\]
这里 \(f_1\) 有界，且层蛋糕给出
\(\|f_0\|_1=\int_0^\infty(f^*(s)-a)_+\,\mathrm ds\)。
由重排的单调性，这等于 \(\int_0^t f^*(s)\,\mathrm ds-ta\)，水平集合上的等号不影响积分。端点 \(a=0\) 直接同样处理；若遇到下确界的极限值，先取略大的截断高度再取极限。故该分解的代价就是所求积分。

对于算子估计，每个分解产生 \(Tf=Tf_0+Tf_1\)，代价至多 \(A_0\|f_0\|_{X_0}+tA_1\|f_1\|_{X_1}\)。对全部分解取下确界即可。
\end{proof}
由此可以定义 \(0<\theta<1\)、\(1\leq q\leq\infty\) 时的实插值量
\[
 \|f\|_{\theta,q;K}=
 \left(\int_0^\infty[t^{-\theta}K(t,f)]^q\frac{\mathrm dt}{t}\right)^{1/q},
\]
无穷指数取上确界。对 \(L^1,L^\infty\)，令 \(1/p=1-\theta\)，它与前面 \(L^{p,q}\) 的拟范数等价。下面验证这一说法，而不把它当作一个新的黑箱。

\ProseParagraphBreak

单调性给出 \(K(t,f)\geq t f^*(t)\)，所以 Lorentz 量不超过 \(K\) 量。反向置 \(h(s)=s^{1/p}f^*(s)\)，则
\[
 t^{-1/p'}K(t,f)=
 \int_0^1u^{-1/p}h(tu)\,\mathrm du.
\]
在乘法测度 \(\mathrm dt/t\) 中，变换 \(t\mapsto tu\) 保持范数；Minkowski 因而给出上界
\(\int_0^1u^{-1/p}\,\mathrm du\cdot\|h\|_{L^q(\mathrm dt/t)}
=p'\|f\|_{L^{p,q}}\)。
在 \(q=\infty\) 时直接取上确界也是同一证明。算子的 \(K\) 估计于是同时传递整个高度分布的控制。这里没有推出任意两个弱型端点间的最佳 Lorentz 指数；弱型端点不是已经给定的 Banach 范数端点，其使用还需另行检查。

\ProseParagraphBreak

对两个端点都给定强型的线性算子，代入上式并换变量 \(tA_1/A_0\)，得到插值常数 \(A_0^{1-\theta}A_1^\theta\)。这个乘法结构与 Riesz--Thorin 相同，而证明只使用分解、重排和 Minkowski。前者自然保持同一个中间 Lebesgue 指数，后者自然看到 Lorentz 的第二个指数；两种方法的适用条件和产生的信息因此各有侧重。

\ProseParagraphBreak

几种插值路线在同一个算子上可以互相核对。对可积核卷积，两个同指数端点由 Young 给出，Riesz--Thorin 重现整个强型区间；若核只在临界弱空间中，集合积分与大值分解更合适；若目标是让输入带有一个细于 \(L^p\) 的第二指数，则 \(K\) 泛函直接记录高度分布。选择方法的依据是手中实际掌握的端点，而不是插值定理名字的先后次序。

\ProseParagraphBreak

实际证明还应先检查共同定义域。比如 Fourier 变换同时有 \(L^1\to L^\infty\) 与 \(L^2\to L^2\) 的界，这两种对象在 \(L^1\cap L^2\) 上一致，才可以把它看成同一个待插值算子。若两端的公式只分别定义在互不相容的空间上，即使形式上写出了两个范数界，也没有一个可供插值的共同线性算子。对奇异积分，这个相容性通常由测试函数上的主值公式与平方延拓来建立。

\ProseParagraphBreak

输出是否线性同样重要。极大算子把所有尺度的绝对平均取上确界，一般不是线性的；把输入换成复幂函数后，不能期待极大函数的矩阵元对参数解析。因此本章对它采用分布函数插值。相比之下，固定的 Fourier 变换或固定卷积是线性的，解析变化发生在测试输入与测试输出中，配对仍可验证为全纯函数。假如把“绝对值”“上确界”和“积分”在解析证明中任意交换，便可能把非解析量误放入三线引理。

\ProseParagraphBreak

常数也可以作为一个有效检查。若用两端范数 \(A_0,A_1\) 推导中间界，缩放算子 \(T\mapsto cT\) 后，中间界必须线性乘以 \(|c|\)。Riesz--Thorin 的 \(A_0^{1-\theta}A_1^\theta\) 正好具有这种齐次性；Marcinkiewicz 的未优化分割可能给出较粗的常数，但仍应在适当重选高度阈值后保持齐次。若某个推导出现两个端点常数相乘、总指数大于一，通常意味着阈值没有随算子大小调整，或只是用了一个合法但很粗的非齐次上界。

\ProseParagraphBreak

插值不能恢复被端点假设遗忘的信息。两个范数界可能来自完全不同的机制，插值只利用其共同的线性或次线性结构；它不会证明核的抵消、主值存在性、最优常数或极值函数分类。例如 Hausdorff--Young 的常数一来自基本插值，而 Gaussian 测试揭示中间指数还有更细的最佳常数问题。正向上界与某类测试函数给出的下界应分别记录，只有两者相等才完成最优常数的计算。

\ProseParagraphBreak

最后，开放指数区间中的有界性并不总能通过取极限传到端点。若估计常数在 \(p\to1\) 时发散，对每个固定输入得到的范数序列就没有统一控制；即使常数不发散，也要说明输入空间与输出空间随指数变化时采用了什么收敛。极大函数的具体反例已经给出最直接的检验：它在每个 \(p>1\) 上有界，然而一个单位区间的像有不可积尾部。弱端点的提出正是为了保留这种有限而有意义的控制。

% 本章小结（不计入编号）
\begin{chaptersummary}
弱形式减少了未知函数所需的导数阶数，却没有丢掉方程：紧支集测试辨认内部散度，测试空间和边界类决定所解的边值问题，界面项记录通量平衡。负阶空间把函数源项与散度型源项放在同一连续对偶中。Lax--Milgram 用有界性构造算子，以强制性控制逆，再排除值域的正交补，得到无需对称性的存在与稳定性。

\ProseParagraphBreak

一般椭圆主部只需有界可测系数即可定义能量弱解；低阶项可能使原型失去强制性，但加上足够大的零阶移位可恢复它。去掉移位时，紧嵌入将问题化为恒等算子的紧扰动，可解性由伴随核的有限个条件决定。整个论证首先产生 \(H^1\) 解；要得到 \(H^2\)、连续梯度乃至经典解，下一章还须逐项增加系数、数据和边界的正则性。
\end{chaptersummary}


\BookExercises

各题均为本书自拟，函数空间均取实数域。题中的 \(H^{-1}(\Omega)\) 是以梯度范数赋予 \(H_0^1(\Omega;\mathbb R)\) 后的连续对偶；一维端点与点值取第一卷命题21.4.1的绝对连续代表。前几题核验空间和数据，随后比较近似、系数扰动与共振；最后一题使用选读小节的相容条件。

% 自拟：一维点源的精确对偶范数。
\Needspace{6\baselineskip}
\begin{exercise}\label{b2:ex:15-point-source}
在 \(I=(0,1)\) 上，令 \(H^{-1}(I)=(H_0^1(I;\mathbb R))^*\)，对偶范数对应 \(\|v'\|_2\)。对 \(0<s<1\)，以绝对连续代表定义 \(\langle\delta_s,v\rangle=v(s)\)。求 \(\delta_s\) 的精确对偶范数，并说明哪个单位范数函数达到对偶上确界。
\end{exercise}
\begin{solution}
令 \(G_s(x)=(1-s)x\)（\(x\le s\)）、\(G_s(x)=s(1-x)\)（\(x\ge s\)）。对 \(v\in H_0^1(0,1)\) 的绝对连续代表，零端点值给出
\[
 \int_0^1G_s'v'=(1-s)\int_0^sv'-s\int_s^1v'
 =(1-s)v(s)+sv(s)=v(s).
\]
因此 \(|v(s)|\le\lVert G_s'\rVert_2\lVert v'\rVert_2\)，且取 \(v=G_s/\lVert G_s'\rVert_2\) 达到上界。直接计算 \(\lVert G_s'\rVert_2^2=s(1-s)^2+(1-s)s^2=s(1-s)\)，所以范数为 \(\sqrt{s(1-s)}\)。
\end{solution}

% 自拟：点质量维数结论已移入正文；改求平均点源的精确能量。
\Needspace{6\baselineskip}
\begin{exercise}\label{b2:ex:15-high-dimensional-point}
设整数 \(d\ge2\)、\(0<\varepsilon<1\)，\(B=B(0,1)\subset\mathbb R^d\)，记 \(s_d=|S^{d-1}|\)。令 \(H^{-1}(B)=(H_0^1(B;\mathbb R))^*\)，对偶范数对应 \(\|\nabla v\|_2\)。在单位球上取单位质量源
\(f_\varepsilon=\frac{d}{s_d\varepsilon^d}\cdot\mathbf1_{B(0,\varepsilon)}\)。求 \(-\Delta u_\varepsilon=f_\varepsilon\)、零 Dirichlet 边界的径向弱解，并证明
\[
 \|f_\varepsilon\|_{H^{-1}}^2
 =\frac1{s_d}\left(\frac{\varepsilon^{2-d}}{d+2}
                  +\int_\varepsilon^1r^{1-d}\,\mathrm dr\right).
\]
分别写出二维与高维的发散速度，解释它与 \(f_\varepsilon\to\delta_0\) 的分布收敛是否矛盾。
\end{exercise}
\begin{solution}
径向方程为 \(-(r^{d-1}u')'=r^{d-1}f_\varepsilon\)。在零点不留附加点源，先从零积分，得到 \(u'=-r/(s_d\varepsilon^d)\) 于内球，\(u'=-1/(s_dr^{d-1})\) 于外环。令
\[
 u_\varepsilon(r)=
 \begin{cases}
 \displaystyle\frac1{s_d}\int_\varepsilon^1t^{1-d}\,\mathrm dt
            +\frac{\varepsilon^2-r^2}{2s_d\varepsilon^d},&r\leq\varepsilon,\\[3pt]
 \displaystyle\frac1{s_d}\int_r^1t^{1-d}\,\mathrm dt,&\varepsilon\leq r\leq1.
 \end{cases}
\]
函数及径向通量在分界面相接，零点附近为光滑二次函数，边界值为零；分区积分便验证弱方程。由能量唯一性，它就是所求弱解。等式 \(\langle f_\varepsilon,v\rangle=\int\nabla u_\varepsilon\cdot\nabla v\) 与 Cauchy--Schwarz 给对偶范数上界，取 \(v=u_\varepsilon/\|\nabla u_\varepsilon\|_2\) 达到上界。径向积分故给
\[
 \|f_\varepsilon\|_{H^{-1}}^2
 =s_d\int_0^1|u_\varepsilon'(r)|^2r^{d-1}\,\mathrm dr
 =\frac1{s_d}\left(\frac{\varepsilon^{2-d}}{d+2}
                  +\int_\varepsilon^1r^{1-d}\,\mathrm dr\right).
\]
二维结果为 \((2\pi)^{-1}(1/4+\log(1/\varepsilon))\)；当 \(d>2\) 时为
\(s_d^{-1}[\varepsilon^{2-d}/(d+2)+(\varepsilon^{2-d}-1)/(d-2)]\)，范数分别按 \(\sqrt{\log(1/\varepsilon)}\) 与 \(\varepsilon^{(2-d)/2}\) 发散。另一方面，对固定光滑测试函数，球内平均趋于其中心值，所以源项在分布意义下趋于点质量。这种弱得多的收敛不要求能量对偶范数有界，因而没有矛盾。
\end{solution}

% 自拟：分段材料的显式有效导热。
\Needspace{6\baselineskip}
\begin{exercise}\label{b2:ex:15-layered-material}
令 \(a=1\) 于 \((0,1/2)\)、\(a=4\) 于 \((1/2,1)\)。在实空间 \(H^1(0,1)\) 中求 \(-(au')'=0\)、\(u(0)=0,u(1)=1\) 的弱解（端点取绝对连续代表），计算左右斜率及能量 \(E(u)=\frac12\int_0^1a|u'|^2\)，并判断它是否属于 \(H^2(0,1)\)。
\end{exercise}
\begin{solution}
通量为常数 \(q\)，且 \(1=\int_0^1u'=q(1/2+1/8)=5q/8\)，所以 \(q=8/5\)。左斜率为 \(8/5\)，右斜率为 \(2/5\)，解是通过 \((0,0),(1/2,4/5),(1,1)\) 的分段仿射函数。能量 \(\tfrac12\int a|u'|^2=\tfrac12q\int u'=4/5\)。一阶导数在中点有非零跳跃，二阶分布导数含点质量，故不属于 \(L^2\)，解不在 \(H^2\)。
\end{solution}

% 自拟：Céa引理已在正文证明；改用二维模型检验误差常数的最优性。
\Needspace{6\baselineskip}
\begin{exercise}\label{b2:ex:15-cea}
设 \(H=\mathbb R^2\) 取欧氏内积，\(\beta\in\mathbb R\)，并令
\[
 \begin{gathered}
 B_\beta=\begin{pmatrix}1&-\beta\\\beta&1\end{pmatrix},\qquad
 a_\beta(x,y)=(B_\beta x)\cdot y,\\
 \ell_\beta(y)=(-\beta,1)\cdot y,\qquad
 V=\operatorname{span}\{(1,0)\}.
 \end{gathered}
\]
求使 \(|a_\beta(x,y)|\le M\cdot|x|\cdot|y|\)、\(a_\beta(x,x)\ge\alpha|x|^2\) 对全部 \(x,y\) 成立的最小 \(M\) 和最大 \(\alpha\)。求满足 \(a_\beta(u,y)=\ell_\beta(y)\) 对全部 \(y\in H\) 成立的解 \(u\)，以及满足 \(a_\beta(u_V,y)=\ell_\beta(y)\) 对全部 \(y\in V\) 成立的解 \(u_V\in V\)，并计算
\[
 \frac{|u-u_V|}{\inf_{v\in V}|u-v|}.
\]
据此说明\cref{b2:prop:15-cea}中的常数 \(M/\alpha\) 对一般非对称强制型不能统一减小；当 \(\beta\ne0\) 时，指出 Galerkin 解与欧氏正交投影的区别。
\end{exercise}
\begin{solution}
由 \(B_\beta^{\mathsf T}B_\beta=(1+\beta^2)I\)，有 \(|B_\beta x|=\sqrt{1+\beta^2}|x|\)，故最小有界常数为 \(M=\sqrt{1+\beta^2}\)；取单位 \(y=B_\beta x/|B_\beta x|\) 可达到算子范数。又 \(a_\beta(x,x)=|x|^2\)，故最大强制常数为 \(\alpha=1\)。

全空间方程等价于 \(B_\beta u=(-\beta,1)\)，所以 \(u=(0,1)\)。写 \(u_V=(s,0)\) 并用 \((1,0)\) 测试，得到 \(s=-\beta\)。因此
\[
 e=u-u_V=(\beta,1),\qquad |e|=\sqrt{1+\beta^2},\qquad
 \inf_{v\in V}|u-v|=1,
\]
最后的下确界在欧氏投影 \(v=(0,0)\) 达到。误差比正好等于 \(M/\alpha\)，所以一般 Céa 估计中的该常数确实可以达到。另一方面，\(B_\beta e=(0,1+\beta^2)\)，故 \(a_\beta(e,w)=0\) 对全部 \(w\in V\) 成立，但 \(e\cdot(1,0)=\beta\)。当 \(\beta\ne0\) 时，这一测试正交性不是欧氏正交性，\(u_V\) 也不是欧氏最佳逼近；非对称型的准最佳逼近不能改写成 Poisson 梯度内积下的等距投影结论。
\end{solution}

% 自拟：系数在L无穷范数下的稳定性。
\Needspace{6\baselineskip}
\begin{exercise}\label{b2:ex:15-coefficient-stability}
设整数 \(d\ge1\)，\(\Omega\subset\mathbb R^d\) 为非空有界开集，\(A_n,A\in L^\infty(\Omega;\mathbb R^{d\times d})\) 满足共同椭圆下界 \(\xi^{\mathsf T}A_n(x)\xi,\xi^{\mathsf T}A(x)\xi\ge\lambda|\xi|^2\)（\(\lambda>0\)），对几乎处处 \(x\) 和全部 \(\xi\in\mathbb R^d\) 成立，且 \(\lVert A_n-A\rVert_\infty\to0\)，这里矩阵取欧氏算子范数。令 \(f\in H^{-1}(\Omega)=(H_0^1(\Omega;\mathbb R))^*\)，其对偶范数对应梯度范数；\(u_n,u\in H_0^1(\Omega;\mathbb R)\) 分别弱解 \(-\operatorname{div}(A_n\nabla u_n)=f\) 与 \(-\operatorname{div}(A\nabla u)=f\)。证明 \(u_n\to u\) 于 \(H_0^1\)，并给出定量估计。
\end{exercise}
\begin{solution}
以 \(A_n\) 为左端，两方程相减得 \(\int A_n\nabla(u_n-u)\cdot\nabla v=\int(A-A_n)\nabla u\cdot\nabla v\)。取 \(v=u_n-u\)，利用椭圆性与 Hölder 不等式，得到
\[\lVert\nabla(u_n-u)\rVert_2\le\lambda^{-1}\lVert A_n-A\rVert_\infty\lVert\nabla u\rVert_2
 \le\lambda^{-2}\lVert A_n-A\rVert_\infty\lVert f\rVert_{H^{-1}}.\]
右端趋零，Poincaré 不等式再给出完整 \(H^1\) 收敛。
\end{solution}

% 自拟：Neumann归一化与相容性。
\Needspace{6\baselineskip}
\begin{exercise}\label{b2:ex:15-neumann}
令 \(\Omega=(0,1)\cup(2,3)\)，求解类为 \(H^1(\Omega;\mathbb R)\)，纯 Neumann 通量为零，即 \(\int_\Omega u'v'=\int_\Omega fv\) 对每个 \(v\in H^1(\Omega;\mathbb R)\) 成立。证明源项在两个分支分别等于 \(1,-1\) 时，虽然总积分为零，仍无弱解。若源项改为 \(f(x)=2x-1\) 于第一分支、\(f(x)=2(x-2)-1\) 于第二分支，写出全部解，并说明只要求全域零均值为何仍不唯一。
\end{exercise}
\begin{solution}
第一种数据用第一分支的示性函数测试；它在不连通开集上属于 \(H^1\)，梯度为零，右端积分为 \(1\)，所以矛盾。第二种数据在每个分支积分为零。以局部坐标 \(s\in(0,1)\) 积分 \(-u''=2s-1\) 得 \(u'=s-s^2\)，满足两端零通量；再积分得 \(u=s^2/2-s^3/3+c_j\)，\(j=1,2\)。局部多项式均值为 \(1/12\)，全域零均值只给 \(c_1+c_2=-1/6\)，留下一个自由参数。分别要求两个分支零均值才确定 \(c_1=c_2=-1/12\)。这与正文的连通假设相吻合。
\end{solution}

% 自拟：漂移项的反对称结构。
\Needspace{6\baselineskip}
\begin{exercise}\label{b2:ex:15-divergence-free-drift}
设整数 \(d\ge1\)，\(\Omega\subset\mathbb R^d\) 为非空有界开集，\(b\in W^{1,\infty}(\Omega;\mathbb R^d)\) 且 \(\operatorname{div}b=0\) 几乎处处，\(A\in L^\infty(\Omega;\mathbb R^{d\times d})\)，并对几乎处处 \(x\)、全部 \(\xi\in\mathbb R^d\) 有 \(\xi^{\mathsf T}A(x)\xi\ge\lambda|\xi|^2\)，其中 \(\lambda>0\)。令 \(H^{-1}(\Omega)=(H_0^1(\Omega;\mathbb R))^*\) 的范数对应梯度范数。证明零边界问题 \(-\operatorname{div}(A\nabla u)+b\cdot\nabla u=f\) 对任意 \(f\in H^{-1}(\Omega)\) 在 \(H_0^1(\Omega;\mathbb R)\) 中唯一弱可解，其能量估计不含 \(\lVert b\rVert_\infty\)。
\end{exercise}
\begin{solution}
对紧支集光滑 \(u\)，有 \(\int(b\cdot\nabla u)u=\tfrac12\int b\cdot\nabla(u^2)=0\)。双线性连续性和 \(H_0^1\) 密度把该等式推广到全部 \(H_0^1\)。所以型虽可能不对称，却满足 \(a(u,u)\ge\lambda\lVert\nabla u\rVert_2^2\)。有界性仍含 \(\lVert b\rVert_\infty\)，但强制常数不含它。Lax--Milgram 给出 \(\lVert\nabla u\rVert_2\le\lambda^{-1}\lVert f\rVert_{H^{-1}}\) 及唯一性。
\end{solution}

% 自拟：不同共振频率的相容条件与全体解。
\Needspace{6\baselineskip}
\begin{optionalexercise}\label{b2:ex:15-resonance}
给定 \(q\in\mathbb R\)，在 \(H_0^1(0,\pi;\mathbb R)\) 中求解零端点问题 \(-u''-4u=\sin x+q\sin2x\)。这里弱形式是 \(\int_0^\pi(u'v'-4uv)=\int_0^\pi(\sin x+q\sin2x)v\) 对全部 \(v\in H_0^1(0,\pi;\mathbb R)\) 成立。确定哪些 \(q\) 允许弱解，并在可解时写出全部解。
\end{optionalexercise}
\begin{solution}
型对称，齐次核是 \(\operatorname{span}\{\sin2x\}\)。与核配对得相容条件 \(q\int_0^\pi\sin^22x=0\)，即 \(q=0\)。此时特解为 \(-\tfrac13\sin x\)，故全部解为 \(-\tfrac13\sin x+c\sin2x\)。核的刻画可由一维分布方程先推出 \(u''=-4u\in L^2\)，再由常系数常微分方程及端点条件得到，亦可直接使用正弦展开。因此没有遗漏较粗糙的能量解。
\end{solution}

% !TeX root = ../../Book2.tex
\chapter{奇异积分与 Calderón–Zygmund 理论}
\label{b2:ch:18}
\BookChapterNumber{b2:ch:18}
\begin{chapterabstract}
若卷积核在原点附近像距离的负维数次幂那样增长，取绝对值后的积分通常发散，Young 不等式便不能直接使用。然而，核的正负部分可能相互抵消，使对称截断仍有极限。本章从 Hilbert 变换的具体计算出发，说明怎样把这种抵消转化为范数估计：先将输入分成有界的好部与零平均的坏块，再把坏块的零平均移到核的差分上。由此得到弱端点、中间指数的有界性与主值实现，最后用于共轭调和函数和椭圆方程的二阶导数估计。
\end{chapterabstract}

\begin{learninggoals}\begin{itemize}
\item\label{b2:item:18-goal-1}计算 Hilbert 变换的主值与 Fourier 符号，区分平方范数有界和点态极限。
\item\label{b2:item:18-goal-2}构造 Calderón--Zygmund 分解，说明坏块的零平均怎样改善远场衰减。
\item\label{b2:item:18-goal-3}核对标准核的大小、光滑性及离支集表示，并辨认平方有界性这一独立假设。
\item\label{b2:item:18-goal-4}完成弱型、插值、伴随与极大截断的论证，把稠密类上的主值传给一般输入。
\item\label{b2:item:18-goal-5}用 Riesz 变换描述 Poisson 共轭结构、频率投影与椭圆二阶估计。
\end{itemize}\end{learninggoals}

对区间示性函数作 Hilbert 变换，答案是一个对数。输入有界且可积，输出却在区间端点无界，并在无穷远只按倒数衰减。另一方面，它的 Fourier 符号模长为一，平方范数完全保持。同一个算子为什么能够同时具有这些性质？要回答这个问题，必须分清两个作用：绝对值估计控制大小，正负抵消则能消去绝对值估计中最危险的部分。

\ProseParagraphBreak

\chapref{b2:ch:17}已经准备了三件工具。二进微分定理让局部平均反映函数值，极大立方体分解把大平均的位置收集起来，Marcinkiewicz 插值则把弱端点与平方有界性接成中间的强型估计。本章调用它们时仍会写出所需版本及输入条件。Fourier 部分沿用\secref{b2:sec:0401}的 Plancherel 延拓，归一化固定为 \(e^{-2\pi\mathrm i x\cdot\xi}\)；涉及分布时，只使用\chapref{b2:ch:05}中的测试函数配对、分布导数和 Fourier 变换。除\subsecref{b2:subsec:180502}的周期比较外，底空间均为带 Lebesgue 测度的 \(\mathbb R^d\)。

\ProseParagraphBreak

首次阅读可先沿 Hilbert 变换、好坏分解、标准核、弱型证明这一顺序走到\subsecref{b2:subsec:180403}。这样先看清一个具体核怎样进入一般定理。随后读极大截断与主值传递，再到\secref{b2:sec:1805}看 Riesz 变换的实现和应用。向量值估计与 BMO 端点解释的是同一方法怎样服务于后续章节，可以在掌握标量证明后回读；它们不作为本章标量 \(L^p\) 有界性的前提。

\ProseParagraphBreak

进一步阅读可参考 Elias M. Stein 的《Singular Integrals and Differentiability Properties of Functions》\cite{Stein1970Singular}与 Loukas Grafakos 的《Classical Fourier Analysis》\cite{Grafakos2014Classical}。前者适合继续理解核抵消与微分性质的联系，后者可用于扩展奇异积分、乘子与函数空间的背景。本章采用的标准核版本在正文中证明，较弱核条件下的进一步结论另行说明出处和适用范围。

% !TeX root = ../../Book2.tex
\section{Hilbert 变换}
\label{b2:sec:1801}

在实直线上，核 \(1/x\) 的绝对积分在原点和无穷远都发散。然而它是奇函数，原点两侧的对称截断可以相互抵消。Hilbert 变换把这种抵消写成一个 Fourier 乘子，是研究奇异积分的第一个完整模型。

\subsection{主值}
\label{b2:subsec:180101}

对 \(f\in\mathcal S(\mathbb R)\)，定义\term[hilbert-bianhuan]{Hilbert 变换}[Hilbert transform]
\[
 Hf(x)\coloneqq\lim_{\varepsilon\downarrow0}
 \dfrac1\pi\int_{|x-y|>\varepsilon}\dfrac{f(y)}{x-y}\,\mathrm dy.
\]
\symidx{\(Hf\)：Hilbert 变换}
这里的极限称为\term[zhuzhi-jifen]{主值积分}[principal value integral]。换元 \(t=x-y\) 后，近原点的积分可以写成
\[
 \int_{\varepsilon<|t|<1}\dfrac{f(x-t)-f(x)}{t}\,\mathrm dt.
\]
均值定理使差商有界，所以极限存在；远处因 \(f\) 快速衰减而绝对可积。减去 \(f(x)\) 只在对称区间上进行，不能对一个无穷的常数积分作没有解释的减法。

\ProseParagraphBreak

若 \(f=\mathbf1_{(0,1)}\)，避开两个端点可直接计算
\[
 Hf(x)=\dfrac1\pi\log\left|\dfrac{x}{x-1}\right|.
\]
在区间内部，原点附近的对称抵消仍给出同一公式；在两个端点主值发散。这个例子同时显示主值不是绝对积分，以及有界输入可以产生无界输出。函数的单个点值改变不影响积分，但主值是否存在是对观察点而言的条件。

\subsection{Fourier 乘子}
\label{b2:subsec:180102}

考虑有限截断核 \(h_{\varepsilon,R}(t)=\mathbf1_{\{\varepsilon<|t|<R\}}/(\pi t)\)。其 Fourier 变换为
\[
 \widehat h_{\varepsilon,R}(\xi)
 =-\dfrac{2\mathrm i}{\pi}\int_\varepsilon^R\dfrac{\sin(2\pi\xi t)}t\,\mathrm dt.
\]
变换的符号由本卷负号指数决定。正、负频率得到相反的虚数值，不存在额外的 \(2\pi\) 幅值因子。
\begin{lemma}\label{b2:thm:18-dirichlet-integral}
上述乘子对 \(\varepsilon,R,\xi\) 一致有界，且在 \(\xi\ne0\) 时当 \(\varepsilon\downarrow0\)、\(R\uparrow\infty\) 收敛到 \(-\mathrm i\operatorname{sgn}\xi\)。
\end{lemma}
\begin{proof}
换元把积分化为 \(\int_a^b\sin u/u\,\mathrm du\)。在 \((0,1)\) 上被积函数的模不超过一；在 \((1,\infty)\) 上分部积分，使任意尾段的积分模不超过 \(3/a\)。这同时给出一致有界性和尾积分收敛。为计算全积分，先加入 \(e^{-\delta u}\)，对频率参数 \(s\) 求导得到
\[
 \dfrac{\mathrm d}{\mathrm ds}\int_0^\infty
 e^{-\delta u}\dfrac{\sin(su)}u\,\mathrm du
 =\int_0^\infty e^{-\delta u}\cos(su)\,\mathrm du
 =\dfrac{\delta}{\delta^2+s^2}.
\]
零参数处积分为零，故在 \(s=1\) 为 \(\arctan(1/\delta)\)。分部积分同样控制有阻尼和无阻尼的远尾，先固定尾端再令 \(\delta\downarrow0\)，得到 \(\pi/2\)。负频率由奇性得到负号。
\end{proof}
乘子在零点的取值不影响 \(L^2\) 算子；这与核在原点的截断约定是两个不同层次的无关零测修改。

\subsection{\texorpdfstring{\(L^2\)}{L2} 有界性}
\label{b2:subsec:180103}

\begin{theorem}\label{b2:thm:18-hilbert-l2}
Hilbert 变换具有唯一的 \(L^2(\mathbb R)\) 延拓，满足
\[
 \widehat{Hf}(\xi)=-\mathrm i\operatorname{sgn}(\xi)\widehat f(\xi),
 \qquad \|Hf\|_2=\|f\|_2.
\]
有限截断 \(h_{\varepsilon,R}*f\) 在 \(L^2\) 中趋于 \(Hf\)。
\end{theorem}
\begin{proof}
由 Plancherel 和上一引理，乘子的绝对值有共同上界，且几乎处处趋于绝对值为一的极限。对 \( |\widehat f|^2\) 使用支配收敛，得到截断在 \(L^2\) 中收敛及等距性。对 Schwartz 函数，主值积分已逐点存在，\(L^2\) 收敛抽出的几乎处处收敛子列识别两种定义。Schwartz 稠密性给出唯一延拓。
\end{proof}
因此 \(H^2f=-f\) 在 \(L^2\) 中成立，且 Hilbert 空间伴随为 \(-H\)。它不是正算子：不交区间之间的核可能为正也可能为负；其有界性来自 Fourier 符号和空间抵消，不能用正核的积分大小证明。

\subsection{\texorpdfstring{\(L^p\)}{Lp} 有界性}
\label{b2:subsec:180104}

核 \(K(x,y)=1/(\pi(x-y))\) 满足
\[
 |K(x,y)|\leq C|x-y|^{-1},\qquad
 |K(x,y)-K(x,y')|\leq C\dfrac{|y-y'|}{|x-y|^2}
\]
当 \(|y-y'|\leq|x-y|/2\) 时成立。第二项由差分分母直接计算。对于紧支集输入，在支集外不需要主值，\(Hf\) 等于普通核积分。这三项事实——\(L^2\) 有界、远离对角线的核表达、核差分控制——是稍后一般定理的全部输入。

\ProseParagraphBreak

由\cref{b2:thm:18-cz-lp}得到 \(1<p<\infty\) 的 \(L^p\) 有界性。这里是后文同章定理的应用位置，证明将在\secref{b2:sec:1804}完成。对任意 \(L^p\) 输入，连续延拓首先定义等价类；主值表示则需要控制截断的最大值，不能从强范数有界性独自推出。对本模型，这项控制在该节末尾给出。

\ProseParagraphBreak

无穷端点不能补上：区间示性函数的对数奇性已经否定 \(L^\infty\to L^\infty\)。在 \(L^1\) 端点，其远处渐近为 \(1/(\pi x)\)，也不绝对可积。后面的弱 \((1,1)\) 结论允许这样的尾部。

% 扩写定位：b2:subsec:180104
Hilbert 变换的两个端点都可以用一个函数看出障碍。对于 \(\mathbf1_{(0,1)}\)，已计算的主值是
\(\pi^{-1}\log|x/(x-1)|\)。在 \(x=0,1\) 附近，它具有对数无界性，故 Hilbert 变换不能把全部 \(L^\infty\) 输入送入 \(L^\infty\)。在 \(|x|\to\infty\) 时，展开 \(\log(1+1/(x-1))\) 得到 \(H\mathbf1_{(0,1)}(x)=1/(\pi x)+O(x^{-2})\)，所以其绝对值不可积，否定强型 \((1,1)\)。

\ProseParagraphBreak

两个失败来自不同的区域：无穷端点失败出现在输入跳跃的边缘，积分端点失败出现在远场。函数本身既可积又有界，不能通过要求输入同时属于两个端点来消去这些障碍。弱型 \((1,1)\) 仍然允许 \(1/x\) 的尾部，而 \(1<p<\infty\) 时对数奇性和 \(1/x\) 尾部都可以积到 \(p\) 次；这解释了强型区间的几何位置。

\ProseParagraphBreak

若输入同时具有零积分，远场的第一项会消失。例如取 \(f=\mathbf1_{(0,1)}-\mathbf1_{(1,2)}\)，则 \(Hf(x)=O(x^{-2})\) 于无穷处。一般紧支集 \(f\) 且 \(\int f=0\) 时，
\[
 Hf(x)=\frac1\pi\int f(y)
 \left(\frac1{x-y}-\frac1x\right)\mathrm dy,
\]
当 \(|x|\) 大于支集半径的两倍，绝对值不超过 \(C|x|^{-2}\int|y f(y)|\,\mathrm dy\)。这就是坏块尾估计的一维具体形式。不过零积分没有消除区间端点的对数奇性，仍不能恢复 \(L^\infty\) 强型。为讨论这类局部平均振荡，需要\appref{b2:app:C}的 BMO 空间。

% !TeX root = ../../Book2.tex
\section{Calderón–Zygmund 分解}
\label{b2:sec:1802}

核的奇性不能靠输入处处小来控制，因为一个 \(L^1\) 函数可能具有任意高的尖峰。Calderón–Zygmund 分解把尖峰限制在可计量的坏区域，并使每块坏函数的总质量为零。

\subsection{好/坏分解}
\label{b2:subsec:180201}

下面采用\term[calderon-zygmund-fenjie]{Calderón--Zygmund 分解}[Calderón--Zygmund decomposition]的下述版本。

\begin{theorem}[Calderón--Zygmund 分解]\label{b2:thm:18-cz-decomposition}
设 \(f\in L^1(\mathbb R^d)\)、\(\alpha>0\)。存在互不相交的二进立方体 \(Q_j\) 以及 \(f=g+\sum_jb_j\)，使
\[
 \sum_j|Q_j|\leq\alpha^{-1}\|f\|_1,\quad
 \|g\|_\infty\leq2^d\alpha,\quad \|g\|_1\leq\|f\|_1,
\]
\[
 b_j=(f-f_{Q_j})\mathbf1_{Q_j},\quad
 \int b_j=0,\quad
 \sum_j\|b_j\|_1\leq2\|f\|_1,
\]
其中 \(f_Q=|Q|^{-1}\int_Q f\)，级数在 \(L^1\) 中收敛。
\end{theorem}
\begin{proof}
对 \(\{M_{\mathcal D}f>\alpha\}\) 取\cref{b2:thm:17-dyadic-level}的极大立方体。定义 \(g=f\) 在其并的补集上，\(g=f_{Q_j}\) 在每个 \(Q_j\) 内。立方体平均上界和二进微分结论分别控制内、外的本性上确界。\(\int_{Q_j}|g|\leq\int_{Q_j}|f|\)，所以 \(L^1\) 界成立。坏块积分为零，且 \(\|b_j\|_1\leq2\int_{Q_j}|f|\)。不交性给出范数和与收敛，逐点除坏区域外全部项为零，坏区域内至多一项非零，故分解等式成立。
\end{proof}
复值 \(f\) 的平均仍是复数，控制来自 \(|f_Q|\leq(|f|)_Q\)，不需要把正部负部分别分解。如果原函数还属于 \(L^2\)，则 \(g\in L^2\)，而坏块级数也在 \(L^2\) 中收敛：其尾部平方积分等于不交支集上平方积分的尾和。

\ProseParagraphBreak

好函数的名称只表示它有受控的高度，并不意味着连续或光滑。坏函数虽保留尖峰，却拥有零平均；这个条件会在远场去掉核的常数项。

\subsection{弱型 \texorpdfstring{\((1,1)\)}{(1,1)}}
\label{b2:subsec:180202}

\begin{lemma}\label{b2:thm:18-cancellation-tail}
设 \(b\in L^1\) 支撑于边长 \(\ell\)、中心 \(c\) 的立方体 \(Q\)，且 \(\int b=0\)。若核满足
\[
 |K(x,y)-K(x,c)|\leq C\dfrac{|y-c|^\delta}{|x-c|^{d+\delta}}
 \quad(|x-c|\geq2|y-c|)
\]
其中 \(\delta>0\)，则在 \(Q^*=8\sqrt d\,Q\) 外有
\[
 \int_{\mathbb R^d\setminus Q^*}
 \left|\int_QK(x,y)b(y)\,\mathrm dy\right|\,\mathrm dx
 \leq C_{d,\delta}\|b\|_1.
\]
\end{lemma}
\begin{proof}
零平均使积分等于 \(\int_Q(K(x,y)-K(x,c))b(y)\,\mathrm dy\)。在扩大立方体外，\(|x-c|\geq4\sqrt d\,\ell\)，而 \(|y-c|\leq\sqrt d\,\ell/2\)，故差分假设适用。先对 \(y\) 取绝对积分，再对 \(x\) 用极坐标，尾核积分为
\[
 C\ell^\delta\int_{|x-c|\geq4\sqrt d\,\ell}|x-c|^{-d-\delta}\,\mathrm dx
 \leq C_{d,\delta}.
\]
Tonelli 完成积分交换。
\end{proof}
若不减去 \(K(x,c)\)，大小条件只有 \(|x-c|^{-d}\)，其远场绝对积分发散。这说明分解中的零平均不是方便的附属性质，而是端点估计中不可替代的衰减来源。扩大倍数的具体值只为了保证差分假设，后续常数允许依赖维数。

\subsection{插值}
\label{b2:subsec:180203}

将好部分交给 \(L^2\) 理论、坏部分交给上段尾估计，就能得到弱 \((1,1)\)。设暂时 \(f\in L^1\cap L^2\)，\(T\) 在 \(L^2\) 上有界，取 \(\alpha=t\)。好部分满足
\[
 |\{|Tg|>t/2\}|\leq4t^{-2}\|T\|_{2\to2}^2\|g\|_2^2
 \leq C t^{-1}\|f\|_1,
\]
因为 \(\|g\|_2^2\leq\|g\|_\infty\cdot\|g\|_1\)。坏立方体扩大后的并也有不超过 \(C t^{-1}\|f\|_1\) 的测度。在其补集上，核差分给出 \(\int|T\sum_jb_j|\leq C\sum_j\|b_j\|_1\)。对这项使用 \(L^1\) Chebyshev 即得同一界。

\ProseParagraphBreak

若坏块有无穷多个，先对有限和应用尾估计，再使用 \(L^2\) 收敛抽取几乎处处收敛子列和 Fatou。不能假定算子输出的无穷级数先在 \(L^1\) 中收敛；结论只在坏区域的补集上有该积分控制。下一节把核与算子的准确假设固定下来。

% !TeX root = ../../Book2.tex
\section{Calderón–Zygmund 核}
\label{b2:sec:1803}

Hilbert 核与 Riesz 核都有两种不同的性质：核本身在对角线附近不可积，移开一个变量后得到的核之差却衰减得更快。前一种性质决定积分为什么需要截断，后一种性质使坏块的零平均能够发挥作用。本节把这些条件分开写出，再说明怎样把核与一个已经有平方估计的算子联系起来。

\subsection{大小条件}
\label{b2:subsec:180301}

设整数 \(d\geq1\)，\(K:(\mathbb R^d\times\mathbb R^d)\setminus\{x=y\}\to\mathbb C\) 可测。若存在有限常数 \(C_0\geq0\)，使对每个 \(x\ne y\) 都有
\[
 |K(x,y)|\leq C_0|x-y|^{-d},
\]
则称 \(K\) 满足标准大小条件。半径为 \(r\) 的环带有 \(r^d\) 量级的体积，所以这个上界在每一倍增环带上只给出常数量级的绝对积分；对角线附近有无穷多个这样的环带，不能直接求和。

\ProseParagraphBreak

若 \(f\in L^2\) 支撑在紧集 \(F\) 内，则 \(f\in L^1\)。在与 \(F\) 有正距离的点 \(x\) 处，大小条件给出
\[
 \int_F|K(x,y)|\cdot|f(y)|\,\mathrm dy
 \leq C_0\cdot\operatorname{dist}(x,F)^{-d}\cdot\|f\|_1<\infty.
\]
因此支集外的核积分具有普通绝对积分的意义。这里同时用了输入的有限支集和观察点到支集的距离；支集内部的积分尚未由此定义。

\ProseParagraphBreak

卷积型核写成 \(K(x,y)=k(x-y)\)。例如取有界球面函数 \(\Omega\)，令 \(k(z)=\Omega(z/|z|)\cdot|z|^{-d}\)。若光滑输入在 \(x\) 附近恒等于一，则局部截断积分含有
\[
 \int_{\epsilon<|z|<1}k(z)\,\mathrm dz
 =\left(\int_{S^{d-1}}\Omega\,\mathrm d\sigma\right)\cdot
    \log\frac1\epsilon.
\]
球面平均不为零时，这一项发散。齐次次数 \(-d\) 给出大小条件，却没有提供抵消；球面上的奇函数平均为零，Hilbert 核就是一维的例子。

\subsection{光滑性条件}
\label{b2:subsec:180302}

坏块的零平均使 \(K(x,y)\) 被替换为 \(K(x,y)-K(x,c_Q)\)。因此需要控制的是移动输入变量 \(y\) 后的变化。对上述离开对角线的核，固定 \(0<\delta\leq1\)。标准光滑条件要求存在有限常数 \(C_1\geq0\)，使
\[
 \begin{aligned}
 |K(x,y)-K(x,y')|&\leq C_1\frac{|y-y'|^\delta}{|x-y|^{d+\delta}}
 &&\left(|y-y'|\leq\frac{|x-y|}2\right),\\
 |K(x,y)-K(x',y)|&\leq C_1\frac{|x-x'|^\delta}{|x-y|^{d+\delta}}
 &&\left(|x-x'|\leq\frac{|x-y|}2\right).
 \end{aligned}
\]
所有出现的原始点均满足 \(x\ne y\)。距离限制使新点仍离开对角线，且新旧距离可互相比较。第一行控制原算子的坏块；取伴随后输入、输出互换，第二行才成为伴随坏块所需的条件。

\ProseParagraphBreak

核的梯度界是一个常用的检查方法。若 \(K\) 在对角线外为 \(C^1\)，且
\[
 |\nabla_xK(x,y)|+|\nabla_yK(x,y)|\leq C|x-y|^{-d-1},
\]
则沿连接新旧点的线段积分导数即可取 \(\delta=1\)。例如 \(|y-y'|\leq|x-y|/2\) 时，线段上每点到 \(x\) 的距离至少为 \(|x-y|/2\)，所以积分中的梯度有同一个上界。

\ProseParagraphBreak

弱型证明实际上只用到了输入变量差的积分控制。设 \(h=|y-y'|>0\)，上述 Hölder 界给出
\[
 \begin{aligned}
 \int_{|x-y|>2h}|K(x,y)-K(x,y')|\,\mathrm dx
 &\leq C_1|S^{d-1}|h^\delta
            \int_{2h}^\infty r^{-1-\delta}\,\mathrm dr\\
 &=\frac{C_1|S^{d-1}|}{\delta\,2^\delta}.
 \end{aligned}
\]
\(h=0\) 时差为零。因而标准光滑性蕴含积分 Hörmander 条件
\[
 \sup_{y,y'}\int_{|x-y|>2|y-y'|}
       |K(x,y)-K(x,y')|\,\mathrm dx<\infty.
\]
它直接对应上一节对坏块远场的积分；点态条件还将用于截断比较和平均振荡估计。

\ProseParagraphBreak

也可以把幂次换成较慢的连续模量。设 \(\omega:[0,1]\to[0,\infty)\) 单调不减、\(\omega(0)=0\)，且
\[
 |K(x,y)-K(x,y')|\leq C|x-y|^{-d}\cdot
 \omega\!\left(\frac{|y-y'|}{|x-y|}\right)
 \quad\left(|y-y'|\leq\frac{|x-y|}2\right).
\]
径向积分换元 \(s=h/r\)，得到
\[
 \int_{|x-y|>2h}|K(x,y)-K(x,y')|\,\mathrm dx
 \leq C|S^{d-1}|\int_0^{1/2}\omega(s)\frac{\mathrm ds}s.
\]
因此 \(\int_0^1\omega(s)\,\mathrm ds/s<\infty\) 也足以控制坏块尾部。若已经给定有界 \(L^2\) 算子及支集外表示，这个积分条件可代替弱型证明中的输入变量 Hölder 条件；伴随也满足对应条件时，再由插值和对偶得到全部有限的中间指数。

\ProseParagraphBreak

取 \(\omega(s)=s^\delta\) 就回到幂次情形。取 \(\omega(s)=\bigl(\log(e/s)\bigr)^{-2}\) 时，换元 \(v=\log(e/s)\) 得到可积尾部 \(\int^\infty v^{-2}\,\mathrm dv\)；改成 \(\bigl(\log(e/s)\bigr)^{-1}\) 则发散。这里区分的是这一尾积分论证能否成立，并未排除某个具体算子还有别的抵消性质。按二进环带估计，同一个条件写成
\(\sum_{j\geq1}\omega(2^{-j})<\infty\)：每个环带的体积消去核的 \(|x-y|^{-d}\)，留下的正是模量。单调性给出
\[
 (\log2)\cdot\omega(2^{-j-1})
 \leq\int_{2^{-j-1}}^{2^{-j}}\omega(s)\frac{\mathrm ds}s
 \leq(\log2)\cdot\omega(2^{-j}),
\]
从而级数与积分具有相同的收敛性。

\subsection{奇异积分算子}
\label{b2:subsec:180303}

核的条件只描述离开对角线的作用。还须指定一个算子，并独立证明其平方有界性。
\Needspace{6\baselineskip}
\begin{definition}\label{b2:thm:18-cz-operator-definition}
设整数 \(d\geq1\)，\(T:L^2(\mathbb R^d;\mathbb C)\to L^2(\mathbb R^d;\mathbb C)\) 为有界复线性算子。设可测核 \(K:(\mathbb R^d\times\mathbb R^d)\setminus\{x=y\}\to\mathbb C\)，并存在有限常数 \(C_0,C_1\geq0\)、\(0<\delta\leq1\)，使
\[
 |K(x,y)|\leq C_0|x-y|^{-d}
 \quad(x\ne y),
\]
\[
 \begin{aligned}
 &|K(x,y)-K(x',y)|+|K(y,x)-K(y,x')|\\
 &\quad\leq C_1\frac{|x-x'|^\delta}{|x-y|^{d+\delta}}
 \quad\left(x\ne y,\ |x-x'|\leq\frac{|x-y|}2\right).
 \end{aligned}
\]
若对每个支撑在某个紧集 \(F\) 内的 \(f\in L^2(\mathbb R^d)\)，在 \(\mathbb R^d\setminus F\) 上几乎处处有
\[
 Tf(x)=\int_FK(x,y)\cdot f(y)\,\mathrm dy,
\]
则称 \(T\) 为具有核 \(K\) 的\term[calderon-zygmund-suanzi]{Calderón--Zygmund 算子}[Calderón--Zygmund operator]。上述核条件称为标准核条件；第二个显示式同时包含两个变量的光滑性。
\end{definition}

平方有界性没有被核条件包含。例如零核既对应零算子，也对应任意有界乘法算子 \(M_af=af\)、\(a\in L^\infty\)：输入支集外两种输出都为零。后面的有界性理论适用于这些算子，但若只知道支集外的核，就不能判定给定算子是否等于硬截断的主值。

\ProseParagraphBreak

伴随所需的核可以直接从内积求出。取紧支集 \(f,g\in L^2\)，并假设两个支集有正距离。它们均可积，且
\[
 \int\!\int |K(x,y)|\cdot|f(y)|\cdot|g(x)|\,\mathrm dy\,\mathrm dx
 \leq C\|f\|_1\cdot\|g\|_1<\infty.
\]
这里 \(C\) 还依赖两个支集的距离。Fubini 与 \(\langle Tf,g\rangle=\langle f,T^*g\rangle\) 给出
\[
 \int f(y)\cdot\overline{T^*g(y)}\,\mathrm dy
 =\int f(y)\cdot
       \overline{\int\overline{K(x,y)}\cdot g(x)\,\mathrm dx}\,\mathrm dy.
\]
固定 \(g\)，让 \(f\) 遍历任一离开 \(\operatorname{supp}g\) 的有界紧区域上的 \(L^2\) 函数，即可识别两边的局部 \(L^2\) 函数。再以可数个这样的区域穷竭支集外，得到伴随的核
\[
 K^*(x,y)=\overline{K(y,x)}.
\]
它的大小与两个变量光滑常数仍可取 \(C_0,C_1\)，且 \(\|T^*\|_{2\to2}=\|T\|_{2\to2}\)。输入变量的光滑性在伴随后变成输出变量的光滑性，两者不能只保留其中一条。

\ProseParagraphBreak

下面从一个具体核构造算子。这里的目标是先取得独立的 \(L^2\) 界，再使用下一节的一般有界性定理。
\Needspace{6\baselineskip}
\begin{theorem}\label{b2:thm:18-homogeneous-kernel}
设整数 \(d\geq2\)，\(\Omega\in C^1(S^{d-1};\mathbb C)\)，且 \(\int_{S^{d-1}}\Omega\,\mathrm d\sigma=0\)，其中 \(\sigma\) 为球面面积测度。则对每个 \(f\in\mathcal S(\mathbb R^d)\)、每个 \(x\in\mathbb R^d\)，主值
\[
 T_\Omega f(x)=\lim_{\epsilon\downarrow0}
 \int_{|y|>\epsilon}\frac{\Omega(y/|y|)}{|y|^d}\cdot f(x-y)\,\mathrm dy
\]
存在。它给出兼容的有界线性算子 \(T_\Omega:L^p(\mathbb R^d)\to L^p(\mathbb R^d)\)、\(1<p<\infty\)，且
\[
 \|T_\Omega f\|_p\leq C_{d,p}\|\Omega\|_{C^1(S^{d-1})}\cdot\|f\|_p.
\]
其 \(L^2\) Fourier 乘子在 \(\xi\ne0\) 时为
\[
 m(\xi)=-\int_{S^{d-1}}\Omega(\theta)\cdot
 \left(\log|2\pi\xi\cdot\theta|
       +\frac{\pi\mathrm i}{2}\cdot\operatorname{sgn}(\xi\cdot\theta)\right)\mathrm d\sigma(\theta).
\]
对数奇性按球面上的可积函数解释，\(m\) 为零次齐次；其在 \(\xi=0\) 的取值可任意指定。
\end{theorem}
\begin{proof}
先固定 \(f\in\mathcal S\)、\(x\)。在 \(0<|y|<1\) 的积分中减去 \(f(x)\)，球面平均零使常数部分消失。由微积分基本定理，
\[
 |f(x-y)-f(x)|\leq |y|\sup_{|z-x|\leq1}|\nabla f(z)|,
\]
所以近原点的绝对积分由 \(C_x\int_0^1\mathrm dr\) 控制。远处 \(|x-y|\) 随 \(|y|\) 增大，Schwartz 衰减使径向尾部可积。因此主值在每个点存在。

为了求平方界，先在径向乘上 \(e^{-\eta r}\)、\(\eta>0\)，得到同样按球面平均零定义的算子 \(T_\eta\)。对 \(\operatorname{Re}z>0\)，有
\[
 \int_0^\infty\frac{e^{-zr}-e^{-r}}r\,\mathrm dr=-\log z.
\]
右半平面取连续对数：在其每个紧子集上，对 \(z\) 求导的积分由一个指数衰减函数控制，导数为 \(-1/z\)；再由 \(z=1\) 的值为零确定常数。减项 \(e^{-r}\) 与角变量无关，乘 \(\Omega\) 后的积分为零。由此得到
\[
 m_\eta(\xi)=-\int_{S^{d-1}}\Omega(\theta)\cdot
                  \log(\eta+2\pi\mathrm i\xi\cdot\theta)\,\mathrm d\sigma(\theta).
\]
具体地，先在两个有限径向截断之间计算，再与 Schwartz 频率测试函数配对。近 \(r=0\) 时，差
\(e^{-(\eta+2\pi\mathrm i\xi\cdot\theta)r}-e^{-r}\)
除以 \(r\) 后由 \(C_\eta(1+|\xi|)\) 控制；乘测试函数后可积。远处则由 \(e^{-\eta r}+e^{-r}\) 控制。因而可以在这些配对中换序并放开截断，得到 \(\widehat{T_\eta f}=m_\eta\cdot\widehat f\) 的温和分布恒等式。

需要的估计是对 \(\eta\) 与频率同时一致的。对 \(\xi\ne0\)，记 \(\omega=\xi/|\xi|\)、\(a=\eta/|\xi|\)。平均零先去掉 \(\log|\xi|\)。若 \(a\geq1\)，再去掉 \(\log a\)，余下的 \(\log(1+2\pi\mathrm i\omega\cdot\theta/a)\) 一致有界。若 \(0<a\leq1\)，对数虚部的绝对值不超过 \(\pi/2\)，实部的正部一致有界，负部则由
\(C+\bigl|\log|\omega\cdot\theta|\bigr|\) 控制。旋转不变性给出
\[
 \sup_{\omega\in S^{d-1}}
 \int_{S^{d-1}}\bigl|\log|\omega\cdot\theta|\bigr|\,\mathrm d\sigma(\theta)<\infty:
\]
旋转到一个坐标方向后，在赤道附近令 \(s=\omega\cdot\theta\)，面积密度有界，该部分归结为
\(\int_{-1/2}^{1/2}\bigl|\log|s|\bigr|\,\mathrm ds<\infty\)；其余球面上对数有界。因此
\[
 \sup_{\eta>0}\|m_\eta\|_\infty\leq C_d\|\Omega\|_\infty.
\]
对几乎每个球面点令 \(\eta\downarrow0\)，复对数趋向
\(\log|2\pi\xi\cdot\theta|+\mathrm i\pi\cdot\operatorname{sgn}(\xi\cdot\theta)/2\)。上述控制给出 \(m_\eta(\xi)\to m(\xi)\)。Plancherel 与频率上的支配收敛于是给出有界 \(L^2\) 乘子算子 \(T\)，且 \(T_\eta f\to Tf\) 于 \(L^2\)。

在物理空间中，减去 \(f(x)\) 后的近原点积分和 Schwartz 远尾，都有与 \(0<\eta\leq1\) 无关的可积控制，所以 \(T_\eta f(x)\to T_\Omega f(x)\) 对每个 \(x\) 成立。取 \(L^2\) 收敛的几乎处处子列，便把这个主值识别为 \(Tf\)。平均零也从乘子公式中消去 \(\log|\xi|\)，给出零次齐次性。

核的大小与梯度常数由 \(C_d\|\Omega\|_{C^1}\) 控制。它对光滑紧支集输入的支集外表示已由主值公式给出；对一般紧支集 \(L^2\) 输入，在任意到其支集有正距离的有界区域上，用共用一个稍大紧支集的光滑 \(L^2\) 逼近。该区域到共同支集有正距离，核积分的差由 \(L^1\) 误差一致控制；\(L^2\) 输出收敛再抽子列，即得所需支集外表示。以可数个这样的区域穷竭支集外，说明 \(T\) 满足标准算子定义。\cref{b2:thm:18-cz-lp}因此给出全部 \(1<p<\infty\) 的界。先将非零的 \(\Omega\) 除以其 \(C^1\) 范数应用该结论，再还原尺度，就得到陈述中的线性范数因子；\(\Omega=0\) 时直接成立。
\end{proof}

乘子公式的两部分来自不同的对称性。若 \(\Omega\) 为奇函数，对数项的球面积分为零；若 \(\Omega\) 为偶函数，符号函数项的积分为零。再假设 \(\Omega\) 实值时，前一种乘子为纯虚函数，后一种为实偶函数。二维的 \(\Omega(\theta)=\theta_1\theta_2\) 给出偶核 \(x_1x_2/|x|^4\)，而 \(\Omega(\theta)=\theta_j\) 给出 Riesz 型奇核。相反，\(\Omega=1\) 产生前面的径向对数发散；人为扣除发散常数属于另外一种归一化，不能据此声称原硬截断主值已经存在。

\subsection{截断}
\label{b2:subsec:180304}

取整数 \(d\geq1\)，设核在对角线外为 Borel 可测，且满足 \(|K(x,y)|\leq C_0|x-y|^{-d}\)。前面的标准核满足这一可测性：在远离对角线的小邻域内，先移动 \(x\) 再移动 \(y\)，两个 Hölder 界共同给出连续性。对 \(\epsilon>0\) 定义硬截断
\[
 T_\epsilon f(x)=\int_{|x-y|>\epsilon}K(x,y)\cdot f(y)\,\mathrm dy.
\]
若 \(f\in L^p(\mathbb R^d)\)、\(1<p<\infty\)，Hölder 给出
\[
 \int_{|x-y|>\epsilon}|K(x,y)|\cdot|f(y)|\,\mathrm dy
 \leq C_{d,p}C_0\epsilon^{-d/p}\|f\|_p.
\]
这里径向控制核属于 \(L^{p'}\)，因为 \(dp'>d\)。\(p=1\) 时改用 \(|K|\leq C_0\epsilon^{-d}\)，仍有绝对收敛。故每个固定截断在这些输入上有普通积分的意义；一般有界输入却可能在无穷远处失去可积性。

\ProseParagraphBreak

把截断半径移向零，还需另外的论证。给定的标准 \(L^2\) 算子可以含有一个支集外核看不见的乘法项，所以截断的极限不必等于它。常见的识别办法，是先对 \(C_c^\infty\) 输入验证截断的逐点极限，再证明这些截断在 \(L^2\) 中趋于给定算子；平方收敛的子列才能把两个极限识别起来。下一节先证明极大截断的有界性，再把这项识别从光滑类传给一般 \(L^p\) 输入。

\ProseParagraphBreak

不同截断约定的比较也应写出实际差值。取实函数 \(\chi\in C^\infty([0,\infty))\)，在 \([0,1]\) 为零、在 \([2,\infty)\) 为一，且 \(0\leq\chi\leq1\)，令
\[
 T^\chi_\epsilon f(x)=\int K(x,y)\cdot
                 \chi(|x-y|/\epsilon)\cdot f(y)\,\mathrm dy.
\]
与硬截断的差仅在 \(\epsilon<|x-y|<2\epsilon\)，所以对有限指数的 \(L^p\) 输入，
\[
 |T^\chi_\epsilon f(x)-T_\epsilon f(x)|
 \leq C_dC_0\epsilon^{-d}\int_{B(x,2\epsilon)}|f|
 \leq C_d'C_0\cdot Mf(x).
\]
此界能传递 \(1<p<\infty\) 的极大范数估计，但右端不随 \(\epsilon\) 消失，不能单凭这一界识别极限。

\ProseParagraphBreak

若硬截断在某点已有有限极限，则两种截断其实有相同极限，而不必再要求核的抵消。由 \(\int_1^2\chi'(s)\,\mathrm ds=1\)，对上述绝对收敛积分换序，得
\[
 T^\chi_\epsilon f(x)=\int_1^2\chi'(s)\cdot T_{s\epsilon}f(x)\,\mathrm ds.
\]
因此它与硬截断极限 \(L\) 的差不超过
\[
 \left(\int_1^2|\chi'(s)|\,\mathrm ds\right)
 \sup_{\epsilon\leq r\leq2\epsilon}|T_rf(x)-L|\longrightarrow0.
\]
对球面平均零的齐次核，还能在光滑输入上直接比较：环带差核的积分为零，减去 \(f(x)\) 后，差不超过
\(C_d\|\Omega\|_\infty\epsilon\sup_{B(x,2\epsilon)}|\nabla f|\)。这给出共同的光滑类极限；向一般输入推广仍须控制极大算子的逼近误差。这里比较的是两种截断；它们与预先给定的算子 \(T\) 是否相等，仍是前一段所说的识别问题。

\ProseParagraphBreak

截断还可以先保留一个有限的外半径，再逐次放开。对 \(L^p\)、\(1\leq p<\infty\)，固定内半径后的远场已有上述绝对积分控制，外半径趋于无穷不会改变结果。对常函数或一般有界输入，远场可能发散或仅有条件收敛，需要另行指定归一化；有限指数上建立的定义不能不加说明地用于这些输入。

% !TeX root = ../../Book2.tex
\section{\texorpdfstring{\(L^p\)}{Lᵖ} 有界性}
\label{b2:sec:1804}

分解给出弱型，插值给出下半段指数，伴随再给出上半段。三个步骤各使用不同的假设；把它们连起来才得到完整的 \(L^p\) 范围。

\subsection{\texorpdfstring{\(L^2\)}{L2} 假设}
\label{b2:subsec:180401}

全章的主要假设是 \(L^2\) 有界性。对卷积算子，它通常由有界 Fourier 乘子得到；对非平移不变的算子，需要另外的估计或测试条件。标准核条件只控制远场差分，不能单独排除局部能量增长。

\ProseParagraphBreak

在分解中，\(L^2\) 只用于好函数：其高度有限，且 \(L^1\) 范数有限，所以它落在 \(L^2\)。这样将原输入的 \(L^1\) 大小与平方估计连接起来。坏函数并没有更好的平方范数，若把整个 \(f\) 直接代入 \(L^2\) 估计，会失去所需的端点常数。

\subsection{弱型 \texorpdfstring{\((1,1)\)}{(1,1)}}
\label{b2:subsec:180402}

\begin{theorem}[Calderón--Zygmund 弱型]\label{b2:thm:18-cz-weak}
设 \(T\) 满足\cref{b2:thm:18-cz-operator-definition}。则对 \(f\in L^1\cap L^2\)、\(t>0\)，有
\[
 |\{|Tf|>t\}|\leq C t^{-1}\|f\|_1,
\]
其中常数只依赖维数、光滑指数、核常数和 \(L^2\) 算子范数。\(T\) 有唯一的 \(L^1\to L^{1,\infty}\) 连续延拓。
\end{theorem}
\begin{proof}
取高度 \(t\) 的分解。好部分、扩大坏区域及补集上的坏部分估计已经在\secref{b2:sec:1802}完成，其和给出结论。对一般 \(L^1\) 输入取 \(f_n\in L^1\cap L^2\) 范数逼近。弱估计用于差 \(f_n-f_m\)，使 \(Tf_n\) 在每个有限测度集合上依测度 Cauchy。通过选择差的超水平集合测度可求和的子列，在整数半径球上同时得到几乎处处极限；这些局部极限在交集上一致，拼成可测函数 \(Tf\)。Fatou 的集合版本把弱界传给极限。对任意另一逼近列，同样的差估计使极限相等几乎处处，因此延拓唯一。这里没有把弱拟范数误当成满足常数一的三角不等式。
\end{proof}
弱端点一般不能升级为 \(L^1\) 有界，Hilbert 变换已给反例。下一步只在开区间内插值。

\subsection{\texorpdfstring{\(1<p<\infty\)}{1<p<∞}}
\label{b2:subsec:180403}

\begin{theorem}[Calderón--Zygmund]\label{b2:thm:18-cz-lp}
设 \(T\) 满足\cref{b2:thm:18-cz-operator-definition}。则对所有 \(1<p<\infty\)，有唯一的有界延拓 \(T:L^p\to L^p\)，且
\[
 \|Tf\|_p\leq C_p\|f\|_p.
\]
\end{theorem}
\begin{proof}
在 \(1<p<2\) 时，弱 \((1,1)\) 与强 \((2,2)\) 经\cref{b2:thm:17-marcinkiewicz}得到所需界。\(T^*\) 也是标准核算子，所以同样在 \(1<p'<2\) 有界。对 \(p>2\) 的 \(f\in L^p\cap L^2\) 与测试 \(g\in L^{p'}\cap L^2\)，
\[
 \left|\int Tf\,\overline g\right|
 =\left|\int f\,\overline{T^*g}\right|
 \leq C_{p'}\|f\|_p\cdot\|g\|_{p'}.
\]
对测试 \(g\) 取上确界得到 \(L^p\) 界；稠密性给出唯一延拓。公共输入上的各个延拓由 \(L^2\) 或局部依测度极限识别，因此它们是一个兼容算子的不同空间实现。
\end{proof}
核是否齐次、算子是否卷积都没有进入这项证明。其可迁移的部分是大小、光滑、局部平方理论和零平均分解，而不是某一个特定 Fourier 符号。

\subsection{对偶性}
\label{b2:subsec:180404}

最后说明主值的点态问题。若 \(T_\varepsilon\) 在 \(L^2\) 中实现 \(T\)，定义 \(T_*f=\sup_{\varepsilon>0}|T_\varepsilon f|\)。对 \(1<q<\infty\)，核差分和 \(T\) 的 \(L^q\) 有界性给出
\begin{proposition}\label{b2:thm:18-maximal-truncation}
在上述主值实现条件下，
\[
 T_*f(x)\leq C_q\left[
 M(|Tf|^q)(x)^{1/q}+M(|f|^q)(x)^{1/q}\right]
\]
几乎处处成立。因此 \(T_*\) 在每个 \(L^p\)、\(1<p<\infty\) 上有界。
\end{proposition}
\begin{proof}
固定 \(\varepsilon\)，写 \(f=f_1+f_2\)，其中 \(f_1=f\mathbf1_{B(x,2\varepsilon)}\)。对 \(z\in B(x,\varepsilon/4)\)，\(Tf_2(z)\) 是支集外的核积分。比较它与 \(T_\varepsilon f(x)\)，中间环带的积分由大小条件控制为 \(C Mf(x)\)；远场核差分分成半径 \(2^k\varepsilon\) 的环带，每带不超过 \(C2^{-k\delta}Mf(x)\)，可求和。故
\[
 |T_\varepsilon f(x)|\leq |Tf(z)|+|Tf_1(z)|+CMf(x).
\]
对 \(z\) 取 \(q\) 次平均并用 Minkowski。第一项被 \(M(|Tf|^q)^{1/q}\) 控制；第二项由 \(L^q\) 有界性不超过
\[
 C\varepsilon^{-d/q}\|f\mathbf1_{B(x,2\varepsilon)}\|_q
 \leq C M(|f|^q)(x)^{1/q}.
\]
用 Hölder 将 \(Mf\) 吸收到同一项，再对截断半径取上确界。对给定 \(p>1\)，选 \(1<q<p\)，极大算子在 \(L^{p/q}\) 上有界，而 \(T\) 在 \(L^p\) 上有界，得到结论。
\end{proof}
若某个稠密测试类上截断逐点收敛，以上最大估计使截断差的振荡集合由逼近误差控制，按\cref{b2:thm:17-differentiation}的相同误差集合方法得到一般 \(L^p\) 输入几乎处处收敛；范数极限识别它为 \(Tf\)。这适用于 Hilbert 变换。弱 \(L^1\) 的最大截断理论需要额外端点论证，留在\appref{b2:app:C}讨论。

% 扩写定位：b2:subsec:180404
后面的平方函数需要一次处理许多输出。逐个标量估计再求平方和，通常会产生与输出个数有关的常数，因而不足以处理无限频带。
\begin{proposition}\label{b2:prop:18-hilbert-valued-cz}
设 \(\mathcal H_1,\mathcal H_2\) 为有限维 Hilbert 空间。令 \(T:L^2(\mathbb R^d;\mathcal H_1)\to L^2(\mathbb R^d;\mathcal H_2)\) 有界，核 \(K(x,y)\) 为两空间之间的线性映射，满足统一的算子范数大小与两变量 Hölder 条件，并在支集外表示 \(T\)。则 \(T\) 满足弱型 \((1,1)\) 和所有 \(1<p<\infty\) 的强型估计；常数只依赖上述界、\(p,d,\delta\)，不依赖 Hilbert 空间维数。
\end{proposition}
\begin{proof}
对标量函数 \(\|f(x)\|_{\mathcal H_1}\) 作二进坏立方体选择。好函数在每个坏块内取 Bochner 平均 \(f_Q\)，坏函数为 \(b_Q=(f-f_Q)\mathbf1_Q\)。凸性给出
\(\|g\|_\infty\leq2^d\lambda\)、\(\|g\|_1\leq\|f\|_1\)，从而
\(\|g\|_2^2\leq2^d\lambda\|f\|_1\)。
且 \(\int b_Q=0\)、\(\sum_Q\|b_Q\|_1\leq2\|f\|_1\)，坏块的测度总和不超过 \(\|f\|_1/\lambda\)。

在放大的坏块外，把积分写成
\[
 Tb_Q(x)=\int_Q[K(x,y)-K(x,c_Q)]b_Q(y)\,\mathrm dy.
\]
先取 Hilbert 范数再用算子范数控制核差，得到
\(\int_{\mathbb R^d\setminus Q^*}\|Tb_Q(x)\|\,\mathrm dx\leq C\|b_Q\|_1\)。
好部分用 \(L^2\) Chebyshev，坏部分用刚才的积分界和 \(L^1\) Chebyshev，便重复得到弱型估计。每一步常数只来自标量几何与给定的算子范数界。

对 \(1<p<2\)，将输出的 Hilbert 范数作为次线性输出；输入按 \(\|f(x)\|\) 的高低截断，本章之前的插值证明仍逐行成立。对 \(p>2\)，在 Hilbert 对偶配对中使用伴随算子；其核为 \(K(y,x)^*\)，具有同一组界。因而得到全部指数。
\end{proof}
例如，若 \(T_j\) 是一族具有统一常数的卷积型 Calderón--Zygmund 算子，则对任意有限族 \(f_j\)，
\[
 \left\|\left(\sum_j|T_jf_j|^2\right)^{1/2}\right\|_p
 \leq C_p\left\|\left(\sum_j|f_j|^2\right)^{1/2}\right\|_p.
\]
取对角算子 \(T(f_j)_j=(T_jf_j)_j\)：其核的算子范数为 \(\sup_j|k_j(x-y)|\)，其 \(L^2\) 界由逐项 Plancherel 或逐项 \(L^2\) 界相加得到。有限维一致常数和单调收敛允许推广到可数族。这里只要求对角核统一；一般矩阵核不能以每个矩阵元有界代替算子范数有界，因为矩阵大小会影响其范数。

\ProseParagraphBreak

这一证明也适用于坏块估计中只给出统一积分 Hörmander 条件的情形，前提是算子及伴随都满足对应条件。最大截断的 Cotlar 估计仍需另外检查点态条件或其他替代工具；抽象 \(L^p\) 有界性与主值几乎处收敛是两件不同的事。

无穷输入端点也存在一个可完整验证的替代估计。对局部可积 \(h\)，记
\(\|h\|_{\mathrm{BMO}}=\sup_Q|Q|^{-1}\int_Q|h-h_Q|\)，其中上确界取所有立方体，且 \(h_Q\) 为平均。该空间将在\appref{b2:app:C}系统定义；这里先证明与本章核条件直接相关的结论。
\begin{proposition}\label{b2:prop:18-infinity-bmo}
标准 Calderón--Zygmund 算子 \(T\) 对每个有界紧支集 \(f\) 满足
\(\|Tf\|_{\mathrm{BMO}}\leq C\|f\|_\infty\)，常数与输入支集大小无关。
\end{proposition}
\begin{proof}
固定立方体 \(Q\)、中心 \(c_Q\)，取充分大的固定扩大块 \(Q^*\)。分解 \(f=f_1+f_2\)，其中 \(f_1=f\mathbf1_{Q^*}\)。局部部分由平方界控制：
\[
 \frac1{|Q|}\int_Q|Tf_1|
 \leq |Q|^{-1/2}\|T\|_{2\to2}\|f_1\|_2
 \leq C\|f\|_\infty.
\]
远场在 \(Q\) 内具有支集外核表示。取比较常数
\(c=\int_{(Q^*)^c}K(c_Q,y)f(y)\,\mathrm dy\)，该积分因紧支集及离开对角线而有限。核的第一变量光滑性给出
\[
 |Tf_2(x)-c|
 \leq C\|f\|_\infty\ell(Q)^\delta
 \int_{|y-c_Q|>c_d\ell(Q)}|y-c_Q|^{-d-\delta}\,\mathrm dy
 \leq C'\|f\|_\infty.
\]
故 \(Q\) 上相对于常数 \(c\) 的平均振荡受控。换成真正平均 \((Tf)_Q\) 至多增加因子二，取全部 \(Q\) 的上确界即可。
\end{proof}
两个核变量的职责在此再次分开：弱积分端点的坏块需要输入变量的差分，无穷端点的局部振荡需要输出变量的差分。仅给其中一种光滑性，不能把另一方向的结论也视为已经证明。

\ProseParagraphBreak

当输入仅有界而不紧支集时，上面常数 \(c\) 的远场积分可能发散。不过差
\(\int[K(x,y)-K(c_Q,y)]f(y)\,\mathrm dy\) 仍绝对收敛，因其多一个衰减幂。由此可以按模常数构造局部相容的输出；附录中的 BMO 对偶也自然只见到模常数的对象。这个归一化说明了为什么端点的正确目标是平均振荡，而不是一个具有绝对数值的全局有界函数。

% !TeX root = ../../Book2.tex
\OptionalSection{Schauder 估计}{b2:sec:1605}

待撰写。

\subsection{Hölder 空间与缩放}
\label{b2:subsec:160501}

待撰写。

\subsection{Poisson 方程的内部估计}
\label{b2:subsec:160502}

待撰写。

\subsection{冻结系数与一般内部估计}
\label{b2:subsec:160503}

待撰写。

\subsection{边界估计的条件与理论边界}
\label{b2:subsec:160504}

待撰写。


\begin{chaptersummary}
Calderón--Zygmund 分解在每个坏立方体上保留平均值，把剩余部分变成零平均函数。远离这个立方体时，核可以减去中心值；核差分比核本身多衰减一个正幂次，于是坏部的远场可以绝对积分。好部由平方估计控制，坏区由体积估计控制，这三项合起来得到弱 \((1,1)\)。再经插值与伴随，标准 Calderón--Zygmund 算子在全部 \(1<p<\infty\) 上有兼容的延拓。

\ProseParagraphBreak

使用这一结论时，还要分别核对核条件、\(L^2\) 有界性和截断实现。离支集核并不决定算子在对角线上的作用，\(L^p\) 范数延拓也不自动等于逐点主值。极大截断提供统一的误差控制，具体核在光滑稠密类上的收敛提供起点，两者结合才把主值传给一般输入。Hilbert 与 Riesz 变换满足这些条件，它们的零阶 Fourier 符号又把共轭调和结构和二阶椭圆估计连接起来。区域问题仍须经过截断和局部化，所得估计中的低阶项不能在没有额外条件时删去。
\end{chaptersummary}

\BookExercises


% 自拟：由区间模型精算分布函数与弱范数，训练端点的定量判断。
\Needspace{6\baselineskip}
\begin{exercise}\label{b2:ex:18-new-1}
在带 Lebesgue 测度的 \(\mathbb R\) 上，令
\[
 h=H\mathbf1_{(0,1)},\quad
 h(x)=\frac1\pi\log\left|\frac{x}{x-1}\right|\quad(x\ne0,1),
\]
并在两个端点任意补有限值。对 \(t>0\)，求出 \(\{|h|>t\}\) 的区间组成，证明其分布函数与弱 \(L^1\) 准范数分别为
\[
 d_h(t)=\frac2{\sinh(\pi t)},
 \quad \|h\|_{L^{1,\infty}}
       \coloneqq\sup_{t>0}t\cdot d_h(t)=\frac2\pi.
\]
分别解释 \(t\downarrow0\) 和 \(t\uparrow\infty\) 的分布行为，并据此判断 \(h\) 的有限指数可积性。
\end{exercise}
\begin{solution}
固定 \(t>0\)，记 \(q=e^{\pi t}>1\)。条件 \(|h(x)|>t\) 等价于
\[
 \left|\frac{x}{x-1}\right|<q^{-1}
 \quad\text{或}\quad
 \left|\frac{x}{x-1}\right|>q.
\]
在 \(x<0\) 上，比值的绝对值为 \(-x/(1-x)<1\)，第一项不等式给出
\(-1/(q-1)<x<0\)。在 \(0<x<1\) 上，比值为 \(x/(1-x)\)，两项分别给出
\(0<x<1/(q+1)\) 和 \(q/(q+1)<x<1\)。在 \(x>1\) 上，比值为 \(x/(x-1)>1\)，第二项给出
\(1<x<q/(q-1)\)。因此，忽略有限多个点后，超水平集就是
\[
 \left(-\frac1{q-1},\frac1{q+1}\right)
 \ \cup\
 \left(\frac q{q+1},\frac q{q-1}\right).
\]
两段长度之和为
\[
 d_h(t)=2\left(\frac1{q-1}+\frac1{q+1}\right)
       =\frac{4q}{q^2-1}=\frac2{\sinh(\pi t)}.
\]
由 \(\sinh s\geq s\)（\(s>0\)），有 \(t\cdot d_h(t)\leq2/\pi\)；又由
\(\sinh s/s\to1\)（\(s\downarrow0\)），这个上界在极限中达到。所以弱范数等于 \(2/\pi\)，其上确界无需在某个正高度取到。

更具体地，低高度时 \(d_h(t)\sim2/(\pi t)\)，记录的是远处 \(1/x\) 尾部占据的长区间；高高度时 \(d_h(t)\sim4e^{-\pi t}\)，记录的是两个对数奇点附近迅速缩小的区域。用分布函数恒等式
\[
 \|h\|_p^p=p\int_0^\infty t^{p-1}\cdot d_h(t)\,\mathrm dt
 \quad(0<p<\infty),
\]
无穷端因指数衰减总可积，零端则等价于 \(\int_0^1t^{p-2}\,\mathrm dt\)。故恰在 \(p>1\) 时 \(h\in L^p\)，在 \(0<p\leq1\) 时不属于 \(L^p\)。此外高高度的分布函数始终为正，故 \(h\notin L^\infty\)。
\end{solution}

% 自拟：将 Hilbert 变换的代数恒等式与正负频率结合，判断可逆性及精确平方范数。
\Needspace{6\baselineskip}
\begin{exercise}\label{b2:ex:18-new-2}
设 \(1<p<\infty\)，\(a,b\in\mathbb C\)，\(H\) 取本章兼容的 \(L^p(\mathbb R;\mathbb C)\) 延拓。证明 \(A=aI+bH\) 为可逆有界算子，当且仅当 \(a^2+b^2\ne0\)，并在可逆时求其逆。再在 \(L^2\) 上证明
\[
 \|A\|_{2\to2}=\max\{|a-\mathrm ib|,|a+\mathrm ib|\},
 \quad
 \|A^{-1}\|_{2\to2}
 =\frac1{\min\{|a-\mathrm ib|,|a+\mathrm ib|\}}.
\]
必要性与范数下界须给出实际输入，不能只在一个频率点上检验符号。
\end{exercise}
\begin{solution}
先将正文的平方恒等式传到所用空间。在 \(L^2\) 中已有 \(H^2=-I\)。对 \(\varphi\in\mathcal S\)，其 \(L^2\) 与 \(L^p\) 实现相容，且 \(H\varphi\in L^2\cap L^p\)，所以该恒等式也在 \(L^p\) 中成立。对一般 \(f\in L^p\)，取 \(\varphi_n\in\mathcal S\) 在 \(L^p\) 中逼近；两次使用 \(H\) 的有界性，就得到 \(H^2f=-f\)。因此
\[
 (aI+bH)(aI-bH)=(aI-bH)(aI+bH)=(a^2+b^2)I.
\]
若 \(a^2+b^2\ne0\)，则
\[
 A^{-1}=\frac{aI-bH}{a^2+b^2}
\]
是所需的有界双侧逆。

由于 \(a^2+b^2=(a-\mathrm ib)(a+\mathrm ib)\)，退化时至少有一个因子为零。取非零 \(\psi_+\in C_c^\infty\bigl((1,2)\bigr)\)，令 \(f_+=\mathcal F^{-1}\psi_+\in\mathcal S\)。其频率全为正，故
\(Hf_+=-\mathrm if_+\)，从而 \(Af_+=(a-\mathrm ib)f_+\)。同样取非零频率测试函数支撑于 \((-2,-1)\)，得到 \(f_-\in\mathcal S\) 且
\(Af_-=(a+\mathrm ib)f_-\)。每个 Schwartz 函数都属于当前的 \(L^p\)，所以任一因子为零都使 \(A\) 有非零核，不能可逆。

在 \(L^2\) 中，Plancherel 给出
\[
 \|Af\|_2^2
 =|a-\mathrm ib|^2\int_{\xi>0}|\widehat f(\xi)|^2\,\mathrm d\xi
  +|a+\mathrm ib|^2\int_{\xi<0}|\widehat f(\xi)|^2\,\mathrm d\xi.
\]
这给出最大值所对应的上界，而选取 \(f_+\) 或 \(f_-\) 就得到相应下界。可逆时，逆算子在两个半轴上的符号分别为
\((a-\mathrm ib)^{-1}\) 和 \((a+\mathrm ib)^{-1}\)，用同样的输入得到其精确范数。正、负半轴都是正测度集，所以这些范数由真正的 \(L^2\) 函数实现。
\end{solution}

% 自拟：区分算子范数最优与非零函数能否取等，训练频率集中构造。
\Needspace{6\baselineskip}
\begin{exercise}\label{b2:ex:18-new-3}
设整数 \(d\geq2\)、\(1\leq j\leq d\)，\(R_j\) 为乘子
\(-\mathrm i\xi_j/|\xi|\) 定义的 \(L^2(\mathbb R^d;\mathbb C)\) Riesz 变换。证明
\(\|R_j\|_{2\to2}=1\)，但每个非零 \(f\in L^2\) 都满足
\(\|R_jf\|_2<\|f\|_2\)。构造范数为一的 Schwartz 函数族，使其输出范数趋于一。
\end{exercise}
\begin{solution}
由 Plancherel，
\[
 \|f\|_2^2-\|R_jf\|_2^2
 =\int_{\mathbb R^d}
   \frac{\sum_{k\ne j}\xi_k^2}{|\xi|^2}
       \cdot|\widehat f(\xi)|^2\,\mathrm d\xi\geq0.
\]
权重只在第 \(j\) 条坐标轴上为零，原点处的补值不影响积分。由于 \(d\geq2\)，这条轴的 \(d\) 维 Lebesgue 测度为零。若 \(f\ne0\)，则 \(|\widehat f|>0\) 的集合有正测度，故上式严格为正。特别地，范数等式不能由某个非零输入取到。

为逼近上界，记 \(e_j\) 为第 \(j\) 个单位坐标向量，取非零
\(\psi\in C_c^\infty\bigl(B(0,1)\bigr)\)。对 \(0<\varepsilon<1/2\)，定义
\[
 \widehat f_\varepsilon(\xi)
 =\frac{\varepsilon^{-d/2}}{\|\psi\|_2}
       \cdot\psi\left(\frac{\xi-e_j}{\varepsilon}\right).
\]
该函数光滑且紧支撑，所以 \(f_\varepsilon\in\mathcal S\)；换元给出
\(\|f_\varepsilon\|_2=1\)。在其频率支集上，\(\xi_j\geq1-\varepsilon\)、\(|\xi|\leq1+\varepsilon\)，所以
\[
 \left(\frac{1-\varepsilon}{1+\varepsilon}\right)^2
 \leq\|R_jf_\varepsilon\|_2^2\leq1.
\]
令 \(\varepsilon\downarrow0\)，输出范数便趋于一。这说明最优算子范数可以由集中族逼近，即使没有任何非零取等函数。维数条件也有实质作用：在一维，\(R_1=H\) 在全部 \(L^2\) 上等距。
\end{solution}

% 自拟：将坏块零平均推广到有限阶矩消失，并用有限差分构造精确阶数。
\Needspace{6\baselineskip}
\begin{exercise}\label{b2:ex:18-new-4}
设 \(m\geq0\) 为整数、\(R>0\)，\(b\in L^1(\mathbb R;\mathbb C)\) 在 \([-R,R]\) 外几乎处处为零。定义
\[
 \mu_\ell\coloneqq\int_{\mathbb R}y^\ell\cdot b(y)\,\mathrm dy
 \quad(0\leq\ell\leq m),
\]
并假设 \(\mu_\ell=0\) 对 \(0\leq\ell<m\) 成立。当 \(m=0\) 时这一组消失条件为空。此题只在支集外定义普通绝对积分
\[
 Hb(x)\coloneqq\frac1\pi\int_{\mathbb R}\frac{b(y)}{x-y}\,\mathrm dy
 \quad(|x|>R),
\]
不预设一般 \(L^1\) 输入在支集内的主值。证明当 \(|x|>2R\) 时，
\[
 \left|Hb(x)-\frac{\mu_m}{\pi x^{m+1}}\right|
 \leq\frac{2R^{m+1}}{\pi|x|^{m+2}}\|b\|_1.
\]
再对
\[
 b_m(y)=\sum_{k=0}^m(-1)^k\binom mk\cdot\mathbf1_{(k,k+1)}(y)
\]
证明前 \(m\) 个矩为零，而 \(\mu_m=(-1)^m m!\)，从而该尾部阶数可以实际出现。
\end{exercise}
\begin{solution}
紧支集保证每个所用的矩都绝对收敛。在 \(|x|>2R\)、\(|y|\leq R\) 时，有限几何展开给出恒等式
\[
 \frac1{x-y}=\sum_{\ell=0}^m\frac{y^\ell}{x^{\ell+1}}
                  +\frac{y^{m+1}}{x^{m+1}(x-y)}.
\]
乘 \(b(y)\) 积分，前 \(m\) 项因矩条件消失。余项中
\(|x-y|\geq|x|/2\)，故
\[
 \begin{aligned}
 \left|Hb(x)-\frac{\mu_m}{\pi x^{m+1}}\right|
 &\leq\frac2{\pi|x|^{m+2}}
             \int|y|^{m+1}\cdot|b(y)|\,\mathrm dy\\
 &\leq\frac{2R^{m+1}}{\pi|x|^{m+2}}\|b\|_1.
 \end{aligned}
\]
估计对正、负两侧的 \(x\) 都成立，首项的符号保留在 \(x^{m+1}\) 中。

对给定的 \(b_m\)，它支撑于 \([0,m+1]\)，因而属于 \(L^1\)。将每个区间换元为 \(y=z+k\)，得到
\[
 \int y^\ell\cdot b_m(y)\,\mathrm dy
 =\int_0^1\sum_{k=0}^m(-1)^k\binom mk(z+k)^\ell\,\mathrm dz.
\]
记前向差分为 \(\Delta P(z)=P(z+1)-P(z)\)。反复展开可得
\[
 \Delta^mP(z)=\sum_{k=0}^m(-1)^{m-k}\binom mk\cdot P(z+k).
\]
差分每次将多项式次数降低一，且把最高次系数乘以原来的次数。因此对 \(P(z)=z^\ell\)，\(\ell<m\) 时 \(\Delta^mP=0\)，\(\ell=m\) 时 \(\Delta^mP=m!\)。上面的矩积分于是分别为零和 \((-1)^m m!\)；\(m=0\) 时同样得到 \(\mu_0=1\)。由于第 \(m\) 个矩非零，所得尾部的首项确为 \(x^{-(m+1)}\) 阶，而非仅有一个可能不尖锐的上界。
\end{solution}

% 自拟：从离支集局部性识别对角线自由度，补足零核乘法例子的结构解释。
\Needspace{6\baselineskip}
\begin{exercise}\label{b2:ex:18-new-5}
设整数 \(d\geq1\)，\(T_1,T_2\) 是 \(L^2(\mathbb R^d;\mathbb C)\) 上的有界复线性算子。设可测核 \(K\) 在 \(x\ne y\) 上满足
\(|K(x,y)|\leq A|x-y|^{-d}\)，其中 \(A<\infty\)；并且对每个紧集 \(F\) 与每个在 \(F\) 外几乎处处为零的 \(f\in L^2\)，有
\[
 T_1f(x)=T_2f(x)=\int_FK(x,y)\cdot f(y)\,\mathrm dy
 \quad\text{对几乎处处的 }x\notin F.
\]
证明差算子 \(S=T_1-T_2\) 必为乘法算子：存在唯一的 \(a\in L^\infty(\mathbb R^d)\) 等价类，使
\[
 Sf=a\cdot f\quad(f\in L^2),
 \quad\|S\|_{2\to2}=\|a\|_\infty.
\]
证明中须从紧支集输入出发，解释为什么 \(S\) 对任意可测集合都保持支撑，而不能只对开集作形式判断。
\end{exercise}
\begin{solution}
若 \(f\) 支撑于紧集 \(F\)，共同核表示相减便知 \(Sf=0\) 在 \(F\) 外几乎处处成立。先把这个局部性推广到任意 Lebesgue 可测集 \(E\)。设 \(f\in L^2\) 在 \(E\) 外几乎处处为零，选定有限值可测代表，并令
\[
 E_n=E\cap B(0,n)\cap\{|f|\leq n\}.
\]
这些集合递增且覆盖 \(f\) 的非零部分，至多相差零测集。由有限测度集合的内正则性，选紧集 \(K_n\subset E_n\)，使
\(|E_n\setminus K_n|<n^{-4}\)。令 \(f_n=f\cdot\mathbf1_{K_n}\)，则
\[
 \|f-f_n\|_2^2
 \leq\int_{E\setminus E_n}|f|^2+n^2|E_n\setminus K_n|
 \leq\int_{E\setminus E_n}|f|^2+n^{-2}\longrightarrow0.
\]
每个 \(Sf_n\) 都在 \(K_n\subset E\) 外为零，且 \(Sf_n\to Sf\) 于 \(L^2\)。乘以 \(\mathbf1_{E^c}\) 是有界算子，故极限 \(Sf\) 也在 \(E\) 外为零。这里不要求 \(K_n\) 递增。

对任意 \(f\in L^2\)，将它分成 \(\mathbf1_E\cdot f\) 与 \(\mathbf1_{E^c}\cdot f\)，两部分的输出分别留在 \(E\) 与 \(E^c\)。因此
\[
 S(\mathbf1_E\cdot f)=\mathbf1_E\cdot Sf.
\]
现在令 \(D_n=(-n,n)^d\)、\(a_n=S\mathbf1_{D_n}\)。当 \(m<n\) 时，
\[
 \mathbf1_{D_m}\cdot a_n
 =S(\mathbf1_{D_m}\cdot\mathbf1_{D_n})
 =S\mathbf1_{D_m}=a_m.
\]
取去掉可数个零测例外集的共同满测集，这些函数就在各 \(D_n\) 上相容，可以拼成一个可测函数 \(a\)，满足 \(a=a_n\) 于 \(D_n\)。对每个可测 \(E\subset D_n\)，又有
\[
 S\mathbf1_E=\mathbf1_E\cdot a,
 \quad\|a\cdot\mathbf1_E\|_2
      \leq\|S\|_{2\to2}\cdot|E|^{1/2}.
\]
若 \(\{|a|>\|S\|_{2\to2}+\varepsilon\}\cap D_n\) 有正测度，以上式中的 \(E\) 取为此集合，平方后就与 \(L^2\) 积分下界矛盾。因此对每个 \(n\) 和 \(\varepsilon>0\)，该集合测度为零，得到 \(\|a\|_\infty\leq\|S\|_{2\to2}\)。

由线性，\(Sf=a\cdot f\) 对支撑在某个 \(D_n\) 内的简单函数成立。这样的简单函数在 \(L^2\) 中稠密：先截去空间远尾和过高幅度，再将有界的复值范围作有限网格近似。两边均为有界算子，故等式推广到全部 \(L^2\)。反向不等式
\(\|S\|_{2\to2}=\|M_a\|_{2\to2}\leq\|a\|_\infty\)
来自直接积分，于是范数相等。若另一个 \(\widetilde a\) 给出相同乘法算子，作用于各 \(\mathbf1_{D_n}\) 即得 \(a=\widetilde a\) 几乎处处，证明唯一性。可见离支集核所遗漏的，恰好是这样的有界乘法项。
\end{solution}

% 自拟：组合方向乘子得到向量投影，连接奇异积分与无散条件。
\Needspace{6\baselineskip}
\begin{exercise}\label{b2:ex:18-new-6}
设整数 \(d\geq2\)、\(1<p<\infty\)，\(R_j\) 取本章兼容的 \(L^p(\mathbb R^d;\mathbb C)\) Riesz 变换。对向量函数 \(F=(F_1,\ldots,F_d)\in L^p(\mathbb R^d;\mathbb C^d)\)，定义
\[
 (\mathbb PF)_i=F_i+\sum_{j=1}^dR_iR_jF_j
 \quad(1\leq i\leq d),
\]
其中向量范数采用逐点欧氏范数。证明 \(\mathbb P\) 在该向量值 \(L^p\) 空间上有界，\(\mathbb P^2=\mathbb P\)，且
\(\operatorname{div}(\mathbb PF)=0\) 于分布意义；对每个 \(\varphi\in\mathcal S(\mathbb R^d)\)，还有 \(\mathbb P\nabla\varphi=0\)。再证明在 \(p=2\) 时，\(\mathbb P\) 是到闭子空间
\[
 \mathcal V=\{F\in L^2(\mathbb R^d;\mathbb C^d):
                         \operatorname{div}F=0\text{ 于分布意义}\}
\]
的正交投影，且 \(\|\mathbb P\|_{2\to2}=1\)。
\end{exercise}
\begin{solution}
先看有界性。每个分量的 \(L^p\) 范数不超过 \(\|F\|_{L^p(\mathbb C^d)}\)，且 \(R_i,R_j\) 均为有界 \(L^p\) 算子。有限次使用三角不等式，例如
\[
 \begin{aligned}
 \|\mathbb PF\|_{L^p(\mathbb C^d)}
 &\leq\sum_{i=1}^d\|(\mathbb PF)_i\|_p\\
 &\leq\sum_i\|F_i\|_p
   +\sum_{i,j}\|R_i\|_{p\to p}\cdot\|R_j\|_{p\to p}\cdot\|F_j\|_p\\
 &\leq C_{d,p}\|F\|_{L^p(\mathbb C^d)},
 \end{aligned}
\]
就给出足够的界，不要求这个常数最优。

在 \(L^2\) 中，\(R_iR_j\) 的符号为 \(-\xi_i\xi_j/|\xi|^2\)，所以 \(\mathbb P\) 的矩阵符号是
\[
 P(\xi)=I-\frac{\xi\xi^{\mathsf T}}{|\xi|^2}
 \quad(\xi\ne0).
\]
它是到 \(\xi^\perp\subset\mathbb C^d\) 的欧氏正交投影，因此
\(P(\xi)^2=P(\xi)=P(\xi)^*\)、\(\xi^{\mathsf T}P(\xi)=0\)，且矩阵范数不超过一。Plancherel 把这些等式传给 \(\mathbb P\)，得到幂等性、自伴性、平方范数上界及无散性。

还须识别投影的像。对 \(F\in L^2\)，分布恒等式 \(\operatorname{div}F=0\) 等价于
\(\xi\cdot\widehat F(\xi)=0\) 几乎处处：Fourier 变换将散度化为 \(2\pi\mathrm i\xi\cdot\widehat F\)，后者局部可积，一个局部可积函数给出零分布当且仅当它几乎处处为零。这样的 \(F\) 满足 \(P(\xi)\widehat F=\widehat F\)，即 \(\mathbb PF=F\)。而任意输出已经无散，所以像恰为 \(\mathcal V\)。若一列无散向量在 \(L^2\) 中收敛，与每个紧支集光滑测试函数的梯度配对都收敛，极限仍无散，故 \(\mathcal V\) 闭。结合自伴与幂等，得到所述正交投影。

范数确实为一，而不只是至多一。取 \(\psi\in C_c^\infty(\mathbb R^d)\)，使
\((\partial_1\psi,\partial_2\psi)\) 不恒为零，并令
\[
 G=(\partial_2\psi,-\partial_1\psi,0,\ldots,0).
\]
混合偏导相等使 \(\operatorname{div}G=0\)，且 \(G\ne0\)，故 \(\mathbb PG=G\) 给出下界一。另一方面，\(\widehat{\nabla\varphi}=2\pi\mathrm i\xi\cdot\widehat\varphi\) 沿 \(\xi\) 方向，矩阵 \(P(\xi)\) 将其消去，所以 \(\mathbb P\nabla\varphi=0\) 在 \(L^2\) 中成立。

最后处理一般 \(p\)。对 Schwartz 向量，上述等式先在 \(L^2\) 成立；它和其中的 \(\mathbb P\) 输出均属于 \(L^2\cap L^p\)，各个 Riesz 变换的相容性使幂等式也在 \(L^p\) 中成立。取 Schwartz 向量在 \(L^p\) 中逼近任意 \(F\)，利用有界性即可传递 \(\mathbb P^2=\mathbb P\)。对每个紧支集光滑测试函数，其梯度属于 \(L^{p'}\)，所以输出的无散分布等式也能在该极限中传递。梯度消去公式同样由相容性成立。这里的正交性与范数一专指 \(p=2\)，未对其他指数作相同断言。
\end{solution}

% 自拟：把热核积分、预解方程与二阶 Riesz 乘子连接，并区分有质量与零质量估计。
\Needspace{6\baselineskip}
\begin{exercise}\label{b2:ex:18-new-7}
设整数 \(d\geq1\)、\(1<p<\infty\)、\(\lambda>0\)，\(f\in L^p(\mathbb R^d;\mathbb C)\)。令
\[
 H_t(x)=(4\pi t)^{-d/2}e^{-|x|^2/(4t)},
 \quad u_\lambda=\int_0^\infty e^{-\lambda t}\cdot(H_t*f)\,\mathrm dt,
\]
时间积分按 \(L^p\) 范数中的积分理解。证明它收敛，且
\[
 (\lambda-\Delta)u_\lambda=f\text{ 于分布意义},
 \quad\|u_\lambda\|_p\leq\lambda^{-1}\|f\|_p.
\]
先对 \(f\in\mathcal S\) 求出 \(\widehat u_\lambda\) 并证明 \(u_\lambda\in\mathcal S\)，再由
\[
 \partial_i\partial_ju_\lambda
 =R_iR_j(f-\lambda u_\lambda)
\]
证明全部二阶分布导数属于 \(L^p\)，且
\[
 \sum_{i,j=1}^d\|\partial_i\partial_ju_\lambda\|_p
 \leq C_{d,p}\|f\|_p,
\]
其中常数与 \(\lambda>0\) 无关。最后用缩放说明不存在对所有 \(u\in C_c^\infty(\mathbb R^d)\) 都成立的统一估计
\(\|u\|_p\leq C\|\Delta u\|_p\)。
\end{exercise}
\begin{solution}
先核对时间积分的连续性。固定 \(t_0>0\)，当 \(t_0/2\leq t\leq2t_0\) 时，热核的显式公式给出共同可积控制 \(C_{d,t_0}e^{-|x|^2/(8t_0)}\)。支配收敛说明 \(\|H_t-H_{t_0}\|_1\to0\)，Young 不等式因而给出
\[
 \|H_t*f-H_{t_0}*f\|_p
 \leq\|H_t-H_{t_0}\|_1\cdot\|f\|_p\longrightarrow0.
\]
零时的 \(L^p\) 连续性由\cref{b2:thm:12-heat-convolution}给出。热核非负且积分为一，Young 不等式还给出
\[
 \|H_t*f\|_p\leq\|f\|_p,
 \quad\int_0^\infty e^{-\lambda t}\|H_t*f\|_p\,\mathrm dt
                  \leq\lambda^{-1}\|f\|_p.
\]
因而可先在紧时间区间内积分，再放开两端，得到收敛的 \(L^p\) 值积分及所求零阶界。特别地，\(f\mapsto u_\lambda\) 是范数不超过 \(\lambda^{-1}\) 的线性算子。

若 \(f\in\mathcal S\)，还可在 \(L^1\) 中作同一积分，因为
\(\int_0^\infty e^{-\lambda t}\|H_t*f\|_1\,\mathrm dt<\infty\)。Fubini 因而允许先取 Fourier 变换，得到
\[
 \widehat u_\lambda(\xi)
 =\int_0^\infty e^{-t(\lambda+4\pi^2|\xi|^2)}
                 \cdot\widehat f(\xi)\,\mathrm dt
 =\frac{\widehat f(\xi)}{\lambda+4\pi^2|\xi|^2}.
\]
分母严格为正，其倒数是光滑函数，各阶导数至多作多项式增长，所以乘上 Schwartz 函数仍为 Schwartz 函数。Fourier 反变换说明 \(u_\lambda\in\mathcal S\)，并立即给出
\((\lambda-\Delta)u_\lambda=f\)。

记 \(g=f-\lambda u_\lambda\)。由刚求出的公式，
\[
 \widehat g(\xi)=\frac{4\pi^2|\xi|^2}{\lambda+4\pi^2|\xi|^2}
                       \cdot\widehat f(\xi).
\]
两个 Riesz 符号相乘为 \(-\xi_i\xi_j/|\xi|^2\)，故
\[
 \widehat{R_iR_jg}(\xi)
 =-\frac{4\pi^2\xi_i\xi_j}{\lambda+4\pi^2|\xi|^2}
                       \cdot\widehat f(\xi)
 =\widehat{\partial_i\partial_ju_\lambda}(\xi).
\]
零频处任意补值不影响等式。使用 \(R_iR_j\) 的 \(L^p\) 有界性与零阶界，有
\[
 \begin{aligned}
 \sum_{i,j}\|\partial_i\partial_ju_\lambda\|_p
 &\leq C_{d,p}\|f-\lambda u_\lambda\|_p\\
 &\leq C_{d,p}(\|f\|_p+\lambda\|u_\lambda\|_p)
 \leq2C_{d,p}\|f\|_p.
 \end{aligned}
\]
吸收因子二即可得到与 \(\lambda\) 无关的常数。

对一般 \(f\in L^p\)，取 \(f_n\in\mathcal S\) 在 \(L^p\) 中逼近。零阶界使对应的 \(u_{\lambda,n}\to u_\lambda\) 于 \(L^p\)，而
\(R_iR_j(f_n-\lambda u_{\lambda,n})\)
也在 \(L^p\) 中收敛到 \(R_iR_j(f-\lambda u_\lambda)\)。对任意 \(\varphi\in C_c^\infty\)，Hölder 不等式保证
\[
 \int u_{\lambda,n}\cdot\partial_i\partial_j\varphi
 \longrightarrow\int u_\lambda\cdot\partial_i\partial_j\varphi.
\]
所以二阶分布导数的极限正是 \(\partial_i\partial_ju_\lambda\)，并由上述 \(L^p\) 极限表示。对方程作同样的测试配对，得到其分布版本；范数估计也随 \(L^p\) 收敛传递。

零质量参数不能提供同样的零阶控制。固定非零 \(\phi\in C_c^\infty\)，使 \(\Delta\phi\ne0\)，令 \(v_R(x)=\phi(x/R)\)、\(R>0\)。换元给出
\[
 \|v_R\|_p=R^{d/p}\|\phi\|_p,
 \quad\|\Delta v_R\|_p=R^{d/p-2}\|\Delta\phi\|_p.
\]
若存在所述统一常数，则 \(R^2\|\phi\|_p\leq C\|\Delta\phi\|_p\) 对所有 \(R\) 成立，与 \(R\uparrow\infty\) 矛盾。二阶估计中相同的缩放因子可以相消，零阶估计却不能，这解释了 \(\lambda>0\) 在前一估计中的作用。
\end{solution}

% 自拟：用光滑乘法系数的差消去对角线奇性，比较差核的逐点极限与平方范数极限。
\Needspace{6\baselineskip}
\begin{exercise}\label{b2:ex:18-new-8}
设 \(d\geq1\) 为整数，\(T\) 是 \(\mathbb R^d\) 上\cref{b2:thm:18-cz-operator-definition}的标准 Calder\'on--Zygmund 算子，核为 \(K\)。另设硬截断
\[
 T_\epsilon g(x)=\int_{|x-y|>\epsilon}K(x,y)\cdot g(y)\,\mathrm dy
\]
对每个 \(g\in C_c^\infty(\mathbb R^d;\mathbb C)\)，当 \(\epsilon\downarrow0\) 时在 \(L^2\) 中趋于 \(Tg\)。给定 \(a,f\in C_c^\infty(\mathbb R^d;\mathbb C)\)，令 \(C_af=a\cdot Tf-T(a\cdot f)\)。证明
\[
 C_af(x)=\int_{\mathbb R^d}\bigl(a(x)-a(y)\bigr)\cdot K(x,y)\cdot f(y)\,\mathrm dy
\]
几乎处处成立，且右端对每个 \(x\) 都绝对收敛；解释乘法系数的差如何改善对角线附近的奇性。再证明对每个 \(1<p<\infty\)，有
\[
 \|C_af\|_p\leq2\|a\|_\infty\cdot\|T\|_{L^p\to L^p}\cdot\|f\|_p.
\]
这里 \(T\) 的 \(L^p\) 实现取\cref{b2:thm:18-cz-lp}的兼容延拓，不预先假定 \(T_\epsilon f\) 单独几乎处处收敛。
\end{exercise}
\begin{solution}
设核的大小常数为 \(A\)，\(a\) 的全局 Lipschitz 常数为 \(L_a\)。当 \(|x-y|<1\) 时，
\[
 \Bigl|\bigl(a(x)-a(y)\bigr)\cdot K(x,y)\cdot f(y)\Bigr|
 \leq A L_a\|f\|_\infty\cdot|x-y|^{1-d}.
\]
这在单位球上可积，径向积分为 \(\int_0^1\mathrm dr\)；系数之差提供的 \(|x-y|\) 消去了一阶奇性。在 \(|x-y|\geq1\) 时，绝对积分不超过 \(2A\|a\|_\infty\cdot\|f\|_1\)。因此对每个 \(x\)，差核的全积分绝对收敛，并且
\[
 a(x)\cdot T_\epsilon f(x)-T_\epsilon(a\cdot f)(x)
 =\int_{|x-y|>\epsilon}\bigl(a(x)-a(y)\bigr)\cdot K(x,y)\cdot f(y)\,\mathrm dy
\]
随 \(\epsilon\downarrow0\) 趋于该积分。另一方面，\(f,a\cdot f\in C_c^\infty\)，题设的 \(L^2\) 收敛及乘法算子的有界性使左端在 \(L^2\) 中趋于 \(a\cdot Tf-T(a\cdot f)\)。从任意趋零截断序列中抽取一个 \(L^2\) 收敛的逐点子列，就把后者与已存在的全参数差核极限识别，得到所求几乎处处等式。这一步只比较两项之差，没有预设两项各自的逐点极限。

最后利用 \(T\) 的 \(L^p\) 有界性及 \(L^2\)、\(L^p\) 实现的相容性，
\[
 \|a\cdot Tf-T(a\cdot f)\|_p
 \leq\|a\|_\infty\cdot\|Tf\|_p+\|T(a\cdot f)\|_p
 \leq2\|a\|_\infty\cdot\|T\|_{L^p\to L^p}\cdot\|f\|_p.
\]
绝对积分表示使用了 \(a\) 的 Lipschitz 正则性，而这个范数估计只用其有界性，两项结论的假设作用不同。
\end{solution}

